% 第14章：大规模与在线重建
\chapter{大规模与在线重建}
\label{ch:14}

\paragraph*{本章阅读锚点}
本章的锚点是规模。原始 3DGS 面向单场景训练；大规模方法要解决分块、合并、流式加载、分布式训练和跨区域一致性。


3D Gaussian Splatting 在单个物体或小规模室内场景中已展现出令人印象深刻的效果——数十万至数百万高斯基元即可实现照片级真实感的实时渲染。然而，当我们将视野从一栋建筑扩展到整个街区、从一条街道延伸至一座城市、从单次拍摄延伸到持续数月的在线更新时，3DGS 面临的根本挑战发生了质变。具体而言，参数规模从百万级跃升至十亿级，训练数据从数千张图像增长到数万甚至数十万张，场景覆盖范围从几十米扩展到数公里。与此同时，时间维度也从一次性的离线优化，变为需要持续适应场景变化的在线增量更新。

这些挑战可以归纳为四个核心问题。

\textbf{内存墙}（memory wall）：3DGS 的每个高斯基元携带 59 个参数（3 位置 + 4 四元数 + 3 缩放 + 1 不透明度 + 48 球谐系数），十亿个高斯需要约 55 GB 参数存储。加上梯度和优化器状态，总内存消耗轻松超过单张 GPU 的容量。Grendel（2024.06）~\cite{2406_18533} 的实验表明，在 Rubble 数据集上，单 GPU 仅能容纳约 1120 万个高斯，而 16 GPU 分布式训练可将高斯数量扩展至 4040 万，PSNR 从 26.28 提升至 27.28。

\textbf{计算墙}（compute wall）：大规模场景的训练时间从数十分钟增长到数十小时甚至数天。例如，COLMAP 的 SfM（运动恢复结构，Structure from Motion，即从多视角图像恢复三维结构和相机位姿的技术）处理万张图像可能需要十余小时，3DGS 优化本身也需要数小时。

\textbf{质量墙}（quality wall）：大规模城市场景涉及天空、远景、弱纹理区域、动态物体、跨视角外观差异等诸多复杂因素。在这些条件下，标准 3DGS 的渲染质量会显著退化。

\textbf{时效墙}（timeliness wall）：现实场景不是静止的。车辆移动、建筑施工、季节更替都在不断改变场景外观，因此需要持续更新而非一次性重建。

这四面墙并非彼此独立——内存墙限制了可训练的高斯数量，进而制约了质量墙的突破；计算墙阻碍了时效墙所需的快速更新。

本章围绕上述四个核心问题，系统梳理了 106 篇相关论文。第~\ref{sec:14_1}~节讨论城市场景重建，覆盖分治策略、几何约束、跨视角融合和语义辅助四条技术路线。第~\ref{sec:14_2}~节介绍在线与增量重建，涵盖流式重建、场景更新、变化检测和持续学习。第~\ref{sec:14_3}~节深入分层与 LOD（Level-of-Detail）表示，探讨八叉树、层次高斯和虚拟内存等核心数据结构。第~\ref{sec:14_4}~节分析分布式与并行训练，覆盖模型并行、数据并行和联邦学习三种范式。第~\ref{sec:14_5}~节讨论其他大规模方法，包括核外训练（out-of-core training）、多模态融合和特殊场景重建。第~\ref{sec:14_summary}~节总结全章。

\begin{figure}[htbp]
\centering
\resizebox{\textwidth}{!}{\begin{tikzpicture}[
  route/.style={rectangle, rounded corners=5pt, draw=#1, fill=#1!10,
    text width=3.2cm, minimum height=0.9cm, align=center, font=\small},
  method/.style={rectangle, rounded corners=3pt, draw=#1, fill=#1!6,
    font=\scriptsize, text width=3.0cm, align=center, minimum height=0.6cm},
  arr/.style={-{Stealth[length=4pt]}, techgray!60},
  >=Stealth
]

  \node[route=techgray, text width=4cm, font=\small\bfseries] (root) at (22.2, 5.6) {大规模与在线重建方法分类};
  \node[route=techblue] (b0) at (7.3, 3.1) {城市场景重建};
  \node[method=techblue] (b0_l0) at (1.5, 0.6) {分治重建\\{\tiny CityGaussian, BlockGaussian, LoBE-GS}};
  \node[method=techblue] (b0_l1) at (5.3, 0.6) {跨视角语义\\{\tiny Horizon-GS, GS4City, DRAGON}};
  \node[method=techblue] (b0_l2) at (9.2, 0.6) {外观优化\\{\tiny GS-Diff, ADM-GS, CityGS-X}};
  \node[route=techgreen] (b1) at (16.4, 3.1) {在线增量重建};
  \node[method=techgreen] (b1_l0) at (12.8, 0.6) {流式更新\\{\tiny Gaussian On-the-Fly, Splatfacto-W}};
  \node[method=techgreen] (b1_l1) at (16.4, 0.6) {增量融合\\{\tiny TideGS, MegaGS, FreeGS}};
  \node[route=techorange] (b2) at (22.2, 3.1) {分层与LOD};
  \node[method=techorange] (b2_l0) at (20.3, 0.6) {层次化结构\\{\tiny Octree-GS, HUG, Scaffold-GS}};
  \node[method=techorange] (b2_l1) at (24.1, 0.6) {核外渲染\\{\tiny Grendel, Virtual Memory GS}};
  \node[route=techpurple] (b3) at (27.7, 3.1) {分布式并行};
  \node[method=techpurple] (b3_l0) at (27.7, 0.6) {并行训练\\{\tiny Grendel, Distributed 3DGS}};
  \node[method=techpurple] (b3_l1) at (31.4, 0.6) {联邦协同\\{\tiny PhotoReg, GaussReg, SIREN}};
  \node[route=techgray] (b4) at (35.0, 3.1) {其他};
  \node[method=techgray] (b4_l0) at (35.0, 0.6) {核外训练\\{\tiny Out-of-core GS, VastGaussian}};
  \node[method=techgray] (b4_l1) at (38.3, 0.6) {多模态\\{\tiny Neural-MMGS, LiDAR-GS}};
  \draw[arr] (root) -- (b0);
  \draw[arr] (root) -- (b1);
  \draw[arr] (root) -- (b2);
  \draw[arr] (root) -- (b3);
  \draw[arr] (root) -- (b4);
  \draw[arr] (b0) -- (b0_l0);
  \draw[arr] (b0) -- (b0_l1);
  \draw[arr] (b0) -- (b0_l2);
  \draw[arr] (b1) -- (b1_l0);
  \draw[arr] (b1) -- (b1_l1);
  \draw[arr] (b2) -- (b2_l0);
  \draw[arr] (b2) -- (b2_l1);
  \draw[arr] (b3) -- (b3_l0);
  \draw[arr] (b3) -- (b3_l1);
  \draw[arr] (b4) -- (b4_l0);
  \draw[arr] (b4) -- (b4_l1);
\end{tikzpicture}}
\caption{大规模与在线重建方法分类}
\label{fig:ch14_large_scale_taxonomy}
\end{figure}

% ======================================================================
\section{城市场景重建}
\label{sec:14_1}

城市场景重建是大规模 3DGS 研究中论文数量最多、技术路线最丰富的方向。城市环境具有几个显著特点：场景尺度大（数平方公里）、视角跨度大（从无人机俯瞰到地面街景）、外观变化大（不同时段的光照差异）、动态元素多（车辆、行人）、弱纹理区域多（建筑墙面、道路表面）。这些特点使得直接将标准 3DGS 应用于城市场景效果不佳，催生了大量针对性方法。

\subsection{分治重建与几何约束}
\label{ss:14_1}

本节介绍两条技术路线：（1）分治策略与场景分割；（2）几何约束与表面重建。

\subsubsection{分治策略与场景分割}

本路线介绍大规模城市场景的分治重建方法，通过空间分割、分块训练和场景合并实现可扩展的高质量重建。

将大规模场景分割为可管理的子区域，分别训练后合并，是处理城市场景最主流的策略。这类方法需要回答三个核心问题：如何分割、如何训练、如何合并。

\paragraph{全局先验引导的分治训练与块级 LOD 渲染}

CityGaussian（2024.04）~\cite{2404_01133} 是城市场景分治重建的代表性工作，也是后续大量分治方法的基础参照。

\textbf{Baseline 与问题。}标准 3DGS（2024.03）~\cite{2403_10103} 在大规模场景中面临根本困难：场景需要超过两千万个高斯基元才能充分表达，单 GPU 显存远不足以容纳如此规模的参数（位置、旋转、缩放、不透明度、48 阶球谐系数），训练过程中频繁出现内存溢出。即使训练成功，渲染时对全部高斯做投影排序也极为缓慢。直观的分治策略是将 COLMAP 稀疏点云直接按空间划分后分别训练，但这种方式缺乏深度约束和全局几何感知。各子区域在独立训练中会产生大量几何不准确的浮游体（floaters）来过拟合块外区域，使得后续块间融合几乎不可行。

\textbf{改了什么。}CityGaussian 将大规模场景的分治重建拆解为三个有序阶段：\emph{全局先验生成}（Global Prior Generation）、\emph{基元与数据分块}（Primitives and Data Division）、\emph{分层细节渲染}（Level-of-Detail）。全局先验作为分割和训练的锚点，取代了对 COLMAP 点云直接分割的朴素策略；数据分配以结构相似性度量驱动，取代了简单空间投影分配；LoD 渲染通过块级细节选择，取代了对全部高斯的暴力渲染。

\textbf{实现原理。}\emph{第一阶段——全局先验生成：}使用全部图像以标准 3DGS 训练 30{,}000 次迭代，得到粗粒度但几何正确的全局高斯集合 $G_K = \{G_k | k = 1, \ldots, K\}$。该先验的作用是双重的：首先，它提供了场景整体几何分布的合理估计，后续块级微调时高斯点被引导至几何正确的位置，大幅减轻融合时的块间干扰；其次，相比 COLMAP 稀疏点云渲染的模糊图像，全局高斯能生成更干净、更准确的渲染结果，为下一步的数据筛选提供了可靠的评估基础。

\emph{第二阶段——基元与数据分块：}由于城市场景通常是无界的（unbounded），直接在原始三维空间做均匀网格划分会产生大量几乎为空的块，导致严重的负载不均衡。CityGaussian 先对全局高斯位置做归一化和收缩映射（contraction），将前景区域线性映射至 $[-1, 1]^3$，将背景区域非线性压缩至该立方体的外表面附近，使无界空间转化为有界立方体。在收缩空间上做均匀网格划分，即可得到高斯数量更为均衡的分块方案。

数据分配的核心问题是：对于第 $j$ 个块所包含的高斯子集 $G_{K_j}$，哪些相机视角 $\tau_i$ 应被纳入该块的训练？CityGaussian 设计了双条件分配策略。第一条件基于结构相似性：将完整全局高斯渲染的图像 $I_{G_K}(\tau_i)$ 与去掉第 $j$ 块高斯后渲染的图像 $I_{G_K \setminus G_{K_j}}(\tau_i)$ 做比较，若二者之间的 SSIM（Structural Similarity Index Measure，结构相似性度量）损失超过阈值 $\epsilon$，说明第 $j$ 块对渲染结果有显著贡献，该视角应被纳入：

\begin{equation}
B_1(\tau_i, G_{K_j}) = \begin{cases} 1, & L_{\text{SSIM}}(I_{G_K}(\tau_i),\; I_{G_K \setminus G_{K_j}}(\tau_i)) > \epsilon \\ 0, & \text{otherwise} \end{cases}
\end{equation}

第二条件处理块边缘的"外向观察"视角——当相机位于某块内部但朝向块外时，仅靠 SSIM 条件无法捕获这些视角，因此额外纳入相机位置落入块内的视角：

\begin{equation}
B_2(\tau_i, G_{K_j}) = \begin{cases} 1, & b_{j,\min} < \text{contract}(\mathbf{p}_{\tau_i}) < b_{j,\max} \\ 0, & \text{otherwise} \end{cases}
\end{equation}

最终分配为二者的并集 $B(\tau_i, G_{K_j}) = B_1 \lor B_2$。若某块的高斯数量仍然过少，则自动扩展其空间边界直到满足最低训练数据要求，以防止过拟合。

分块完成后，每个块以全局先验高斯作为初始化，在原始未收缩空间中独立微调，损失函数为 $\mathcal{L}_1$ 和 $\mathcal{L}_{\text{SSIM}}$ 的加权和。得益于全局几何先验的引导，块间干扰被大幅抑制，最终融合只需简单地拼接（concatenation）各块内属于自身空间边界的高斯即可。

\emph{第三阶段——分层细节渲染：}大规模场景渲染的核心瓶颈是每次视角切换都需处理全部高斯。CityGaussian 利用 LightGaussian 压缩策略，从训练好的完整高斯模型出发，生成多个不同容量的细节层次（更少点数、更低特征维度、更低数据精度）。渲染时采用块级细节选择而非逐点选择：对每个块，用中位绝对偏差（Median Absolute Deviation, MAD）算法估计其鲁棒空间边界，避免浮游体对边界估计的干扰；然后计算该块包围盒与视锥体的交截，根据块中心到相机的最小距离选取对应细节层次的高斯集合，直接拼接送入光栅化器。这一块级策略避免了逐点距离计算的开销，同时 MAD 鲁棒边界估计有效排除了异常高斯对块体积的虚增。

\textbf{设计动机。}CityGaussian 的三项关键设计各有其深层考量。全局先验的必要性源于一个观察：COLMAP 稀疏点云的渲染结果模糊且不准确，无法可靠地评估特定视角对某子区域的贡献程度。而经过 30K 次迭代训练的粗粒度全局高斯能提供干净得多的渲染结果，使后续的 SSIM 数据筛选变得有意义。收缩映射的引入则是因为城市场景的无界特性——远处天空和建筑的高斯分布极为稀疏，均匀网格划分会导致大量空块，而收缩映射可以将有效高斯密度重新均衡化。块级 LoD 选择而非逐点选择的动机是计算效率：在数千万高斯的场景中，逐点距离计算和分配的代价不可接受，块级策略则将决策粒度从数千万降至数十至数百。

\textbf{关键参数。}全局先验训练迭代次数 30{,}000；SSIM 分配阈值 $\epsilon$；空间划分的网格分辨率（块的数量）；LoD 的细节层次数和压缩比率；MAD 鲁棒边界估计的超参数 $n_{\text{MAD}}$。实验在 Mill-19 和 UrbanScene3D 等大规模数据集上验证，16 GPU 并行训练可实现当时最优渲染质量，并支持跨尺度的实时渲染。

从上述描述可以看出，CityGaussian 的核心思想是"先全局、后分块、再分层"——全局先验确保几何正确性，空间分块实现并行训练，块级 LOD 保证渲染效率。这一三阶段范式为后续大量分治方法奠定了框架基础。

\paragraph{分解梯度致密化与拉伸过滤器的几何精度突破}

CityGaussianV2（2024.11）~\cite{2411_00771} 将大规模场景重建从"渲染优先"推进到"渲染与几何并重"的新阶段。以 2D Gaussian Splatting 为基础基元，通过分解梯度致密化（DGD）、拉伸过滤器和精简训练流水线三项核心改进，在保持与 CityGaussian 相当渲染质量的同时，F1 几何分数提升约 11\%，训练时间缩短至少 25\%，内存占用降低 50\% 以上。

\textbf{基线做法。}CityGaussian V1 以 3DGS 为基元，虽然渲染质量出色，但 3D 椭球体的体积表示导致表面定位模糊，提取的 mesh 质量远不能满足实际需求。若将 2DGS 直接替换进 CityGaussian 的框架，会遭遇两个致命问题。

（1）\emph{收敛困难}：2DGS 在训练早期比 3DGS 更容易产生模糊重建。这是因为 2DGS 默认的致密化梯度来源于总损失 $\mathcal{L} = (1-\lambda)\mathcal{L}_1 + \lambda \mathcal{L}_{\text{SSIM}}$，而 $L_1$ RGB 损失对模糊不敏感，不优先保持结构完整性，导致致密化驱动力不足。

（2）\emph{高斯爆炸}：2D surfel 从远处观察时投影面积极小（小于 1 像素），经抗锯齿低通滤波后协方差被替换为固定值，无法通过有效梯度调整自身的缩放和旋转。在分块微调中，这些退化的小点被分配到每个块的少量视角频繁观察到，梯度快速累积，触发高斯数量的指数级增长，最终导致内存溢出。

\textbf{改了什么。}为解决上述问题，CityGaussianV2 提出三项核心改进。

（1）\emph{分解梯度致密化（Decomposed-Gradient-based Densification, DGD）}：将致密化梯度从总损失梯度重构为以 SSIM 梯度为主导的形式，优先利用对结构信息高度敏感的 SSIM 损失梯度驱动致密化，有效解决早期训练模糊问题。

（2）\emph{拉伸过滤器（Elongation Filter）}：在每次致密化前检查每个 surfel 的长短轴比率，排除极端拉伸的退化 surfel，从源头切断"退化小点→高梯度→致密化爆炸"的正反馈循环。

（3）\emph{精简训练流水线}：将球谐阶数从 3 降至 2，并在块级微调中引入基于百分位贡献评分的在线剪枝（trimming），绕过 CityGaussian 的后处理剪枝和蒸馏阶段，大幅节省存储和内存。

此外，还引入了基于 Depth-Anything-V2 的深度先验监督作为几何优化的辅助引导。

\textbf{实现原理。}

\emph{分解梯度致密化（DGD）：}标准 3DGS 的致密化梯度为位置梯度的范数 $\nabla_{\text{densify}} = \|\nabla \mathcal{L} / \partial \boldsymbol{\mu}_n\|$。DGD 将其重构为：

\begin{equation}
\nabla_{\text{densify}} = \max\!\bigl(|\nabla \mathcal{L}_1|_{\text{avg}},\; |\nabla \mathcal{L}_{\text{SSIM}}|_{\text{avg}},\; 1\bigr) \cdot \nabla \mathcal{L}_{\text{SSIM}}
\end{equation}

其中 $|\nabla \mathcal{L}_{\text{SSIM}}|_{\text{avg}}$ 将 $\nabla \mathcal{L}_{\text{SSIM}}$ 自动缩放到与原始总损失梯度同量级，使其能与原始致密化阈值直接兼容，无需额外调参。$\max$ 函数中包含常数 $1$ 确保梯度下界不为零。核心思想是 SSIM 损失对结构模糊高度敏感，将其梯度作为致密化的唯一驱动力，迫使致密化过程优先关注结构完整性缺失的区域。

\emph{拉伸过滤器（Elongation Filter）：}对每个 surfel 定义拉伸率：

\begin{equation}
\rho_n = \frac{s_{n,u} - s_{n,v}}{s_{n,u} + s_{n,v}}
\end{equation}

其中 $s_{n,u}$ 和 $s_{n,v}$ 分别为 surfel 的两个切向缩放分量。$\rho_n$ 接近 1 意味着极端拉伸（一个轴远长于另一个轴）。在每次致密化操作前，检查所有 surfel 的 $\rho_n$，将 $\rho_n$ 高于某阈值的 surfel 排除在克隆和分裂操作之外。消融实验表明，该过滤器在预训练阶段对性能几乎没有影响（因为正常视角下退化点占比低），但在分块微调阶段能有效防止高斯数量的指数级膨胀。

\emph{深度先验监督：}利用 Depth-Anything-V2 估计逆深度图 $\hat{D}_k$，对齐到数据集尺度后记为 $D_k$，深度监督损失为 $\mathcal{L}_{\text{Depth}} = |\hat{D}_k - D_k|$。损失权重 $\lambda_{\text{depth}}$ 在训练过程中从 0.5 指数衰减至 0.0025，逐步抑制不完美深度估计的负面影响。

\emph{精简训练流水线：}（i）球谐降阶——从头使用 SH degree 2（特征维度从 48 降至 27），绕过后处理蒸馏步骤，存储和内存额外降低超过 25\%，速度提升 4.2 FPS，对渲染质量影响可忽略；（ii）在线剪枝——在块级微调中引入基于百分位贡献评分的剪枝（trimming），单视角下第 $n$ 个高斯的贡献定义为：

\begin{equation}
C_{n,k} = \frac{1}{|P_k|}\sum_{p \in P_k} \alpha_n(p) \prod_{j=1}^{n(p)-1}(1-\alpha_j)^{1-\gamma}
\end{equation}

其中 $P_k$ 为第 $k$ 个视角下第 $n$ 个高斯的 2D 投影区域，$n(p)$ 为其深度排序，$\gamma=0.5$。对分配给第 $m$ 块的所有视角取平均得到 $C_n$，以百分位阈值（默认 pruning ratio 0.025，激进设置 0.1）决定是否丢弃。该方法在训练中定期评估贡献并自动移除冗余和未被观察到的高斯，无需后处理剪枝阶段。最终对合并后的高斯执行 vectree 量化（码本大小 8192，量化比率 0.4），关键属性的 SH 经激进的向量量化压缩，其余属性以 float16 存储，实现约 75\% 的存储压缩。

\textbf{设计动机。}CityGaussianV2 的每项设计都有明确的针对性。DGD 的动机源于一个关键实证发现：2DGS 在早期训练中比 3DGS 更容易产生模糊，而消融实验证实致密化梯度的来源中，SSIM 梯度的参与最为关键——$L_1$ 损失对模糊不敏感，无法提供足够的致密化驱动力。拉伸过滤器的设计动机来自对 2D surfel 几何特性的深入分析：2D 高斯从远处投影时面积趋近于零，在抗锯齿低通滤波下协方差被固定值替换，丧失了通过梯度自我修正的能力。这些"僵尸"surfel 在分块微调的稀疏视角下反而积累了高梯度，是内存爆炸的根本原因。球谐降阶的动机是简化训练流水线：CityGaussian V1 需要 30K 次额外剪枝迭代和 10K 次蒸馏迭代，耗时且占内存；直接从头使用低阶 SH 配合在线剪枝，可将后处理完全内化到训练过程中。深度先验的引入则基于一个物理直觉：GS 算法需要迭代优化来从每个视角的单目线索中消除歧义，最终收敛到一致的 3D 几何。深度先验在训练早期提供强有力的几何引导，加速收敛；而指数衰减的权重则确保训练后期不会受限于不完美的深度估计。此外，论文还建立了大规模场景的几何评估基准——借鉴 Tanks and Temples 的评估协议，包括点云对齐、重采样、体边界裁剪和 F1 分数测量，通过可见频率过滤边界效应，将 F1 分数的误差棒长度从 0.1 降至 0.003。

\textbf{关键参数。}DGD 的梯度缩放因子 $\beta = 0.9$；拉伸过滤器阈值（$\rho_n$ 的下界）；深度先验损失权重从 0.5 指数衰减至 0.0025；表面法线损失权重 0.0125（原始 2DGS 的 1/4，降低对复杂场景的过强约束）；深度畸变损失权重设为 0（实证发现对性能有害）；在线剪枝的百分位比率默认 0.025，激进设置 0.1；球谐阶数固定为 2；vectree 量化码本大小 8192，量化比率 0.4；mesh 提取使用 TSDF（Truncated Signed Distance Function）算法，体素大小 1m，SDF 截断距离 4m。所有实验在 8 块 A100 GPU 上进行。在 GauU-Scene 上，Residence 场景分为 42 个块，Russian Building 和 Modern Building 各 33 个块。实验表明，DGD 策略可提升 1.0 PSNR、0.04 SSIM 和近 0.02 F1 分数；完整流水线（ours-s）训练时间缩短 25\%、内存降低超过 50\%。从 7K 次迭代结果出发训练的紧凑版本（ours-t）可将总训练时间从 3 小时降至 2 小时，模型大小低于 300 MB，适合部署到智能手机或 VR 头显等低端设备。

\paragraph{内容感知二叉树分割与伪视图几何约束}

BlockGaussian（2025.04）~\cite{2504_09048} 在分治框架的三个核心环节（分割、训练、合并）上均提出了新的解决方案，特别是针对"分割不均衡""监督不匹配""空间浮游体"三个关键问题给出了系统性回应。

\textbf{Baseline 与问题。}在 BlockGaussian 之前，CityGaussian 等方法在收缩空间上做均匀网格划分，虽然缓解了无界场景的空块问题，但未考虑场景内容复杂度的空间差异。建筑密集区域和空旷道路被等同对待，导致各块的重建难度和训练时间差异巨大。此外，分块训练存在两个固有缺陷：（1）\emph{监督不匹配}——某块的训练视角只能看到该块高斯渲染的部分图像，与完整真值图像做损失计算时，图像中不属于该块的区域会引入错误的梯度信号；（2）\emph{空间浮游体}——各块独立优化时，缺乏对其占据空间以外的区域（如天空、远距离建筑之间的大段空气空间）的监督，导致大量浮游体产生，合并后严重降低渲染质量。

\textbf{改了什么。}BlockGaussian 提出三项核心改进：（1）基于稀疏点云密度分布的内容感知二叉树分割，取代均匀空间划分；（2）辅助高斯 + 小批量优化策略，解决监督不匹配问题；（3）伪视图几何约束损失，在不引入额外训练视角的前提下对空气空间施加间接深度监督，消除浮游体。

\textbf{实现原理。}\emph{内容感知分割（Content-Aware Scene Partition）：}BlockGaussian 利用 SfM 稀疏点云的密度分布作为场景复杂度的估计器。具体流程为：首先基于 Manhattan 世界假设估计地面法向量并沿 $y$ 轴对齐，将稀疏点云投影至 $x$-$z$ 平面并手动定义重建感兴趣区域（Region of Interest, RoI）；然后以二叉树结构递归地分割场景——给定最大树深 $M$ 和叶节点最大点数阈值 $N_b^t$，当某节点深度 $d < M$ 且包含的点云数量 $N_b > N_b^t$ 时，沿该节点包围盒的最长边一分为二，生成两个子节点；否则标记为叶节点。这种自适应分割使得高密度区域（复杂建筑群）被细分为更多小块，低密度区域（开阔路面）保持为大块，自然实现了负载均衡。

视角分配则基于 SfM 可见性信息：对每个块统计其边界内的点云数量 $N_b$，对每张训练视角统计其可见的关键点总数 $N_v$ 及其中落入该块边界的数量 $N_b^v$，若比值 $N_b^v / N_v$ 超过阈值（默认 0.3），则将该视角分配给该块。

\emph{可见性感知块优化（Visibility-Aware Block Optimization）：}分割后，每块只包含自身的块高斯 $G_b = \{g_b^n\}$，但训练视角的真值图像覆盖范围远超块边界。仅用 $G_b$ 渲染的图像 $I_i^b$ 只能看到块内区域，与完整真值 $I_i^{gt}$ 计算损失时会产生错误监督。BlockGaussian 引入\emph{辅助高斯} $G_a = \{g_a^m\}$，由分配给该块的训练视角中、落在块范围外的稀疏点云初始化，负责建模块外可见场景。渲染时联合使用 $G_b$ 和 $G_a$：

\begin{equation}
I_i^r = \mathcal{R}(R_i, t_i, G_b, G_a)
\end{equation}

辅助高斯的梯度同样参与反向传播，但只对块高斯 $G_b$ 执行致密化（densification）操作，不对 $G_a$ 致密化，以减少冗余。此外，引入深度先验损失增强几何质量：

\begin{equation}
\mathcal{L}_{\text{depth}}(D^{e}, D^r) = \mathcal{L}_1\!\left(\frac{1}{D^{e}},\; \frac{1}{D^r}\right)
\end{equation}

其中 $D^{e}$ 由 DepthAnythingV2 估计，$D^r$ 为渲染深度，尺度对齐后在逆深度空间计算 $\mathcal{L}_1$ 损失。深度先验在训练早期提供强有力的几何引导，加速收敛；而在训练后期指数衰减的权重确保不会受限于不完美的深度估计。由于辅助高斯 $G_a$ 的监督严重不足，直接应用标准 3DGS 优化会导致其退化并干扰 $G_b$ 的优化。因此 BlockGaussian 采用小批量（mini-batch）策略，累积多个视角的梯度后再做参数更新，以提升梯度稳定性。

\emph{伪视图几何约束（Pseudo-View Geometry Constraint）：}这是 BlockGaussian 最具特色的贡献。各块独立训练时，空气空间中缺乏直接监督信号，导致浮游体在合并后暴露。直观地说，空气空间中没有物体表面，光度损失不会对空气中的高斯产生约束——它们可以自由漂浮在任何位置而不影响渲染图像的像素颜色。解决这一问题的思路不是引入更多训练视角（代价高且可能不可行），而是利用已有视角的渲染深度构造自监督信号。具体做法是：对每个训练视角（参考视角），通过微扰其相机位置生成伪视角。扰动量 $\Delta t$ 在视差空间中控制，以保证扰动幅度与场景深度自适应匹配：

\begin{equation}
\Delta t = \left[\frac{\text{median}(D_{\text{ref}}^r) \cdot \delta_p}{f},\; 0,\; 0\right]^T
\end{equation}

其中 $\delta_p$ 为视差扰动超参数，$f$ 为相机焦距。利用伪视角的渲染深度图 $D_{\text{pse}}^r$，将伪视角渲染图像 warp 回参考视角，得到重投影图像 $I_{\text{warp}}^r$ 和有效掩码 $M$。伪视图几何约束损失为：

\begin{equation}
\mathcal{L}_{\text{pse}} = \frac{1}{|M|} \sum_{M} \mathcal{L}_1(I_{\text{ref}}^{gt},\; I_{\text{warp}}^r)
\end{equation}

该损失的巧妙之处在于：如果空气空间中存在浮游体，其渲染深度会偏离真实表面深度，导致 warp 后的图像与真值不匹配，从而产生惩罚梯度。这实现了对渲染深度的间接监督，无需额外采集数据即可有效消除浮游体。

\emph{场景合并与渲染：}训练完成后，裁剪掉辅助高斯 $G_a$，直接拼接所有块高斯即可得到完整场景表示。得益于伪视图几何约束对浮游体的有效抑制，合并后的渲染质量得到保障。

\textbf{设计动机。}内容感知分割的动机是观察到的负载不均衡问题：均匀划分下不同块的训练时间可相差数倍，多 GPU 并行效率低下。辅助高斯的引入源于一个被忽视的事实——分块训练的损失计算本身就有问题。如果渲染只覆盖部分图像区域，与完整真值比较必然引入噪声梯度。伪视图几何约束的灵感来自多视角立体视觉中视角扰动的一致性约束，但其关键区别在于不需要真实的多视角数据，而是利用渲染深度构造自洽的几何约束信号。

\textbf{关键参数。}二叉树最大深度 $M$；叶节点最大点数阈值 $N_b^t$；视角分配比例阈值（默认 0.3）；视差扰动超参数 $\delta_p$；小批量优化的累积视角数。该方法在多个大规模场景基准上实现了平均 1.21 dB 的 PSNR 提升和 5 倍优化加速。

\paragraph{八叉树层级神经高斯与层次化加权训练}

HUG（2025.04）~\cite{2504_16606} 专注于大规模航拍城市场景，提出基于可见性的高效数据分割方法和层次化加权训练策略。其可见性分割通过统计每张图像中各块内稀疏点云的重投影可见点数进行训练视角分配，仅需 1 分钟即可完成（CityGS 约需 2 小时），同时基于 SVM（Support Vector Machine）可见性掩码使各块聚焦于区域内像素训练。在层次化表示方面，HUG 基于 Scaffold-GS 和 Octree-GS 构建\emph{八叉树层级神经高斯}架构——锚点初始化于覆盖三维空间的八叉树节点中心，八叉树最大层级 $K$ 由块内所有点到相机中心距离的 0.95 和 0.05 分位数之比确定：$K = \lfloor\log_2(d_{\max}/d_{\min})\rfloor + 1$。每个训练视角下，仅选择预测层级 $L \leq \hat{L}$ 的锚点参与优化，其中 $\hat{L} = \text{clip}(\lfloor\log_2(d_{\max}/d)\rfloor, 0, K-1)$，$d$ 为锚点到相机中心的距离，实现距离自适应的细节选择。

层次化加权训练是 HUG 的核心创新。该方法对八叉树中不同 LOD 层级的锚点分别渲染并施加指数衰减权重的独立监督 $\ell_{\text{hier}} = \alpha \ell_1(\text{gt}, R_{\text{full}}) + (1-\alpha)\ell_s + \sum_{k=1}^{K}(\alpha/2^{K+1-k}\,\ell_1 + (1-\alpha)/2^k\,\ell_s)$，确保高层粗锚点专注于近似色彩、低层细锚点逐步恢复细节，并配合不透明度掩码（保留不透明度 $> 0.5$ 的像素区域）处理高层渲染中的大面积透明区域，避免空白区域对梯度的稀释。HUG 还提出动态梯度阈值衰减策略（阈值每 $M=5000$ 次迭代按 $\tau_g(i) = \tau_g \cdot \gamma^{i/M}$ 衰减，默认 $\gamma = 0.8$）和基于梯度的锚点层级跃迁机制 $L' = \text{clip}(K-1,\; L + 0.01 \cdot \mathbb{I}(\nabla_g > \beta \cdot \tau_g(i)))$（$\beta = 4$），使锚点在训练中自适应调整所属 LOD 层级。

在场景融合阶段，HUG 设计了贪心投票锚点重滤波策略——即使锚点中心位于块外，若其生成的神经高斯中有超过一半的中心落入块内，该锚点仍被保留，解决了朴素空间裁剪导致的边界伪影问题。此外，由于不同块的神经高斯由各自独立训练的小型 MLP 生成，HUG 重写了光栅化管线以支持多 MLP 并行加载。渲染时根据锚点所属块动态选择对应 MLP 生成神经高斯属性，实现了跨块无缝渲染。在四个真实数据集（Mill-19 Building/Rubble、UrbanScene3D Residence/Sci-Art）和一个合成数据集（MatrixCity）上达到当时最优结果。

\paragraph{负载均衡 KD-Tree 与 DepthBackProj 快速分配}

LoBE-GS（2025.10）~\cite{2510_01767} 系统性地诊断了已有分治方法的效率瓶颈，从分割策略、相机分配、块级训练三个层面提出针对性改进，在不牺牲重建质量的前提下将端到端训练时间缩短最多 2 倍。

\textbf{Baseline 与问题。}以 CityGaussian 为代表的粗到细分治管线存在两个关键效率瓶颈。

\emph{第一，负载不均衡。}CityGaussian 在收缩空间上做均匀网格划分，虽然缓解了无界场景的空块问题，但未考虑相机分布的空间异质性。建筑密集区域集中了大量重叠视角，空旷道路只有稀疏的航拍覆盖，导致不同块的优化需求差异巨大。在并行训练中，最慢的块决定了整体训练时间（即长尾瓶颈，straggler effect），均匀划分无法缓解这一问题。

\emph{第二，相机分配开销过大。}CityGaussian 对 $M$ 个块和 $N$ 个相机逐一计算 SSIM：对每个块，先滤除块外高斯后渲染，再与完整渲染结果比较——这需要至少 $(M+1)N$ 次投影操作，在实验中可占端到端运行时间的近一半。

\textbf{改了什么。}LoBE-GS 在分治管线的三个阶段分别引入改进：（1）\emph{负载均衡 KD-Tree 场景分割}——取代均匀网格划分，通过最小最大化目标精化切割线位置，使各块的预估工作量趋于均衡；（2）\emph{DepthBackProj 快速相机分配}——取代 SSIM 渲染比较，利用粗模型深度反投影实现单次深度渲染即可完成全部分配；（3）\emph{可见性裁剪与选择性致密化}——两项轻量技术分别减少块内无关高斯和不必要的致密化操作，降低每块的训练成本。

\textbf{实现原理。}\emph{负载均衡 KD-Tree 场景分割：}LoBE-GS 以 KD-Tree 的叶节点作为最终的分块单元，每个节点存储轴对齐包围盒 $\Omega^{(b)}$ 和分配的相机集合 $\mathcal{V}^{(b)}$。工作负载定义为该区域分配的相机数量：

\begin{equation}
L(\Omega^{(b)}) = |\mathcal{V}^{(b)}|
\end{equation}

KD-Tree 的构建采用\emph{最佳优先}（best-first）策略：每次分裂当前负载最大的叶节点，而非深度优先递归。这一策略在异构相机分布下表现更优——它优先分解高负载区域，确保最终叶节点的负载方差最小化。分裂轴的选择基于几何启发式（取包围盒最长边），以避免产生退化的狭长块并保持空间紧凑性。

分割线的精化是负载均衡的核心机制。初始化时，切割位置设为粗模型高斯中心在分裂轴上的中位数投影，以获得合理的空间分离。然后通过最小最大化目标精化切割线：

\begin{equation}
q^* = \arg\min_{q \in [q_{\min},\, q_{\max}]} \max\!\big(L(\Omega_L(q)),\; L(\Omega_R(q))\big)
\end{equation}

其中 $q_{\min} = 0.1$，$q_{\max} = 0.9$ 以防止退化的极端分割。当切割线 $q$ 从左向右移动时，左子区域扩大、捕获更多相机，右子区域缩小、相机减少，两侧负载呈近似单调的此消彼长关系。因此，LoBE-GS 采用轻量的\emph{二分搜索}迭代：每次将切割线向负载较重的一侧移动，直到两侧工作量趋于可比。该精化过程在每个分裂节点独立执行，逐层减少全局最大工作量。

\emph{DepthBackProj 快速相机分配：}该方案将 SSIM 渲染比较的 $O(MN)$ 复杂度降至 $O(|\mathcal{C}|)$（相机总数）。具体流程分三步：首先，对每张训练图像，利用粗模型通过 $\alpha$-混合 方程渲染深度图 $D$：

\begin{equation}
D = \sum_{i=1}^{|N|} d_i \alpha_i \prod_{j=1}^{i-1}(1-\alpha_j)
\end{equation}

其中 $d_i$ 为第 $i$ 个交点深度，$\alpha_i$ 为其不透明度。然后，将深度图反投影至三维空间，生成密集点云 $\mathcal{P}^{(c)} = \{p_{c,k}\}_{k=1}^{N_p}$。最后，对每个块 $b$ 计算相机 $c$ 的可见性比率：

\begin{equation}
V_{c,b} = \frac{1}{N_p}\sum_{k=1}^{N_p} \mathbb{1}[p_{c,k} \in \Omega^{(b)}]
\end{equation}

当 $V_{c,b} > \tau$（默认 $\tau = 0.1$）时，将该相机分配给块 $b$。此外，在 KD-Tree 递归分裂过程中，子节点仅继承父节点已分配的相机子集并重新计算可见性比率，无需回溯评估全部相机，使分裂和精化的开销与父节点相机集大小成正比，实现高效的递归分割。

\emph{可见性裁剪（Visibility Cropping）：}标准粗到细管线在块级微调时加载完整粗模型，其中大量高斯位于当前块边界之外且不被任何分配相机观察到。虽然视锥体剔除（frustum culling）在光栅化阶段排除了不可见高斯，但 Adam 优化器仍为所有加载的高斯维护动量状态并执行参数更新，造成不必要的内存和计算开销。可见性裁剪仅保留被分配相机中至少一个可见的高斯，大幅减少参与优化的高斯数量。这与朴素的空间裁剪（仅保留中心在块内的高斯）不同——块外但被分配相机可见的高斯仍被保留，以确保渲染监督的完整性。

\emph{选择性致密化（Selective Densification）：}可见性裁剪保留了部分块外高斯以支持渲染监督，但这些高斯在合并阶段会被丢弃，允许其参与致密化（生成新高斯）是多余的。选择性致密化将致密化操作限制在块内高斯上：块外高斯仍可参与渲染和损失计算贡献梯度，但不会分裂或克隆产生新的高斯。这在不影响重建保真度的前提下减少了内存使用和优化开销。

\textbf{设计动机。}LoBE-GS 的各项设计决策背后有明确的实证支撑。

\emph{为什么用相机数量而非块面积或可见点数作为工作负载代理？}论文进行了系统性的代理预测分析：按块优化时间排序后，在 Top-5\% 和 Top-10\% 最重负载块的召回率上，相机数量（0.529/0.471）显著优于块面积（0.353/0.412）和可见点数（0.117/0.265）。这一结果的直觉解释是：块级微调的工作量与贡献梯度和触发致密化事件的训练视角数强相关，相机数量最直接地反映了这一负载来源。

\emph{为什么需要切割线精化？}即使使用 KD-Tree 取代均匀网格，朴素的中位数分割在异构相机分布下仍会产生显著的工作量失衡。精化通过沿切割轴微调分割位置，逐层压缩最慢块的时间。

\emph{为什么采用最佳优先而非深度优先构建？}深度优先递归会先完全展开一侧子树，可能导致另一侧的高负载节点迟迟得不到分裂。最佳优先每次选择全局最大负载叶节点，确保分裂资源始终投入在最需要均衡化的位置。

\textbf{关键参数。}KD-Tree 切割线搜索范围 $[q_{\min}, q_{\max}] = [0.1, 0.9]$；相机可见性分配阈值 $\tau = 0.1$；块配置与 CityGaussian 一致以便公平对比（MatrixCity-Aerial 36 块，Building 20 块，Residence 20 块，Rubble 9 块，Sci-Art 9 块）。实验在 10 节点 $\times$ 8 NVIDIA L40 GPU 集群上评测，LoBE-GS 在五个大规模场景（Mill-19 Building/Rubble、UrbanScene3D Residence/Sci-Art、MatrixCity-Aerial）上达到与 CityGaussian 和 VastGaussian 可比的 PSNR/SSIM，同时在大多数场景上 LPIPS（Learned Perceptual Image Patch Similarity）更优；端到端训练时间（$T_{\text{e2e}} = T_{\text{coarse}} + T_{\text{partition}} + \max_b T_{\text{fine}}^{(b)}$）相比基线提升最多 2 倍，其中分割时间 $T_{\text{partition}}$ 的降低尤为显著。消融实验表明，移除可见性裁剪一致增加最大块训练时间和峰值显存，移除选择性致密化增加训练时间但不影响质量，移除切割线精化在所有场景上增加最慢块时间，替换 DepthBackProj 为 SSIM 相机选择则使分割时间急剧上升但不影响训练时间——说明四项改进各司其职、互为补充。

\paragraph{属性视角图聚类与遮挡感知分割}

OccluGaussian（2025.03）~\cite{2503_16177} 关注到现有分割方法忽略遮挡关系——每个子区域可能包含严重遮挡区域，导致区域内的相机相关性低、对整体重建贡献小。该方法提出遮挡感知场景分割策略，其核心是构建\emph{属性视角图}（attributed view graph）$G = (V, E, X)$：节点 $V$ 对应训练相机，边权 $a_{ij}$ 为 SfM（运动恢复结构，Structure from Motion）特征匹配点数（反映共可见性），节点特征矩阵 $X$ 包含相机三维坐标的位置编码。

分割阶段采用属性图聚类算法。先对特征矩阵执行图卷积平滑（滤波器 $G = (I - \frac{1}{2}L_s)^r$，$L_s$ 为对称归一化拉普拉斯矩阵），再对平滑后的相似度矩阵 $\tilde{S} = \frac{1}{2}(|H| + |H^T|)$（$H = \tilde{X}^T\tilde{X}$）进行谱聚类。聚类数通过自适应精化确定：初始值 $K=10$，然后递归分裂过大簇、忽略过小簇或被完全覆盖的簇，直到各簇相机数收敛至 $[M_c - \epsilon_c M_c,\; M_c + \epsilon_c M_c]$（$\epsilon_c = 0.5$）。边界线由 SVM 线性分类器拟合，因遮挡感知聚类而天然对齐场景物理布局。

在训练相机选择方面，OccluGaussian 区分三组相机：\emph{基础集}（位于区域内的相机）、\emph{扩展集}（区域外但对区域内有充分可见内容的相机）和\emph{边界集}（朝向区域但被遮挡的相机）。边界集的引入至关重要——缺少这些相机时，边界附近的高斯会变得过大或拉长，导致浮游体和伪影。合并时裁剪掉中心位于区域外的高斯后拼接。

同时提出的\emph{区域化渲染}技术通过预计算可见性掩码 $M_j = \{m_i^j\}_{i=1}^{N_g}$（$m_i^j = 1$ 当且仅当高斯 $i$ 在区域 $R_j$ 任一训练视角下的渲染贡献超过 0.01），在渲染时仅加载当前视点所在区域的可见高斯子集，将遮挡关系转化为显式剔除。进一步的\emph{区域细分}将每个区域分为内部子区域和若干边界子区域（向内收缩 $0.1 d_{\max}$），生成更紧凑的可见性掩码，从而在不损失质量的情况下显著加速渲染。实验表明 OccluGaussian 在自建的 OccluScene3D 遮挡数据集、Zip-NeRF 数据集及航拍数据集上均取得最优重建质量和渲染速度。

\paragraph{相机连通图分割与渐进式渲染}

TraGraph-GS（2025.06）~\cite{2506_08704} 解决了已有分治方法难以泛化到任意相机轨迹的问题。刚性空间分割技术在面对不同飞行轨迹时表现不佳，子区域合并时的高斯重叠也会导致纹理细节失真。TraGraph-GS 的场景分割基于\emph{相机连通图} $G = (V, E, w)$ 构建：以图像为顶点，特征匹配数为边权，匹配数越高表示两幅图像的共可见区域越广。分割采用 BFS（Breadth-First Search）图分割算法——选取度最高的未分配节点作为种子，通过广度优先搜索扩展至 $|V|/k$ 个节点形成初始区域。对剩余节点计算其与各区域的得分 $s_i(u) = \sum_{v \in P_i} w(u,v)$，并分配至得分超过中位数的区域中节点数最少的一个，同时保证连通性和各区域节点数的均衡性。相比谱分解（节点分布不均易导致内存溢出）和 Kernighan-Lin 算法（不保证区域连通性），BFS 分割在这两个维度上取得了平衡。

在正则化约束方面，TraGraph-GS 利用图结构为无序图像自动识别邻域帧，引入两项关键约束：（1）\emph{多视角约束}——对每个像素建立空间平面 $(d, \mathbf{n})$，通过单应变换将参考视角的局部图块（patch）投影到邻域视角，以归一化互相关（Normalized Cross-Correlation, NCC）衡量颜色一致性并加入优化目标 $\mathcal{L}_{mv} = \sum_{p \in P_i}(1 - \text{NCC}(p, p'))$，显著改善弱纹理区域的几何精度；（2）\emph{多尺度约束}——根据高斯到相机中心的距离 $\|\boldsymbol{\mu} - \mathbf{c}\|_2$ 与场景半径 $r$ 的比值自适应调整致密化策略：近距离（$< 2r$）使用小高斯克隆，远距离使用大高斯分裂，并移除 $\max(s) > b(\cdot)$ 的过大高斯，兼顾近景细节和远景质量。

TraGraph-GS 还提出\emph{渐进式渲染策略}解决高斯重叠伪影：首先训练一个下采样点云和图像的粗糙全局高斯（跳过致密化以节省内存），渲染时分别用局部高斯和全局高斯渲染新视角，计算匹配得分 $S = (1-\lambda)\mathcal{L}_1 + \lambda\mathcal{L}_{SSIM}$ 对局部高斯排序；然后按贡献从高到低逐步叠加，当像素位置已渲染局部高斯的累积不透明度超过阈值 $\theta$（默认 0.5）时停止，避免低贡献区域的高斯引入浮游体。在航拍数据集（Mill-19、UrbanScene3D）和地面数据集（SmallCity、Campus、自建 Road/General）上，相比之前最优方法分别平均提升 1.86 dB 和 1.62 dB PSNR。

\paragraph{城市数据集与基准}

MatrixCity（2023.09）~\cite{2309_16553} 是城市场景神经渲染领域的里程碑式数据集。

\textbf{动机。}此前的大规模城市场景数据集要么规模不足（Mill-19 仅覆盖单栋建筑），要么缺少视角多样性（仅有航拍或仅有街景），无法全面评测城市级新视角合成方法。该数据集利用 Unreal Engine 5 的 City Sample 项目，开发了灵活的数据采集管线，支持航拍和街景视角，包含 67K 航拍图像和 452K 街景图像，覆盖总面积 $28\,\text{km}^2$ 的两座城市。MatrixCity 不仅提供了灵活的环境控制（光照、天气、人群密度），还建立了全面的基准评测，揭示了城市场景神经渲染的独特挑战。TUM2TWIN（2025.05）~\cite{2505_07396} 则提供了首个综合性的多模态城市数字孪生基准数据集，包含地理参考的室内外采集、多传感器数据（地面、移动、航拍、卫星），覆盖约 $100{,}000\,\text{m}^2$、767 GB 数据，支持从新视角合成到太阳能潜力分析等多种下游任务。

\subsubsection{几何约束与表面重建}

本路线介绍通过外部几何先验增强城市场景几何精度的方法，涵盖 LiDAR 融合、深度正则化、法向约束和结构基元选择等技术。

城市场景的弱纹理区域（建筑墙面、道路表面）和重复结构（窗户阵列）对 3DGS 的几何质量提出严峻挑战。标准 3DGS 仅靠光度损失驱动优化，在这些区域容易产生"墙壁凹陷""路面起伏"等几何失真。原因在于：缺乏纹理信息时，光度损失无法区分正确与错误的几何。大量工作通过引入外部几何先验（LiDAR 深度、单目深度估计、法向约束、语义引导）来补充光度损失的不足，显著提升几何精度。

\paragraph{LiDAR 平面约束 GMM 的连续几何监督}

% 改动概要：LI-GS 通过将 LiDAR 点云建模为平面约束多模态 GMM 并结合 2D Gaussian Surfels，为大规模场景提供从初始化到优化的连续几何监督，显著提升几何重建精度。
LI-GS（2024.09）~\cite{2409_12899} 提出一种 LiDAR-视觉融合的大规模场景重建系统，核心贡献在于将 LiDAR 点云转化为连续概率模型，为 Gaussian Splatting 的初始化、优化和网格提取全流程提供精确几何约束。

\textbf{基线做法与问题。}~大规模场景中标准 3DGS 面临三重几何瓶颈：（1）高斯椭球体具有厚度，违背薄表面假设，导致表面拟合差；（2）缺乏沿相机主轴方向的精确深度监督，光度损失过度影响高斯的几何属性（位置、形状），使其偏离正确表面；（3）自动驾驶等场景中监督视角稀疏，3DGS 易对单视角过拟合而丧失多视角几何一致性。现有 2DGS 虽将三维椭球压平为二维 surfel 改善表面贴合，但在缺乏精确深度信息的大规模场景中仍难以取得理想效果。而已有的 LiDAR-3DGS 融合方法因监督视角稀疏，LiDAR 难以对全场景提供充分约束。

\textbf{改了什么。}~LI-GS 做出三项关键设计：（1）以 2D Gaussian Surfel（而非 3D 椭球）作为地图表示，从根本上消除了厚度对表面贴合的干扰；（2）提出\emph{平面约束多模态 GMM（Gaussian Mixture Model）}建模方法，将离散 LiDAR 点云转化为连续概率分布，用于初始化和优化阶段的几何约束；（3）设计基于 GMM 的几何感知密度控制（Geometry-Aware Density Control）和由粗到细网格提取策略，精准管理高斯数量并去除伪影。

\textbf{实现原理。}~整个系统分为预处理、初始化、优化和网格提取四个阶段：

\emph{预处理阶段}——系统依次执行 LiDAR-Inertial SLAM（SLICT）估计初始位姿、动态物体去除（M-detector）、全局点云增强（HBA），以及图像位姿估计与联合优化（Colmap-PCD），获得精确的全局点云和一致的图像位姿。

\emph{初始化阶段}——这是 LI-GS 的核心创新。系统将彩色点云帧逐帧送入增量式 GMM 建模管线：首先对每帧进行体素化，在各体素内用 RANSAC（Random Sample Consensus）提取局部平面；将平面上的点集 $P = \{z_i = [p_i, g_i]^\top\}$ 变换到平面局部坐标系，以 4D 空间（3D 位置 + 灰度值）建模多模态 GMM。每个 GMM 分量的概率密度为
\begin{equation}
    p_L(\tilde{z}) = \sum_{l \in L} \pi_l \, \mathcal{N}(\tilde{z} \mid \tilde{\mu}_l, \tilde{\Sigma}_l),
\end{equation}

其中 $\pi_l$、$\tilde{\mu}_l$、$\tilde{\Sigma}_l$ 分别为第 $l$ 个分量的权重、均值和协方差。平面约束使 GMM 分量更薄、更贴合表面，同时 4D 建模（纳入灰度维度）能准确捕捉表面纹理分布。增量更新时，对新帧中的点计算其对全局 GMM 的对数似然，低于阈值 $\tau$ 的点被标记为新有效点并触发 GMM 更新。最终，全局 GMM 按特征值分解直接转换为初始 Gaussian Surfel 的位置、法向、半径和主轴方向。

\emph{优化阶段}——总损失函数包含五项：

\begin{equation}
    \mathcal{L} = \lambda_\text{G} \mathcal{L}_\text{G} + \mathcal{L}_p + \mathcal{L}_\text{sky} + \lambda_d \mathcal{L}_d + \lambda_n \mathcal{L}_n,
\end{equation}

其中 GMM 损失 $\mathcal{L}_\text{G} = \mathcal{L}_\text{pos} + \mathcal{L}_\text{shape} + \mathcal{L}_\text{normal}$ 利用最近邻 GMM 分量对高斯 surfel 的位置、形状控制点和法向进行联合监督。位置损失 $\mathcal{L}_\text{pos}$ 将高斯中心拉向局部表面：

\begin{equation}
    \mathcal{L}_\text{pos} = \frac{1}{G} \sum_{g=1}^{G} d_g(p_g), \quad d_g(p) = \sum_{k=1}^{K} w_{g_k} \| (p - \mu_{g_k})^\top \mathbf{n}_{g_k} \|_1,
\end{equation}

权重 $w_{g_k} = \exp(-\|p_g - \mu_{g_k}\|_2^2 / 2\sigma^2)$ 强调距离更近的 GMM 分量。形状损失 $\mathcal{L}_\text{shape}$ 对 surfel 的控制点施加同样约束，确保大面积 surfel 的形状也贴合表面。在密度控制方面，LI-GS 将加权距离 $d_g(p_g)$ 融入生长和剪枝准则：生长时对梯度阈值加权使得远离表面的高斯更容易被分裂；剪枝时降低远离表面高斯的等效透明度使其更易被移除。

\emph{网格提取阶段}——采用由粗到细策略：先构建体素占用图并移除未占体素中的采样点；再利用 GMM 加权距离 $d(p)$ 在已占体素中精细剔除远离表面的点；最后用 screened Poisson 方法提取网格。

\textbf{设计动机。}~选择 2D Gaussian Surfel 而非 3D 椭球，是因为城市场景中绝大多数表面（墙面、路面、屋顶）可近似为薄表面，surfel 的二维切空间自然贴合这种几何先验，避免椭球厚度引起的表面拟合误差。采用平面约束多模态 GMM 建模 LiDAR 点云，而非直接使用原始点云或 SfM 稀疏点，出于三重考量：（1）GMM 作为连续概率分布，能在 LiDAR 未覆盖的未观测区域提供充分监督，防止对单视角过拟合；（2）GMM 沿相机主轴方向提供额外约束，弥补光度损失在深度方向监督不足的根本缺陷；（3）平面约束去除噪声点，使 GMM 分量更薄，天然适合转换为 Gaussian Surfel。GMM 的多模态特性（而非单高斯）则允许同一区域存在多个纹理模式（如窗户与墙面的不同反射率），避免单峰假设导致的纹理混淆。

\textbf{关键参数与实验结果。}~LI-GS 在自采集的六个室内外场景（配备双 HESAI XT32 LiDAR、Insta360 相机、IMU（Inertial Measurement Unit）的采集平台）上进行评估，以 Chamfer-L1、F1-score（20\,cm 阈值）等指标对比 13 种方法。与纯 LiDAR 方法（SLAMesh、VDBFusion、SHINE-Mapping、NeRF-LOAM）和九种高斯方法（含 2DGS、Gaussian Surfels、PGSR、GOF 等，均以彩色点云初始化并加入天空掩码以保证公平）对比，LI-GS 在所有几何指标上均取得最优。相比 LiDAR-only 方法几何精度平均提升 $52.6\%$，相比高斯方法提升 $68.7\%$。消融实验验证了三方面贡献：GMM 初始化比 SfM 点和原始彩色点云更精确，尤其在地面反射等困难场景中避免了错误收敛；GMM 归一化显著降低噪声并产生更光滑表面；几何感知密度控制平均减少 $45.22\%$ 的高斯分量。渲染质量方面与当前高斯方法可比，训练时间略慢于 2DGS（因最近邻搜索开销），在单张 RTX-4090 上运行。

\paragraph{D-Normal 深度一致性与自适应置信度加权}

UrbanGS（2026.02）~\cite{2602_02089} 提出深度一致性 D-Normal 正则化模块，将 D-Normal 约束与外部深度监督结合，实现对位置、旋转等所有几何参数的综合更新。与此前仅依赖单目法向估计器（能有效更新旋转参数但难以更新位置参数）的方法不同，UrbanGS 通过将 D-Normal 约束与外部深度监督联合建模，使梯度信号同时作用于高斯的位置、旋转和形状参数，从而在大规模城市场景中实现更全面的几何修正。自适应置信度加权机制基于梯度一致性和逆深度偏差设计——当多视角梯度方向一致且逆深度偏差较小时赋予高置信度，反之则降低权重，有效抑制不可靠深度估计区域的错误约束。空间自适应高斯剪枝（SAGP）策略根据局部几何复杂度和可见性动态调整高斯密度：对平坦低复杂度区域执行更激进的剪枝以减少冗余，对几何复杂区域保留更多高斯以维持细节表达。统一的分割和视角分配方案消除边界伪影并优化计算负载。

\paragraph{点图密集初始化与渐进式混合几何优化}

MetroGS（2025.11）~\cite{2511_19172} 以分布式 2D Gaussian Splatting 为基础，提出结构化密集增强方案——利用 SfM 先验和点图（pointmap）模型实现更密集的初始化，结合稀疏补偿机制提高重建完整性。

\textbf{动机。}大规模场景中 SfM 稀疏点云覆盖不均匀，导致高斯初始化出现孔洞；同时不同区域的拍摄时间和光照条件可能截然不同，标准 3DGS 无法同时拟合多种外观。具体而言，点图模型从 SfM 稀疏点云推断密集的三维点对应关系，使高斯基元的初始分布更贴近真实表面；稀疏补偿机制则在 SfM 覆盖不足的区域额外注入新高斯，避免重建中出现孔洞。渐进式混合几何优化策略有机整合单目和多视角优化——先用单目深度和法向约束快速建立粗略几何，再切换至多视角一致性精化阶段，利用跨视角深度聚合和法向融合消除单视角估计的歧义，实现高效且精确的几何精化。针对大规模场景中常见的外观不一致问题，引入深度引导的外观建模方法，通过轻量 MLP 将深度特征映射为外观修正向量，学习具有 3D 一致性的空间特征，有效解耦几何与外观。

\paragraph{大视觉模型语义掩码与几何复杂度引导}

SA-GS（2024.05）~\cite{2405_16923} 利用大视觉模型（SAM、DINO）生成语义掩码，引入几何复杂度测量函数作为软正则化，引导特定语义区域内每个高斯的形状。具体而言，SAM 为每幅训练图像生成像素级语义分区，DINO 特征则提供跨视角一致的语义对应关系，二者融合后为每个高斯赋予语义标签。几何复杂度测量函数基于频谱分析推导：将局部区域的纹理频率与高斯的形状参数关联，高频纹理区域的高斯被约束为扁平以贴合表面细节，低频平坦区域则允许更大的形状自由度。同时提出基于语义区域的期望高斯数量估计方法，通过统计各语义类别在图像中的像素占比和几何复杂度，为不同区域的高斯密度提供自适应下界，防止弱纹理或小目标区域因致密化不足而丢失几何。最后通过基于概率密度的提取方法将高斯转换为点云——对不透明度场进行采样并过滤低置信度点，在几何精度指标上显著优于当时方法。

\paragraph{多视角最优视角选择与深度一致性致密化}

ULSR-GS（2024.12）~\cite{2412_01402} 专注于大规模航拍图像的表面提取任务，提出"点到照片"分割方法结合多视角最优视角匹配原则，为每个子区域选择最佳训练图像。该方法的核心创新在于多视角最优视角选择策略——对每个三维点，从所有可观测到它的相机中计算多视角深度一致性得分，选择深度误差最小的视角作为该点的最优源视角，确保每个子区域的训练图像组合最大化几何约束质量。训练阶段采用基于多视角几何一致性的自适应致密化策略：通过跨视角深度聚合将多个视角的深度预测融合为一致的三维点集，在深度一致性高的区域执行更密集的高斯分裂以增强表面提取细节，在一致性低的区域则保守致密化以避免伪影。在大规模航拍摄影测量基准数据集上显著提升复杂城市环境中的表面提取精度。

\paragraph{法向-梯度引导的结构基元选择与轻量网格提取}

SF-Recon（2025.11）~\cite{2511_13278} 提出一种直接从多视角图像重建轻量建筑表面的方法，跳过传统管线中"密集重建 → 网格化 → 后简化"的繁琐步骤。首先训练初始 3DGS 场获取视角一致表示，然后通过法向-梯度引导的高斯优化选择与屋顶和墙体边界对齐的基元——该优化策略计算每个高斯法向与图像梯度方向的一致性分数，将与建筑结构边界（屋脊、墙角、屋檐）高度对齐的高斯标记为结构基元并赋予更高保留权重。多视角边缘一致性剪枝在无需外部监督的情况下，通过检查每个高斯在多视角渲染中的边缘响应一致性来区分结构边缘高斯和非结构浮游体，跨视角边缘响应不一致的高斯被移除以增强结构锐度。最终通过多视角深度约束的 Delaunay 三角化将结构化高斯场转换为轻量建筑网格——先用深度约束的 3D-2D 可见性验证过滤不可靠连接，再以可见性驱动的 Delaunay 图割提取最终表面，生成的网格面片数和顶点数远少于传统 TSDF 方法。

\subsection{跨视角与语义引导}
\label{ss:14_2}

本节介绍两条技术路线：（1）跨视角融合与航拍-地面对齐；（2）语义辅助与先验引导。

\subsubsection{跨视角融合与航拍-地面对齐}

本路线介绍跨越无人机航拍与地面街景等大视角差异场景的融合重建方法，解决视角差异导致的特征匹配困难和高斯优化不稳定问题。

城市数字孪生需要同时利用无人机航拍和地面街景数据，才能获得完整的城市模型。航拍图像提供屋顶和街区布局信息，地面图像提供建筑立面和街道细节。然而，无人机通常从 50--200 米高空俯视，而地面相机在 1.5--2 米高度平视，两者之间的视角差异可达数十倍。这种巨大的视角差异导致特征匹配困难和高斯优化不稳定，是跨视角融合的核心难点。

\paragraph{航拍-地面双阶段训练与致密化冲突消解}

Horizon-GS（2024.12）~\cite{2412_01745} 系统性地解决航拍和街景图像的统一重建与渲染问题。在大规模城市场景建模中，现有方法要么仅处理航拍视角（如 VastGaussian、CityGaussian），要么仅处理街景视角（如 Hierarchical-3DGS），两者各自优化但无法协同，导致重建结果在视角切换时出现断裂和伪影。

Horizon-GS 通过深入分析发现，将两种视角直接联合训练会引发两大核心冲突：（1）致密化冲突——街景视角促使剪枝遮挡区域的高斯基元，而航拍视角则希望在同一区域重新生长基元，两者梯度方向相互干扰；（2）相机分布不均——街景相机数量远多于航拍相机且覆盖 360 度近景，导致优化过程偏向街景细节、航拍视角严重模糊。

为解决上述冲突，Horizon-GS 提出由粗到精的两阶段训练策略。第一阶段仅从航拍图像累积梯度进行致密化，建立覆盖整个场景的粗糙几何骨架，街景视角仅辅助初始高斯基元的全局对齐。第二阶段冻结第一阶段的 MLP 权重和基元属性以保护骨架结构，转而以街景视角为主精化细节。同时引入相机分布平衡机制——在 $[0,1]$ 区间随机采样，若采样值小于 $1/(R+1)$ 则选择航拍图像，否则选择街景图像（第一阶段 $R\!=\!2$ 优先航拍，第二阶段 $R\!=\!1$ 均衡训练）。

在致密化策略上，第二阶段不再使用 Scaffold-GS 的平均屏幕空间梯度，而是筛选同时满足高梯度、高不透明度、大投影半径的神经高斯基元作为新锚点，以捕获稀疏观测区域的细粒度细节。该方法还引入多分辨率 LOD 构建策略，分别计算航拍层级数 $K_a$ 和全局层级数 $K$，在第一阶段仅激活 $K_a$ 层、第二阶段开放全部 $K$ 层，以适应两种视角在细节层次上的显著差异。

Horizon-GS 的设计动机在于：航拍图像提供屋顶和街区布局的全局覆盖，街景图像提供建筑立面和街道的精细纹理，两者天然互补。如果能有效消除训练冲突，统一模型将在航拍和街景两个视角同时达到或超越各自独立训练的质量，实现真正的无缝跨视角体验。

在可扩展性方面，Horizon-GS 采用分区并行训练策略适配大规模城市场景：先按地面投影相机位置将场景划分为 $m\!\times\!n$ 个区块并扩展重叠区域，对航拍视角添加可见相机和辅助点云，对街景视角则从深度图生成点云以确保完整覆盖；各区块独立训练后合并，并转换为完全显式表示（移除 MLP、用球谐函数预测颜色）以加速渲染。此外，该工作构建了包含 5 个合成场景和 2 个真实场景的高质量航拍-地面配准数据集：合成数据基于自定义采集插件生成（涵盖不同尺度、季节和环境），真实数据由 DJI M350RTK 无人机（五相机配置）和自制六相机头盔分别采集航拍和街景图像，利用 ContextCapture 进行 SfM、Depth-Anything-V2 生成深度图、Grounded-SAM 移除移动物体。在包含 MatrixCity、UC-GS 数据集及自采数据集的 11 个场景上，Horizon-GS 在航拍 PSNR 上比 2D-GS 提升约 4 dB、街景 PSNR 提升约 1.5 dB，在大规模场景（如 Hillside\_Morning 和 Hillside\_Summer）上达到 51.5 FPS 实时渲染，几何重建质量也显著优于 2D-GS 基线。

\paragraph{扩散修复的跨海拔迭代桥接}

DRAGON（2024.07）~\cite{2407_01761} 面向大型建筑重建，处理只有近地面（手机拍摄）和高空（无人机拍摄）图像、缺少中间海拔视角的困难场景。

\textbf{动机。}当高海拔和低海拔图像之间存在大段海拔间隙时，直接联合训练 3DGS 会因缺少中间视角约束而产生严重伪影和几何退化。（Novel View Synthesis，新视角合成）算法本身外推中间海拔图像：先以高空和地面图像分别或联合训练 3DGS 模型，然后在中间海拔的虚拟相机位置渲染粗糙新视角图像，再利用感知正则化（基于 LPIPS 感知损失和结构相似性约束）修复这些低质量外推图像中的伪影和缺失区域。修复后的图像作为新增训练数据重新训练 3DGS，生成更高质量的中间海拔渲染，如此迭代直至航拍与地面视角之间的视觉特征鸿沟被充分桥接，从而实现跨海拔的相机配准和完整重建。DRAGON 在 Google Earth Studio 构建的 9 个大型建筑场景上验证了效果，包含多种建筑风格和尺度变化。

\paragraph{双分支初始化与梯度感知正则化的跨视角融合}

CrossView-GS（2025.01）~\cite{2501_01695} 针对跨视角数据的大视角变化导致 3DGS 优化困难的问题，提出多分支构建与融合方法。其双分支初始化策略分别从航拍和地面视角独立训练 3DGS 模型，将两个模型的高斯基元合并后作为跨视角重建的初始化，为高斯分布提供更完整的空间覆盖。在致密化阶段，已有双分支模型中的高斯分布提供了可靠的先验，指导跨视角训练中的新生成高斯向正确位置聚集。梯度感知正则化策略解决核心优化难题——当航拍视角的显著梯度被地面视角的大梯度变化过度平滑时，该方法对各视角的梯度分别计算并自适应加权，保留各视角中结构信息丰富的梯度信号，有效缓解大视角差异引起的细节丢失。高斯补充策略则将双分支中各自独有但另一分支缺失的高斯迁移到跨视角模型中，融合多分支互补信息。

\paragraph{几何感知因果注意力的渐进式扩散引导}

ProDiG（2026.04）~\cite{2604_02003} 解决仅从航拍图像生成地面视角和连贯 3D 场景模型的挑战。

\textbf{动机。}航拍图像虽然覆盖范围广，但从高空俯视与地面平视之间存在数十倍的视角差异，直接将航拍 3DGS 表示外推到地面视角会产生严重的几何失真和纹理模糊。具体而言，该方法在一系列递减的中间海拔上渲染带噪声的新视角图像，利用几何感知因果注意力模块修复这些图像——该模块将极线几何结构注入参考视角的扩散过程中，使扩散模型在生成地面细节时保持与航拍观测的几何一致性。修复后的图像被注入 3DGS 重新训练，然后渲染更低海拔的图像，如此逐步将高斯表示从高空精度推近地面级精度。距离自适应高斯模块通过轻量编码器预测调整因子，根据每个高斯的缩放特征和到活跃相机的距离动态调整高斯尺度和不透明度，确保跨越大范围视角间隙时重建的稳定性。这一渐进式框架无需额外地面真值视角即可实现几何一致的精化。

\paragraph{混合捕获训练策略}

\textbf{动机。}当无人机和地面相机同时拍摄同一场景时，两者的视角、尺度和运动模式截然不同，将它们混合训练会导致梯度方向冲突和优化不稳定。Cho（2026.05）~\cite{2605_00052} 从方差分解的角度系统分析了这一问题，发现双模态相机分布下，跨模态梯度方差远小于模态内方差，因此最简单的结构改变——每步渲染两个视角——就是最有效的策略，而方向感知投影、幅度校正、置信度门控等复杂方法均未超出随机双视角配对的种子方差范围。

\subsubsection{语义辅助与先验引导}

本路线介绍利用语义信息和结构化城市模型先验辅助 3DGS 重建的方法，通过 CityGML 建筑模型、大视觉模型和跨模态融合提升重建质量。

城市场景包含丰富的语义结构（建筑、道路、植被、车辆），这些语义信息可以显著提升重建质量和效率。例如，已知某区域是平坦道路，就不需要在那里分配大量高斯基元；已知某区域是建筑立面，就可以施加更强的法向约束。更重要的是，城市 GIS 数据库中通常已存在结构化的建筑模型（CityGML LoD2/3），这些现成的几何先验远比 SfM 稀疏点云更为丰富和可靠。

\paragraph{CityGML 结构化先验的两遍光线投射语义重建}

现有语义 3DGS 方法主要依赖 2D 基础模型（SAM、CLIP 等）进行特征蒸馏或掩码提升，在通用物体发现上有效，但面对城市立面的重复纹理、反射材质和严重遮挡时，语义边界常常模糊不清，且难以保持建筑层级的部分-整体关系（窗户→墙面→整栋建筑）。

GS4City（2026.04）~\cite{2604_11401} 将 CityGML LoD3 标准的结构化城市模型作为几何与语义双重先验引入语义高斯重建，核心改变包括三方面：（1）从 LoD3 城市模型出发，通过两遍光线投射（two-pass raycasting）生成可靠的图像对齐掩码，并利用 CityGML 父子关系恢复单遍投影容易丢失的精细立面元素（如窗户、门洞）；（2）将城市模型掩码与基础模型（SAM + CLIP）生成的非建筑掩码进行优先级融合，通过跨视角几何重叠与视觉-语言语义验证建立场景一致的实例对应关系；（3）为每个高斯学习一个紧凑的 $D$ 维身份编码（identity code），通过 $\alpha$-混合 渲染后由线性分类器预测实例身份，辅以 3D 空间正则化约束近邻语义一致性。

实现原理是一条完整的"城市模型→掩码→融合→身份编码"管线。首先将 CityGML 多边形三角化并对齐到重建坐标系，为每个三角面赋予三级层级标签（建筑级 id\_feat、表面级 id\_surf、部件级 id\_part）。第一遍光线投射获取前表面层级标签；第二遍仅对部件级几何投射，通过深度容差 $\epsilon=0.5$\,m 和父表面 ID 一致性验证恢复被墙面遮挡的窗户/门洞。城市模型掩码覆盖建筑区域但不涵盖全场景，因此用 SAM 生成原始掩码后，通过质量过滤和建筑区域重叠比检测，再用 CLIP 视觉-语言模型以边界余量区分前景与建筑。跨视角关联以预训练高斯场作为几何桥梁，结合几何重叠率与语义特征相似度将逐视角掩码聚合为场景一致的全局实例组。融合时采用图像掩码优先规则（非建筑前景优先），并通过固定 ID 偏移量保持城市模型实例与图像实例不冲突。每个融合实例还关联一个跨视角平均的 CLIP 特征，支持开放词汇查询。身份编码阶段冻结高斯的几何与外观参数，仅优化每个高斯的 $D$ 维身份向量，通过 $\alpha$-混合 渲染后由线性分类器映射到实例空间，训练目标为逐像素交叉熵损失加上 3D 空间近邻 KL 散度正则项。

设计动机在于：城市模型（CityGML LoD3）与 SfM 稀疏点云或 2D 基础模型有本质区别——它提供的是显式的、带层级语义的、地理参考的边界表示（B-Rep），明确编码了窗户、门洞等立面元素及其与墙面、屋顶的父子关系，这是任何纯图像方法都无法自动恢复的结构化知识。直接将城市模型的 3D 标签投射到高斯基元上并不可靠，因为城市模型定义在显式表面上，而 3DGS 使用 $\alpha$-混合 的基元、缺乏严格表面归属，浮点和遮挡进一步加剧了 3D 分配的不稳定性。因此 GS4City 选择在 2D 图像空间构建可靠监督，再提升到高斯场。

在 TUM2TWIN 和 Gold Coast LoD3 两个城市数据集上，GS4City 在粗粒度建筑/非建筑分割上达到 93.1\% 和 95.4\% IoU（分别超过 LangSplat 和 Gaga 最多 15.8 个百分点），在细粒度语义分割（墙面、屋顶、窗户、门、树木、行人、车辆）上将 mIoU 从 22.3\%/21.0\% 提升至 30.5\%（TUM2TWIN）和从 23.5\%/30.2\% 提升至 44.4\%（Gold Coast），提升幅度最高达 14.2 个百分点。增益最大的类别是具有稳定结构先验的屋顶、墙面、窗户和车辆，而树木等不规则类别上 LangSplat 的纯特征蒸馏仍有优势。GS4City 在时间归一化 mIoU（mIoU/min）上也取得最优的精度-效率平衡。

\paragraph{建筑模型深度-法向先验引导的表面初始化}

GS4Buildings（2025.08）~\cite{2508_07355} 利用低级语义 3D 建筑模型（LoD2）替代传统 SfM 初始化高斯。具体而言，该方法从 CityGML LoD2 模型的平面建筑几何出发，通过光线投射为每个训练视角生成先验深度图和法向图——深度先验损失将渲染深度拉向建筑模型投影深度，法向先验损失约束高斯法向与建筑表面法向对齐，二者在优化过程中提供强大的几何引导，有效解决弱纹理墙面和遮挡区域的表面定位问题。以 2D Gaussian Splatting 为基础基元，高斯基元直接从建筑模型表面采样初始化而非依赖 SfM 稀疏点云，使初始分布天然贴合建筑结构。可选的建筑聚焦模式利用建筑模型掩码将重建限制在建筑区域，排除植被、道路等非建筑元素，减少 71.8\% 的高斯基元，实现重建完整性提升 20.5\%、几何精度提升 32.8\%。

\paragraph{Language-SAM 建筑分割与边界法向约束}

PG-SAG（2025.01）~\cite{2501_01677} 利用跨模态模型 Language Segment Anything 分割建筑掩码，根据可见性检查将建筑区域分组为子区域并行优化。语义感知数据分组的具体流程为：Language Segment Anything 以文本提示"building"分割出像素级建筑掩码，对每张图像统计其建筑像素占比和与各子区域的空间重叠度，通过可见性检查（被建筑掩码覆盖且与子区域有空间交集）将图像分配到对应子区域。边界感知法向损失针对掩码边缘重新构造——标准法向损失在语义边界处产生歧义（建筑与天空的交界处法向方向不明确），PG-SAG 对检测到的掩码边缘像素降低法向约束权重或替换为深度一致性约束，有效缓解边缘法向歧义。梯度约束均衡负载损失则通过分析各子区域的渲染计算复杂度，动态调整像素并行渲染阶段的工作负载分配，最小化线程等待时间和重建损失。

\paragraph{图像-LiDAR-语义三模态紧凑嵌入融合}

Neural-MMGS（2025.09）~\cite{2509_17762} 将图像、LiDAR 和语义三种模态融合为每个高斯的紧凑可学习嵌入，通过轻量级神经解码器映射回各模态的高斯参数。与传统方法将各模态属性作为独立参数附加到高斯不同，Neural-MMGS 的融合嵌入隐式编码了光学、物理和语义特征，以更低的内存开销实现更好的可扩展性。在 Oxford Spires 和 KITTI-360 数据集上验证了该方法的有效性。

\subsection{外观优化}
\label{ss:14_3}

本节介绍外观一致性与正则化技术，解决大规模城市场景中跨视角光照不一致的重建难题。

\subsubsection{外观一致性与正则化}

本路线介绍处理大规模城市场景中跨视角外观不一致的方法，包括外观变换模块、前景-背景分离和扩散引导约束等技术。

大规模城市场景中，不同视角的拍摄时间、光照条件、天气状况可能不同，导致严重的跨视角外观不一致。举例来说，一条街道的两侧照片：一侧在清晨阳光下拍摄，另一侧在午后阴天拍摄，光照条件截然不同。标准 3DGS 假设场景外观固定不变，试图用一组球谐系数同时匹配两种截然不同的外观，结果往往是两边都重建不好。

\paragraph{两阶段室外场景重建}

Pintani et al.（2025.10）~\cite{2510_09489} 提出两阶段框架，显式分离和优化前景与背景。

\textbf{动机。}室外场景中的动态物体（行人、车辆）和天空背景在 3DGS 中往往互相干扰——动态物体留下的残影污染背景高斯，天空区域的浮游体遮挡前景细节。第一阶段将背景基元初始化在球壳内，用背景专属光度项和几何正则化器优化；第二阶段从 SfM 重建添加前景高斯，用标准渲染损失精化，同时固定背景集合。这种显式背景分离不仅减少伪影，还实现了自动、无物体的环境贴图估计。

除显式前景-背景分离外，另一种思路是通过可学习的外观变换模块直接建模跨视角外观差异。

\paragraph{外观变换模块}

Yuan et al.（2025.07）~\cite{2507_23006} 提出鲁棒高效的城市场景重建框架，通过外观变换模块缓解多视角拍摄间的外观不一致性，同时支持灵活调整。

\textbf{动机。}城市场景通常在数小时甚至数天内采集完成，不同时段的光照和天气变化导致同一位置在不同图像中呈现截然不同的外观，标准 3DGS 的固定球谐表示无法建模这种跨视角外观差异。该方法结合场景分割并行训练、可控 LoD 策略和多种增强模块（深度正则化、尺度正则化、抗锯齿），实现城市尺度场景的快速重建和实时渲染。

前述方法主要处理已知视角间的外观不一致。当训练视角本身稀疏时，生成式模型可提供额外约束。

\paragraph{扩散引导的约束重建}

GS-Diff（2025.04）~\cite{2504_01960} 引入多视角扩散模型生成伪观测，将欠约束的 3D 重建问题转化为良定问题，在稀疏数据下实现鲁棒优化。

\textbf{动机。}当训练视角稀疏或分布不均匀时，3DGS 在未被充分观测的区域只能依靠有限的光度约束进行优化，几何和外观都容易出现严重退化。核心机制是利用预训练的多视角扩散模型，以已有输入图像为条件合成补充视角的伪图像，这些伪观测为稀疏覆盖区域提供额外的光度监督，使原本欠定的三维结构推断变得充分约束。为应对野外无约束场景的多重挑战，框架集成了多项增强技术：外观嵌入为每张图像学习低维隐编码以捕获跨视角光照差异；单目深度先验通过 Depth-Anything 估计的深度图施加软几何约束；动态物体建模利用瞬态掩码检测并忽略行人、车辆等运动物体；各向异性正则化惩罚过度拉伸的高斯基元以维持几何合理性；Mip 滤波和绝对梯度等高级光栅化技术减少锯齿和模糊伪影。在 ULTRRA 等四个基准上持续超越基线方法，在稀疏和异质数据条件下尤为显著。

\paragraph{其他城市重建方法}

多个工作从不同角度推进了城市场景重建。

\textbf{效率与压缩。}EfficientGS（2024.04）~\cite{2404_12777} 通过选择性致密化和冗余剪枝，在保持渲染保真度的同时将模型大小缩小约 10 倍，使 4K+ 航拍图像的大规模场景重建更加高效。具体而言，该方法在标准 3DGS 致密化过程中引入基于渲染贡献的冗余检测，剪除对图像质量贡献低于阈值的高斯基元，并配合球谐系数降阶进一步压缩存储。HO-Gaussian（2024.03）~\cite{2403_20032} 提出混合优化方法，将基于网格的体表示与 3DGS 管线结合，消除对 SfM 点初始化的依赖，用高斯方向编码替代球谐函数实现视角依赖颜色表示。CityGS-X（2025.03）~\cite{2503_23044} 放弃传统的合并-分割过程，采用批级多任务渲染实现高效多 GPU 渲染，4 张 4090 即可在 5 小时内训练渲染 5000+ 图像的场景。

\textbf{配准与融合。}PhotoReg（2024.10）~\cite{2410_05044} 提出利用 3DGS 的渲染能力与 3D 基础模型的对偶性，实现多个 3DGS 子地图的自动配准和融合，解决了多机器人协同建图的关键问题。核心做法是从各子地图渲染重叠视角的图像，利用基础特征匹配建立跨地图的 2D-2D 对应关系，再通过鲁棒估计求解子地图间的相对变换。GaussReg（2024.07）~\cite{2407_05254} 提出粗到细的 3DGS 配准框架——粗阶段通过点云配准估计粗略对齐，细阶段利用从 3DGS 渲染的图像提供精细几何信息实现精确对齐，速度比 HLoc 快 44 倍且精度相当。

\textbf{动机。}多个独立训练的 3DGS 子场景需要融合为统一模型，但各子场景的坐标系和尺度可能不一致，直接拼接会产生几何错位。Liu et al.（2025.07）~\cite{2507_20480} 提出自动 3DGS 子地图配准融合方法，通过骨架对齐和椭球感知卷积消除人工干预。GaRField++（2024.09）~\cite{2409_12774} 通过可见性相机选择和渐进点云扩展解决可扩展性问题，利用 ConvKAN 网络解耦外观模块处理光照不均匀，并在无缝拼接阶段合并各子区域高斯。

\textbf{特殊场景与数据。}LetsGo（2024.04）~\cite{2404_09748} 开发了配备 IMU、LiDAR 和鱼眼相机的手持扫描仪 Polar，构建了 GarageWorld 大型车库数据集，提出深度正则化和多分辨率 3D 高斯表示实现高效 LOD 渲染。Mode-GS（2024.10）~\cite{2410_04646} 将单目深度锚定引入地面机器人轨迹数据集，通过残差形式高斯解码器和尺度一致深度损失解决地面场景中高斯漂移问题。DroneSplat（2025.03）~\cite{2503_16964} 针对野外无人机图像中的动态干扰物和有限视角约束，通过局部-全局分割启发式和统计方法自适应调整掩码阈值，结合多视角立体预测和体素引导优化。SplatCo（2025.05）~\cite{2505_17951} 通过跨结构协作模块融合全局三平面和局部上下文网格特征，配合跨视角剪枝和结构-视角共学习机制实现细节保持的大规模场景渲染。CityGo（2025.05）~\cite{2505_21041} 采用混合表示——纹理化代理网格加残差高斯——在移动 GPU 上实现实时渲染，训练速度平均提升 1.4 倍。

Chen et al.（2025.06）~\cite{2506_17636} 引入粗到细策略、解耦外观模型和瞬态掩码模型实现大规模场景的精细表面重建。SMW-GS（2025.06）~\cite{2506_13516} 通过微宏观投影和小波采样增强多尺度细节多样性，配合相机-分割最优分配策略实现可扩展重建。ExploreGS（2025.05）~\cite{2505_10578} 提出基于视觉的低开销无人机 3DGS 框架，利用词袋模型实现板载实时处理。Gao et al.（2024.05）~\cite{2405_11021} 利用 Google Earth 图像构建 3DGS 模型，在 Benchmark 上超越基于 NeRF 的方法。Chen et al.（2024.10）~\cite{2410_11285} 使用无人机搭载 360 度相机采集室内场景，通过分治策略和粗到细对齐实现可扩展的大规模室内重建。EnerGS（2026.04）~\cite{2604_26238} 将部分可观测几何建模为由几何证据诱导的连续能量场，为高斯基元提供软几何引导而非硬约束。Kulkarni et al.（2026.03）~\cite{2603_26638} 开发了端到端管线处理汽车经销商中移动车辆的三维重建，利用 SAM3 实例分割和畸变感知 3DGS 框架。

Guo et al.（2025.11）~\cite{2511_06765} 利用 LiDAR-IMU 里程计提供先验位姿，配合法向约束和有效秩正则化增强弱纹理室外场景的重建。GraphGS（2025.02）~\cite{2502_17377} 设计了基于空间先验的场景结构估计和相机图，通过图引导的多视角一致性约束和自适应采样策略加速 3D 重建并缓解过拟合。St. John 大教堂项目（2026.04）~\cite{2604_24316} 展示了文化遗产保护的大规模应用——7 个夜晚采集 99,000 张图像，结合 DSLR、无人机和 LiDAR 扫描创建包含 250--300 亿三角形的 3D 重建。

\begin{takeaway}
城市场景重建的核心挑战在于效率、配准与特殊场景处理：选择性致密化与剪枝实现约 10 倍压缩，基于渲染的子地图配准解决多机器人协同建图，混合表示（网格+高斯残差）则在移动端实现实时渲染。
\end{takeaway}

% ======================================================================

\subsection{近期方法脉络：数字孪生建模、数字孪生与通信建模、Registration and merging of Gaussian scenes等}

这一组近期工作延续 3DGS/SuGaR、Gaussian splatting、3DGS，主要把改动集中在应用、适配、注册/合并、优化等环节。下面只保留相对既有路线的变化，便于和前文方法直接对照。

\paragraph{RFDT-Channel}
RFDT-Channel（2026.05）~\cite{2606_01261} 延续 3DGS/SuGaR 在“数字孪生建模”上的基本设定，主要改变是：RFDT-Channel：RGB+LiDAR采集，COLMAP/3DGS/SuGaR生成初始mesh，再用LiDAR和Blender做RF规则化与材料绑定。3DGS 的显式高斯表示和可微渲染能力保持为基础。
其设计动机是：提出面向毫米波射线追踪的场景构建工作流，把视觉重建mesh转化为RF可计算几何并绑定电磁材料属性。
具体而言，系统使用RGB视频和LiDAR点云，通过COLMAP、3DGS和SuGaR得到初始三角网格，再在Blender中进行尺度对齐、墙体实体化和门窗开口等RF规则化操作。
适用边界是：3DGS只是流程组件之一；大量RF规则化依赖人工或工具链后处理；与3DGS方法创新关系较弱。

\paragraph{Point-Cloud-Assistant Localized Statistical Channel Prediction by Tangent Gaussian Splatting}
Point-Cloud-Assistant Localized Statistical Channel Prediction by Tangent Gaussian Splatting（2026.06）~\cite{2606_18734} 延续 Gaussian splatting 在“数字孪生与通信建模”上的基本设定，主要改变是：Tangent Gaussian Splatting辅助的localized statistical channel prediction。3DGS 的显式高斯表示和可微渲染能力保持为基础。
其设计动机是：利用点云和高斯局部表示补全未测位置的信道统计信息，服务大规模网络优化。
具体而言，方法围绕参考信号接收功率和点云几何构建局部高斯表示，预测多径角功率谱。
适用边界是：任务不是视觉3DGS重建；高斯splatting更多作为几何/统计建模工具。

\paragraph{Graph-GSReg}
Graph-GSReg（2026.06）~\cite{2606_29782} 可作为“Registration and merging of Gaussian scenes”中的新锚点候选。相较于 3DGS，主要改变是：构建 3D scene graph，将 3DGS registration 转化为图匹配和优化问题。3DGS 的显式高斯表示和可微渲染能力保持为基础。
其设计动机是：避免直接在 Gaussian primitives 间学习对应或依赖昂贵图像优化，把注册提升到对象/关系图层面。
具体而言，从 3DGS 及其渲染图像构建 3D scene graph，再基于图结构寻找跨场景对应和变换，最终合并成统一 Gaussian 表示。
这一锚点判断来自当前分类中的机制关系：候选锚点：注册/合并，且相关性 0.90；后续方法可围绕 面向多个 Gaussian scene 的跨场景配准与统一表示构建。 继续比较和扩展。
适用边界是：依赖场景图质量和可匹配对象关系；低结构或重复场景可能出现错配。

\paragraph{Structure-Aware Gaussian Splatting for Large-Scale Scene Reconstruction}
Structure-Aware Gaussian Splatting for Large-Scale Scene Reconstruction（2026.07）~\cite{2607_01698} 延续 3DGS 在“Structure-aware scheduling for large-scale Gaussian reconstruction”上的基本设定，主要改变是：结构感知调度与频率调制，缓解稀疏初始点导致的失控 densification。3DGS 的显式高斯表示和可微渲染能力保持为基础。
其设计动机是：把 densification target signal frequency 和 image sampling frequency 两类调度统一考虑，以控制高频监督和稀疏点初始化之间的矛盾。
具体而言，在大规模场景中识别 sparsely observed regions，通过结构感知的 densification 和分辨率/采样调度避免用高频图像监督低频稀疏点时产生 uncontrolled densification。
适用边界是：调度策略可能依赖场景结构估计和超参数；极端稀疏区域仍可能需要外部先验。
\section{在线与增量重建}
\label{sec:14_2}

前述方法主要关注离线场景——所有图像在训练开始前就已可用。然而，许多实际应用需要在线（on-the-fly）重建：机器人探索未知环境时需要即时构建地图，自动驾驶车辆需要持续更新环境模型，增强现实应用需要在拍摄过程中实时渲染。在线重建的核心挑战在于：数据按时间顺序到达，无法回溯过去帧；计算必须在有限时间和内存内完成；场景可能随时间变化，需要持续更新。

\subsection{流式与增量更新}
\label{ss:14_4}

本节介绍两条技术路线：（1）流式重建；（2）场景变化检测与更新。

\subsubsection{流式重建}

本路线介绍在图像序列逐步到达过程中实时构建3DGS模型的流式重建方法。

流式重建（streaming reconstruction）的目标是在图像序列逐步到达的过程中实时构建 3DGS 模型。与离线重建的最大区别在于：系统无法回溯已经处理过的图像（除非显式存储为关键帧），因此每张新图像必须在"看到即优化"的模式下增量更新模型。核心挑战是在计算效率、内存管理和重建质量之间取得平衡。具体而言，处理速度必须跟上图像到达速率，内存使用不能无限增长，同时要保证重建质量持续提升。

\paragraph{渐进式局部-半全局在线优化}

Gaussian On-the-Fly Splatting（2025.03）~\cite{2503_13086} 提出渐进式框架，在图像采集过程中实现近实时 3DGS 优化。每张新图像到达时，通过 On-the-Fly SfM 更新位姿和稀疏点，新优化的高斯立即整合到 3DGS 场中。为平衡新旧图像的训练，提出渐进式局部与半全局优化——优先处理新图像及其邻域的重叠图像。自适应学习率调度平衡迭代次数和学习率，稳定跨图像训练。实验表明，每张新图像可在数秒内完成优化。

整个管线分为三个阶段运作。初始阶段以 On-the-Fly SfM 估计前 $N_0$ 张图像的位姿并生成稀疏点云，训练初始 3DGS 场，采用严格迭代限制防止过拟合。渐进训练阶段中，新图像到达后 On-the-Fly SfM 持续更新全部位姿并细化稀疏点云，新稀疏点通过搜索树判别为新增加点后初始化为新高斯融入 3DGS 场；局部+半全局优化基于重叠关系为新图像和相邻图像分配更多迭代，同时随机选择其他图像防止整体退化。所有图像采集完成后，SfM 执行全局光束平差，再经快速额外优化提升整体渲染质量。自适应学习率根据每张图像已获得的训练迭代数动态计算，使新图像获得更高学习率以快速收敛。

\paragraph{增量聚类卸载的大规模在线重建}

Meuleman et al.（2025.06）~\cite{2506_05558} 提出的 on-the-fly 方法在拍摄完成后即可产出相机位姿和训练好的 3DGS，支持密集和宽基线拍摄以及大规模场景。方法引入快速初始位姿估计，利用学习特征和 GPU 友好的迷你光束平差。通过直接采样高斯基元位置和形状，增量生成所需基元，显著加速训练。增量方法通过渐进聚类高斯基元、存储在锚点并卸载出 GPU 来处理大规模场景。这一可扩展辐射场构建方法使 3DGS 在任何视点保持可控规模。

该方法的位姿估计采用 XFeat 快速特征提取（每帧 6144 个关键点），以固定大小的迷你光束平差利用 GPU 并行化高效求解，避免了 Ceres 等通用求解器的开销。在高斯基元生成方面，跳过传统 SfM 稀疏点初始化步骤，直接从三角化关键点创建高斯基元并增量式添加，使训练启动显著加速。可扩展性通过三级策略实现：渐进式聚类将空间邻近高斯归并为锚点，聚类后的基元按需从 GPU 卸载到 CPU 内存，并支持聚类基元的渐进合并以保持任意视点处 3DGS 规模可控，从而在保持竞争力的图像质量的同时覆盖从小尺度到城市尺度的多种场景。

\paragraph{DUSt3R 初始化与 MergeNet 合并的无位姿流式重建}

StreamGS（2025.03）~\cite{2503_06235} 是首个面向无位姿图像流的在线前馈 3DGS 重建方法，能够在不依赖已知相机位姿和内参的条件下，从连续图像流中渐进式构建 3D 高斯表示。

\textbf{基线做法。}在线 3DGS 重建面临三重挑战：其一，传统 SfM 预处理需要完整图像集且计算耗时，无法满足实时流式场景；其二，现有可泛化方法（如 pixelSplat、MVSplat）依赖已知位姿进行立体匹配，且仅支持稀疏视角输入（至多 2--3 帧）；其三，无位姿方法（如 CF-3DGS）将位姿作为可学习参数与高斯联合优化，每帧需要数千次迭代，重建速度极慢。

\textbf{改了什么。}StreamGS 提出完整的无位姿前馈重建管线，包含三个核心设计：（1）基于 DUSt3R 的初始双视角重建模块，同时预测像素级 3D 点图（pointmap）和相机矩阵（内参+外参），无需任何位姿先验；（2）内容自适应精化（adaptive refinement）模块，利用 MASt3R 匹配头在相邻帧间建立稠密像素对应关系，联合精化高斯中心和相机位姿；（3）前馈自适应密度控制（feed-forward ADC）模块，通过跨帧特征聚合将相邻帧的匹配高斯对合并为单一高斯，从而减少约 40\% 的高斯数量，同时保持渲染质量。

\textbf{实现原理。}在初始重建阶段，对相邻帧 $(I^{t-1}, I^t)$ 输入 DUSt3R，获得参考坐标系下的像素级 3D 点图及置信度图，再通过点配准问题求解相对位姿（含缩放因子）。在精化阶段，利用 MASt3R 提取局部 3D 特征，通过最近邻互惠搜索（nearest reciprocal）建立帧间稠密匹配，以此重新估计残差变换并更新点图和位姿（PnP-RANSAC）。在高斯生成阶段，用独立的 2D 图像编码器提取外观特征（因 DUSt3R 冻结），将点图坐标与图像特征拼接后送入轻量解码器预测旋转、缩放、不透明度和球谐系数。在合并阶段，利用已建立的 2D 像素对应关系将 3D 空间的高斯合并转化为高效的 2D 像素级特征变形（warping），再由双层卷积网络（MergeNet）将匹配高斯对聚合为单一高斯——整个合并过程无需渲染损失的梯度反传，运行极快。增量式地将每帧新高斯 $G^t$ 与历史重建 $G^{t-1}$ 合并，最终累积为完整场景表示。

\textbf{设计动机。}之所以能做到无位姿前馈重建，关键洞察在于：DUSt3R 作为预训练的像素级 3D 点预测器，已从大规模 3D 数据中习得了强几何先验，可直接从图像对预测点图而无需已知位姿。之所以比优化方法快 150 倍，是因为整个流程是前馈式（feed-forward）的——所有操作（位姿估计、高斯生成、合并）都在单次前向传播中完成，无需迭代优化。前馈 ADC 的核心观察是：相邻帧的匹配高斯对在形状和颜色上高度相似且空间分布接近，是冗余且可剪枝的；而已有 2D 像素对应关系使得 3D 合并可转化为高效 2D 操作，避免了 $O(n^3)$ 的 3D 网格遍历。精化模块的必要性在于 DUSt3R 冻结后对域外数据（如室内场景）预测不准，精化可将 PSNR 提升 23\%--39\%。

\textbf{关键参数。}在 RE10K 数据集上，StreamGS 的平均 PSNR 为 23.1 dB，仅比优化方法 CF-3DGS 低 0.05 dB，但处理速度达每秒约 9 帧，比 CF-3DGS 快 150 倍。在域外数据集（DL3DV、MVImgNet、ScanNet）上，StreamGS 显著优于依赖已知位姿的 pixelSplat 和 MVSplat，展示了更好的跨域泛化能力。前馈 ADC 在 MVImgNet 和 ACID 上分别剪枝 36.71\% 和 40.48\% 的高斯，PSNR 仅下降 2.53\% 和 3.90\%。方法在 RE10K 上训练 30K 步（A100 GPU，$224\times224$ 分辨率，批大小 14）。

\paragraph{高斯-图像表示与身份感知冗余压缩}

LongSplat（2025.07）~\cite{2507_16144} 专为长序列图像设计，核心是流式更新机制——增量整合当前视角观测同时选择性压缩冗余历史高斯。

\textbf{动机。}标准在线 3DGS 方法逐帧添加高斯但不做冗余管理，高斯数量随序列长度线性增长，长序列场景中很快耗尽内存并导致推理速度急剧下降。关键创新是高斯-图像表示（GIR），将 3D 高斯参数编码为结构化的类图像 2D 格式。GIR 同时实现当前-历史高斯的高效融合和身份感知冗余压缩，使在线重建适应长序列而不会压垮内存或计算资源，高斯数量相比逐像素预测方法减少 44\%。

具体而言，GIR 通过可微透视投影将历史高斯渲染到当前视角的对齐 2D 网格中，使历史信息与当前帧在同一空间内对齐。每帧的流式更新执行两个关键操作：自适应压缩（选择性压缩历史高斯消除冗余）和在线集成（融合当前视角高斯与历史状态）。特征提取管线结合多视角空间特征（来自当前帧和 $N$ 个时序相邻帧）与历史特征（GIR 编码后经浅层 CNN 提取），通过 Transformer 注意力机制融合空间和时间信息，预测更新掩码和压缩权重。该设计借鉴图像压缩方法的紧凑表示思想指导生成更高质量的 3D 高斯，在实时新视角合成任务上达到当时最优的效率-质量平衡。

\paragraph{多相机实时重建}

Guo et al.（2025.12）~\cite{2512_08498} 提出首个面向多相机设备组的在线 3D 重建框架，增量融合多个重叠相机的密集 RGB 流到统一高斯表示中。层次化相机初始化方案在无标定情况下实现粗级相机间对齐，轻量级多相机光束平差在保持实时性能的同时稳定轨迹。冗余消除高斯采样策略和频率感知优化调度减少高斯基元数量和所需优化迭代。该方法仅使用原始多相机视频流即可在 2 分钟内重建数百米的三维场景。

\paragraph{松耦合三角几何-神经高斯外观框架}

PLANING（2026.01）~\cite{2601_22046} 提出松耦合混合表示——显式几何基元（三角面片）与神经高斯分离建模几何和外观。

\textbf{动机。}纯高斯表示在在线流式重建中面临两难：高斯数量太少无法表达细节，太多则内存和计算开销过大；且高斯之间的冗余和竞争在增量更新中难以自动解决。这种解耦支持在线初始化和优化策略，将几何和外观更新分离，产生结构冗余显著降低的稳定流式重建。PLANING 在 ScanNetV2 上 100 秒内完成重建，比 2DGS 快 5 倍以上，同时质量媲美离线逐场景优化。

在技术实现上，可学习三角形基元基于三个可学习顶点定义，构建局部坐标系并引入边缘锐度和平滑度参数，实现边缘保持的贡献函数和可微三角形光栅化器。三角形作为稳定结构锚点，特别适合室内环境中的平面布局。神经高斯外观建模受 Scaffold-GS 启发，高斯属性从三角形特征与 per-Gaussian 特征融合的表示中解码，鼓励局部渲染一致性和平滑度；渲染梯度通过神经高斯反向传播精炼几何。在线重建流水线将几何和外观更新解耦：利用三角形法线和深度先验监督几何，外观通过高斯渲染损失优化，支持渐进式场景扩展和新视角融合。相比 PGSR 密集网格 Chamfer-L2 提升 18.52\%，超越 ARTDECO 1.31 dB PSNR。

\paragraph{在线 RGB-D 密集建图}

OG-Mapping（2024.08）~\cite{2408_17223} 利用稀疏八叉树的结构化场景表示能力，结合结构化 3D 高斯表示实现高效鲁棒的在线密集建图。锚点渐进式精化策略在多个细节层次恢复场景结构。动态关键帧窗口替代固定窗口，更好地处理局部极小值和遗忘问题。HGS-Mapping（2024.03）~\cite{2403_20159} 是首个将高斯表示集成到城市场景在线密集建图的工作，引入混合高斯表示建模场景不同部分，配合混合初始化和自适应更新方法实现高保真快速重建。

\paragraph{活跃-稳定状态演化与语言场联合建模}

OnlineX（2026.03）~\cite{2603_02134} 是前馈框架，使用流式图像同时在线重建 3D 视觉外观和语言场。关键创新是解耦的活跃-稳定状态演化范式——将记忆状态解耦为活跃状态（持续刷新捕获高频局部几何）和稳定状态（保守累积长期全局结构），然后融合前者信息到后者实现保真度与稳定性的兼顾。

在实现层面，活动状态通过当前帧与前一帧的成对交互提取相对的像素级几何和外观信息，稳定状态作为持久全局锚点保存长期全局结构。活动特征与稳定状态融合后计算全局一致的位姿特征，用于预测当前帧的 3D 高斯。此外，OnlineX 联合建模视觉外观和语言场——为每个高斯基元增加低维语言特征（16 维 vs 原始 CLIP 512 维），通过 $\alpha$-混合 同时渲染新视角 RGB 图像和语言图，实现 3D 语义理解。隐式高斯融合模块合并跨多视角的重叠 3D 高斯以减少冗余并保持一致性。整个框架无需离线 SfM，直接从无位姿 RGB 流进行连续重建，在主流数据集上持续超越先前工作的新视角合成和语义理解表现，并具有实时推理速度。

\paragraph{因果 Transformer 与查询驱动解码的流式语义}

S2GS（2026.03）~\cite{2603_14232} 提出严格因果、增量的 3D 高斯语义场框架——不利用未来帧，持续更新场景几何、外观和实例级语义而无需重新处理历史帧。几何-语义解耦双骨干设计中，几何分支执行因果建模驱动增量高斯更新，语义分支利用 2D 视觉基础模型和查询驱动解码器预测分割掩码和身份嵌入。S2GS 处理 1000+ 帧时运行时和内存增长远慢于离线全局处理基线（后者约 80 帧即 OOM）。

S2GS 采用因果 Transformer 编码器处理视频流，通过时间注意力掩码限制每帧只能关注历史前缀，实现并行训练等价于自回归因果模型。推理时缓存历史帧的 key/value 张量以支持高效长序列流式推理。几何流从预训练 3D 基础模型蒸馏几何监督信号（伪深度图和相机参数），通过反投影得到像素对齐的 3D 高斯中心；语义流使用冻结的 2D 视觉基础模型（SigLIP2）提取语义特征，通过查询驱动解码器预测分割掩码和身份嵌入，结合查询级对比对齐和轻量级在线关联机制维护实例记忆。两个流独立运行，避免几何更新干扰语义学习，增量累积的高斯场支持可微渲染和持续的场景更新。

\paragraph{边缘设备在线学习}

Zaino et al.（2026.03）~\cite{2603_08499} 针对变分贝叶斯高斯重建（VBGS）在商用边缘设备上的部署问题，提出精度自适应优化框架。通过内核融合减少中间张量，自动混合精度搜索在有限误差下分配操作级精度。在 Jetson Orin Nano 上首次实现 NVS 训练，峰值内存从 9.44 GB 降至 1.11 GB，训练时间从约 234 分钟降至约 61 分钟。

\subsubsection{场景变化检测与更新}

本路线介绍检测场景变化并仅对变化区域进行局部更新的持续学习方法，避免每次变化都从头重建整个 3DGS 模型。

现实场景不是静止的——物体移动、建筑变化、季节更替。如果每次变化都从头重建整个 3DGS 模型，代价显然不可接受。更合理的策略是：检测哪些区域发生了变化，只对变化区域进行局部更新。这一思路引出三个子问题：（1）如何高效检测变化区域；（2）如何只更新变化部分，同时不破坏未变化区域的质量；（3）如何维护场景的时间历史，以支持时间维度的查询和分析。

\paragraph{DINOv2 变化检测与空间约束局部优化}

\textbf{基线做法与问题。}\,场景发生变化后，朴素策略是对新观测重新运行完整的 3DGS 优化。这存在三个缺陷：（1）丢弃已有重建 $G^{t-1}$ 中未变化区域的信息，导致计算浪费；（2）要求重新采集整个场景，降低灵活性；（3）在观测不足的区域，已有重建被擦除，造成信息丢失。此前的隐式持续学习方法 CL-NeRF（2025.06）~\cite{2506_21117} 虽能处理时间演化，但面临灾难性遗忘和历史状态不可恢复的问题——所有时间步信息被压缩进单一全局表示，容量受限且无法精确回溯。

\textbf{改了什么。}\,CL-Splats 提出基于 3D 高斯的持续学习框架，通过三阶段流水线实现高效局部更新：（1）基于 DINOv2 特征的鲁棒 2D 变化检测；（2）多数投票将 2D 掩码提升为精确的 3D 变化区域；（3）空间约束的局部优化，仅更新变化高斯。此外，该方法支持并发更新（合并多组独立变化）和精确的历史场景恢复。

\textbf{实现原理。}\,变化检测阶段，利用 DINOv2 提取当前观测 $I_i^t$ 与已有重建渲染图 $\hat{I}_i^{t-1}$ 的逐 patch 特征，通过余弦相似度阈值生成 2D 二值变化掩码：

\begin{equation}
    M_i^t = \begin{cases}
        1, & \text{sim}(\hat{F}_i^{t-1}, F_i^t)_1 < \tau_1 \\
        0, & \text{otherwise}
    \end{cases}
\end{equation}

掩码经膨胀填充噪声孔洞后，通过多数投票机制提升到 3D：统计每个高斯在多视角中落入掩码的次数 $c_i$ 与落入画外的次数 $o_i$，当 $0.43\, o_i < |\mathcal{I}^t| < 2\, c_i$ 时归为变化集合 $O^t$。对于新增物体（该区域可能无现存高斯），通过 K-Means 聚类和高斯混合模型在变化区域采样新点。

局部优化阶段，CL-Splats 冻结 $O^t$ 以外的所有高斯，仅对变化区域执行优化。为防止高斯在优化过程中漂移出变化区域，用 HDBSCAN（Hierarchical Density-Based Spatial Clustering of Applications with Noise）对 $O^t$ 聚类后拟合包围球体，每 15 次迭代剪枝离开球体的高斯。同时，修改 3DGS CUDA 核，动态投影变化高斯生成渲染掩码，仅对覆盖的 tile 执行前向和反向计算——确保梯度与全局优化完全一致，但大幅减少计算量。

\textbf{设计动机。}\,选择持续学习而非重新训练的根本原因在于：现实场景的变化通常是局部的，未变化区域占据场景的绝大部分。局部优化利用这一先验，将计算集中在真正需要更新的区域。直接用 2D 掩码约束光度损失（3DGS+M 基线）会导致严重伪影——因为 2D 掩码缺乏 3D 空间感知，会误更新不属于变化区域的高斯。多数投票将变化检测从像素空间提升到高斯空间，是解决此问题的关键。选择球体而非包围盒作为空间约束基元，则因为球体的包含检测比包围盒少约 3 倍 FLOPS。显式存储变化高斯及其索引（而非像 NeRF 方法依赖时间步条件编码）使得历史恢复无需冗余存储，从根本上避免了灾难性遗忘。

\textbf{关键参数与性能。}\,在 CL-Splats 合成数据集上，方法仅用 5K 迭代（40 秒，RTX Quadro 6000）即收敛至 PSNR 40.1 dB，而 CL-NeRF（2025.06）~\cite{2506_21117} 需要 25K 迭代（50 分钟）且仅达 30.1 dB——速度差距约 75 倍。局部优化核平均减少 60\% 优化时间。历史恢复方面，每步仅需额外存储 36 MB（仅保存变化高斯 $O^t$ 及索引 $I^t$），对比朴素全场景存储的 1173 MB，节省约 97\%。方法还支持并发更新：可将多组独立变化并行处理后统一合并到 $G^{t-1}$，这是 NeRF 类隐式方法因表示限制而无法实现的。

\paragraph{时间哈希嵌入与生成式回放的持续更新}

GaussianUpdate（2025.08）~\cite{2508_08867} 将 3D 高斯表示与持续学习结合，显式建模不同类型的变化。多阶段更新策略有效处理新数据更新，可见性感知持续学习方法配合生成式回放（generative replay）实现自我感知更新而无需存储历史图像。

具体而言，GaussianUpdate 包含三个阶段的渐进更新策略，避免外观和几何属性的紧耦合干扰。第一阶段引入时间哈希嵌入建模场景在时间维度的变化，为每个高斯添加时间相关的属性偏移。第二阶段通过可见性分析判断哪些区域发生了变化，只对变化区域进行更新以避免全局重训练。第三阶段的生成式回放机制不存储历史图像，而是利用已学习的高斯模型渲染历史图像作为回放数据，有效防止灾难性遗忘。不同于 4D 建模方法处理连续动态，GaussianUpdate 专注于时间上任意的突发变化（如物体突然出现或消失），设计了自适应的高斯管理机制来处理外观变化和物体级别的出现与消失。

\paragraph{TSDF-3DGS 双图与变化感知密度控制}

VG-Mapping（2025.10）~\cite{2510_09962} 面向机器人反复穿越的半静态场景，提出变体感知密度控制策略——将高斯密度调节与优化解耦，识别变化区域引导初始化和剪枝，避免使用过期信息定义后续优化的起点。

\textbf{动机。}在仓库、工厂等半静态环境中，物体经常被移动或替换，标准 3DGS 的自适应密度控制依赖历史累积的梯度统计，场景变化后这些过期梯度会误导新高斯的生成位置。

该系统采用双线程架构（跟踪+建图）：跟踪线程使用 VINS-Mono 不受半静态变化影响，建图线程结合 TSDF 体素地图和 3DGS 地图。TSDF 地图提供细粒度几何信息，增强后对半静态场景变化的响应——当当前观测与全局地图的 TSDF 值差异超过阈值时，设置负权重快速淘汰过时体素。变化感知密度控制（VDC）通过比较当前观测和已渲染图像的变化掩码识别变化区域，在变化区域引导新高斯初始化（不依赖陈旧梯度）和旧高斯剪枝（移除不再存在的结构）。VDC 的关键创新在于避免了传统自适应密度控制在半静态场景中因陈旧梯度导致的优化起点偏移问题。

\paragraph{动态场景适应与关键帧过时管理}

GaME（2025.06）~\cite{2506_06909} 引入动态场景适应机制持续更新 3DGS 以反映最新变化，并设计关键帧管理机制丢弃过时观测同时保留尽可能多的信息。

\textbf{动机。}当场景中的物体被添加、移除或移动后，已有 3DGS 模型仍然保留旧物体的残影，且用于训练这些残影的旧关键帧继续提供错误的梯度信号，如果不加管理会导致新观测被旧信息覆盖。在 PSNR 上提升 29.7\%，L1 深度误差改善 3 倍。

动态场景适应（DSA）机制处理三种场景变化——添加、移动和移除刚性几何体。检测新增几何时，当观测深度比渲染深度更近，标记对应高斯为被遮挡并种子化新高斯；检测移除或移动几何时，当观测深度比渲染深度更远，降低被遮挡高斯的不透明度。该方法不依赖物体级建模，而是在整体高斯表示上进行"按需"操作。关键帧管理系统使用深度渲染误差和颜色一致性检测过时关键帧，通过加权策略在移除过时信息和保留有用信息之间平衡。在线优化使用 L1+SSIM 颜色损失和 L1 深度损失，加上各向同性正则化防止退化高斯形状。

\paragraph{干扰置信度初始化的跨时间渐进优化}

Cross-Temporal 3DGS（2025.12）~\cite{2512_00534} 提出跨时间相机对齐、基于干扰的置信度初始化和渐进式跨时间优化三阶段框架，使用稀疏图像和先前场景先验高效重建和更新跨时间段的 3D 场景。该方法支持非连续采集，既能用新稀疏视角更新现有场景，也能借助当前采集恢复过去场景。

\paragraph{模板高斯与时变环境精化的长期编年史}

LTGS（2025.10）~\cite{2510_09881} 从稀疏视角更新构建长期高斯场景编年史，将物体构建为模板高斯作为共享物体轨道的结构化可复用先验。物体模板通过精化管线根据少量观测适应时变环境。训练后，框架通过简单变换泛化到多个时间步，显著增强 3D 环境时间演化的可扩展性。

具体流程包括四个步骤：首先使用鲁棒分层定位管线（hloc）估计新图像的相机位姿，结合语义和光度标准比较渲染图像与拍摄图像的差异以检测物体级别变化。其次利用基础模型（MASt3R、SAM）进行密集局部匹配和实例识别，关联不同时间戳的物体。然后分别训练物体模板高斯和静态背景高斯，物体模板作为可重用结构先验通过平移加旋转变换匹配不同时间步的观测。最后通过精化管线，给定少量新视角观测调制模板先验适应时间变化的环境。训练后框架通过简单变换泛化到多个时间步，支持物体插入、移除、替换和迁移等多种变化类型。

\paragraph{在线场景变化检测}

Galappaththige et al.（2025.11）~\cite{2511_12370} 提出首个位姿无关、无标签、多视角一致的在线场景变化检测方法，运行速度超过 10 FPS 且性能超越最佳离线方法。自监督融合损失从多个线索和观测推断场景变化，PnP 快速位姿估计和变化引导的 3DGS 更新策略确保实时性能。

\paragraph{符号距离差分与多视角投票的变化建模}

SCaR-3D（2025.12）~\cite{2512_22830} 利用 3DGS 的一致性性质，通过符号距离的 2D 差分模块和多视角聚合（投票+剪枝）鲁棒分离变化前后的状态。持续场景重建策略选择性更新动态区域同时保留不变区域。

具体而言，SCaR-3D 包含四个步骤：首先使用 SfM 联合估计变化前后图像的相机姿态以统一坐标系，并训练变化前的 3DGS 模型。然后使用 EfficientSAM 提取特征，计算渲染的变化前图像与真实变化后图像之间的 PCA 主导方向上的符号距离，捕获变化的方向性（新增或移除）。接着将 2D 变化检测结果通过 3DGS 的多视角一致性投影到 3D 空间，通过投票和剪枝机制去除噪声和不一致检测，生成稳健的 3D 变化掩码并分离变化前后的对象状态。最后根据变化掩码选择性更新 3DGS 中的动态区域同时保留不变区域。该方法还贡献了 CCS3D 合成数据集用于评测，在复杂变化类型（平移、插入、移除）下表现突出。

\paragraph{增量动态更新}

IDU（2025.08）~\cite{2508_17579} 管线高效更新现有 3DGS 重建——通过相机位姿估计对齐新图像，变化检测定位场景修改，3D 生成式 AI 创建新元素并集成到现有模型。每次更新聚焦单个新物体，显著减少更新时间和人力。

\paragraph{高斯地图共享}

GS-Share（2025.10）~\cite{2510_02884} 提出基于锚点的全局地图构建、虚拟图像增强和增量地图更新的地图共享系统，在 PSNR、LPIPS 和深度 L1 上分别提升 11\%、22\% 和 74\%，地图传输开销减少 36\%。

\paragraph{在线新视角选择}

Lee et al.（2025.08）~\cite{2508_14014} 通过在线重建质量分析选择最优非关键帧进行额外训练，结合关键帧和选定非关键帧从多样视角精化不完整区域，显著提升完整性。

\paragraph{3DGS 场景流式传输}

Siekkinen and K\"{a}m\"{a}r\"{a}inen（2026.04）~\cite{2604_02851} 探索将实时渲染场景作为 3D 高斯流式传输，而非传统 2D 视频流。服务器持续优化 3DGS 模型，通过完整模型快照和增量更新支持重光照和刚体动态。客户端从接收的表示中渲染自身视角，改善延迟补偿的视角灵活性。

\paragraph{结构损伤可视化}

Wang et al.（2026.02）~\cite{2602_16713} 将 3DGS 引入土木基础设施数字孪生，利用 GS 3D 重建可视化 2D 损伤分割结果，开发多尺度重建策略平衡效率与损伤细节，并支持数字孪生随损伤演化持续更新。

\begin{takeaway}
在线与增量重建使 3DGS 从离线处理走向实时场景理解：滑动窗口关键帧选择、增量致密化和遗忘机制解决了数据流式到达下的计算约束，为机器人探索和自动驾驶的实时建图奠定了基础。
\end{takeaway}

% ======================================================================
\section{分层与LOD表示}
\label{sec:14_3}

当场景规模从单栋建筑扩展到整个城市时，在任何视点都加载全部高斯基元既不现实也不必要——远处的建筑只需粗略表示，近处的细节需要精细呈现。这就是细节层次（Level-of-Detail, LOD）技术的核心动机：根据视点距离和可用资源动态选择合适的表示精度，在保持视觉质量的同时控制计算和内存开销。LOD 技术在传统计算机图形学中已有悠久历史，但 3DGS 的半透明高斯基元和可微渲染管线为 LOD 设计带来了独特挑战。

\begin{figure}[htbp]
\centering
\resizebox{!}{0.4\textheight}{\begin{tikzpicture}[
  box/.style={rectangle, rounded corners=4pt, draw=#1, fill=#1!10,
    text width=3.2cm, align=center, font=\small, minimum height=1cm},
  sbox/.style={rectangle, rounded corners=3pt, draw=#1!60, fill=#1!5,
    font=\scriptsize, text width=2.4cm, align=center, minimum height=0.5cm},
  arr/.style={-{Stealth[length=6pt]}, thick, techgray},
  darr/.style={-{Stealth[length=6pt]}, thick, techgray, dashed},
  >=Stealth
]

  \node[box=techblue] (h0) at (3.1, 8.1) {全局场景(城市级)};
  \node[box=techgreen] (h1) at (3.1, 5.6) {区域划分(街区级)};
  \node[box=techorange] (h2) at (3.1, 3.1) {局部细节(建筑级)};
  \node[box=techpurple] (h3) at (3.1, 0.6) {LOD渲染(视距自适应)};
  \draw[arr] (h0) -- (h1);
  \draw[arr] (h1) -- (h2);
  \draw[arr] (h2) -- (h3);
  \draw[darr] (h3) to[bend right=30] node[midway, font=\scriptsize, text=techgray] {核外调度} (h1);
\end{tikzpicture}}
\caption{大规模场景分层表示与LOD渲染}
\label{fig:ch14_scale_hierarchy}
\end{figure}

\subsection{层次化结构设计}
\label{ss:14_5}

本节介绍两条技术路线：（1）层次化高斯结构；（2）LOD 渲染优化。

\subsubsection{层次化高斯结构}

本路线介绍将高斯基元组织为层次结构以支持多级细节渲染的表示方法。

层次化高斯结构是 LOD 实现的基础。核心思想是将高斯基元组织为层次结构，高层节点表示粗略几何，低层节点提供精细细节。

\paragraph{层次合并与中间节点优化的 LOD 方案}

Kerbl et al.（2024.06）~\cite{2406_12080} 提出的层次 3D 高斯表示是 LOD 领域的开创性工作。该方法将高斯基元组织为层次结构，高层节点通过合并低层高斯得到，在渲染远处内容时使用高层粗略表示以提升效率。训练采用分治策略：将大型场景分割为独立块分别训练，然后合并为层次结构并进一步优化中间节点的视觉质量。

层次结构的核心构建方式是：对每个父节点高斯，沿其局部坐标系的主轴方向放置 $K_v$ 个子节点。子节点位置偏移由父节点的旋转矩阵和缩放向量决定，几何尺寸按比例因子 $\gamma < 1$ 逐层缩小，同时继承父节点的朝向、颜色和不透明度，从而形成天然的粗到细层次结构。中间节点的优化是该工作的关键创新——不仅优化叶节点，还对合并产生的内部节点进行独立梯度优化，使其视觉质量远超简单继承父节点参数的朴素策略。

训练阶段引入多项正则化技术应对大规模稀疏数据，包括：基于屏幕空间覆盖率的自适应密度控制、高斯尺寸下界约束（防止退化）、以及针对远处高斯的不透明度衰减策略。LOD 选择基于每个高斯的 2D 投影半径与屏幕空间阈值的比较，实现自动的细节层级切换和层级间平滑过渡。该方法展示了使用简单设备采集数万张图像、覆盖数公里轨迹、持续长达一小时的实时渲染结果。

\paragraph{累积式八叉树 LOD 与视点自适应锚点选择}

\textbf{背景与动机。}标准 3DGS 的一个根本缺陷在于高斯分布与场景真实几何结构的错位（misalignment）——高斯基元的位置由训练视角的光度损失驱动，而非由场景本身的几何结构引导。这种错位带来两个瓶颈问题：一是渲染泛化性差，训练视角之外的视角质量显著下降；二是产生大量冗余重叠基元，无法高效表示场景细节。

\textbf{前驱方法。}Scaffold-GS（2023.12）~\cite{2312_00109} 是首个系统性解决结构对齐问题的方法：将 SfM 稀疏点云体素化为规则网格，每个体素中心作为锚点（anchor），配备 32 维可学习局部特征 $f_v$、缩放因子 $l_v$ 和 $k$ 个可学习偏移量 $\{O_v^k\}$。在渲染时，根据相机-锚点距离和方向动态解码神经高斯的全部属性（位置、不透明度、协方差、颜色），使高斯的生长和剪枝被约束在锚点的局部邻域内，而非无限制地在全场景扩散。

尽管 Scaffold-GS 通过规则间距特征网格改善了结构对齐，Mip-Splatting 通过 3D 平滑和 2D Mip 滤波缓解了高斯冗余，但这些方法在大规模场景中仍然依赖基于可见性的基元筛选——将视锥体内所有基元纳入渲染，而不考虑其投影尺寸。这导致远处的微小细节也被渲染，造成不必要的计算开销和渲染速度的大幅波动。Octree-GS（2024.03）~\cite{2403_17898} 正是针对这一核心矛盾，提出用八叉树层次结构组织锚点，实现视点距离自适应的 LOD 选择。

\textbf{核心思想：LOD 结构化锚点。}Octree-GS 受八叉树数据结构启发，将高斯基元的管理锚点（anchor）组织为层次化的八叉树结构。每个锚点位于不同粒度体素网格的中心——高 LOD 层级的锚点对应更小的体素尺寸，低 LOD 层级对应更大的体素。关键设计是 LOD 的累积性（cumulative）：渲染 LOD $K$ 时，实际光栅化从 LOD 0 到 LOD $K$ 的所有高斯基元。这意味着远处物体使用粗粒度锚点发出的少量大高斯表示，近处物体叠加细粒度锚点提供的精细细节。每个锚点还附带一个可学习的 LOD 偏置（LOD bias）$\Delta L$，用于补偿局部场景复杂性——例如物体锐利边缘区域需要比纯距离判断更高的细节层次。

\textbf{技术细节。}Octree-GS 的技术实现包含四个关键组件。

（1）\textit{锚点初始化与八叉树构建}。给定 SfM 稀疏点云 $P$，首先根据相机-点距离范围确定八叉树层数 $K = \lfloor\log_2(d_{\max}/d_{\min})\rfloor + 1$，其中 $d_{\max}$ 和 $d_{\min}$ 分别为剔除离群值后的最大和最小观测距离（阈值 $r_d = 0.999$）。然后以基础体素大小 $\delta$ 构建各层体素网格，LOD $L$ 的体素大小为 $\delta / 2^L$，SfM 点落入哪个体素就在该体素中心初始化一个锚点。

（2）\textit{视点自适应锚点选择}。对于视角 $i$ 和锚点 $j$，根据观测距离 $d_{ij}$ 动态计算 LOD 级别 $\hat{L}_{ij} = \lfloor\log_2(d_{\max}/(d_{ij} \cdot s))\rfloor + \Delta L_j$，其中 $s$ 为焦距缩放因子（处理多分辨率相机）。直观理解：距离越远，$\hat{L}_{ij}$ 越小，只用粗粒度锚点；距离越近，$\hat{L}_{ij}$ 越大，叠加更多细粒度锚点。锚点被选中当且仅当 $L_j \leq \hat{L}_{ij}$。为确保 LOD 过渡无视觉伪影，采用分段线性不透明度混合（opacity blending）——额外选中满足 $L_j = \hat{L}_{ij} + 1$ 的近满足锚点，并将其高斯不透明度乘以过渡因子 $\hat{L}_{ij}^* - \lfloor\hat{L}_{ij}^*\rfloor$，避免细节层次切换时的突然跳变。

（3）\textit{LOD 感知致密化}。采用"下一层"生长（next-level grow）策略：当某锚点生成的高斯梯度 $\bar{\tau}_g$ 超过阈值时，在更高 LOD 层级的空体素中创建新锚点。为防止过度生长，阈值随层级单调增加 $\tau_g^L = \tau_g \cdot 2^{\Delta L}$。同时引入视角频率（view-frequency）剪枝——移除在训练视角中很少被选中的锚点，有效消除浮游体（floaters）并降低存储。

（4）\textit{渐进式粗到精训练}。从 LOD $\lfloor K/2 \rfloor$ 层开始训练，每隔 $N_i$ 次迭代激活一个新 LOD 层级，粗层级获得更多训练时间（$N_{i-1} = N_i \cdot \gamma_1$，$\gamma_1$ 为增长因子）。这种策略使各 LOD 层级能更忠实地分解到对应的尺度层级，减少锚点冗余。此外，针对野外场景外观不一致问题，为每个锚点引入基于 GLO（Generative Latent Optimization）的可学习外观编码，通过颜色 MLP 解码视角依赖颜色。

Octree-GS 的框架具有通用性——锚点可发射不同类型的高斯基元，包括显式 2D/3D 高斯（带独立可学习属性）或神经高斯（从锚点特征通过 MLP 解码），分别称为 Our-2D-GS、Our-3D-GS 和 Our-Scaffold-GS。

\textbf{实验验证。}Octree-GS 在 21 个小规模场景和 7 个大规模场景上进行了全面评测。大规模场景包括 MatrixCity Block\_Small/Block\_All、MegaNERF 的 Rubble/Building、UrbanScene3D 的 Residence/Sci-Art 以及 Hierarchical-GS 的 SmallCity 街景数据集。

核心发现：（1）渲染质量在所有数据集上一致优于 3D-GS、Scaffold-GS 和 Mip-Splatting 基线，在细粒度细节和弱纹理区域表现尤为突出；（2）渲染效率方面，大规模场景和极端视角序列中的渲染速度比 Scaffold-GS 快近 10 倍，这归功于 LOD 结构设计——远处区域自动使用更少高斯；（3）存储方面，view-frequency 剪枝策略使每视角平均渲染高斯数量显著减少，在大规模场景中存储开销降低尤为明显；（4）在大规模场景中训练时间也更短——因为渐进式训练中不需要同时优化所有 LOD 层级的所有高斯，Our-Scaffold-GS 在 Mip-NeRF 360 上仅需 23 分钟（Scaffold-GS 需 29 分钟）；（5）多尺度（BungeeNeRF 从地面到卫星高度四个尺度）和多分辨率实验证明了 LOD 结构对不同缩放和分辨率变化的自适应鲁棒性；（6）消融实验验证了四个关键组件（next-level grow、LOD bias、渐进式训练、view-frequency）各自的必要性——移除 next-level grow 导致视觉质量显著下降，移除渐进式训练导致 LOD 层级间严重重叠。

\textbf{局限与意义。}Octree-GS 的局限性在于：八叉树构建和渐进式训练的超参数仍需手动调整，各 LOD 层级的锚点平衡和训练迭代激活时机需要经验性设定；依赖精确相机位姿；缺乏几何表面支持。尽管如此，Octree-GS 作为首个系统性地将 LOD 引入高斯表示的方法，其八叉树锚点结构、视点自适应选择和渐进式训练范式为后续大量工作（PyGS、LODGE、FilterGS 等）奠定了基础架构，是 LOD 高斯表示领域的里程碑。

\paragraph{NeRF 初始化金字塔与可学习簇级权重}

PyGS（2024.05）~\cite{2405_16829} 将高斯基元按金字塔层次组织：顶层由少量大高斯组成，每层包含更密集的小高斯。在初始化方面，PyGS 首先训练一个快速基于网格的 NeRF（约 20 分钟）替代 COLMAP 生成点云——从训练相机投射射线并在估计终止点采样，生成粗点云后再以递增采样频率抽取 $L$ 个子集分别初始化不同层级的高斯，天然形成多尺度金字塔结构。这解决了 COLMAP 在大规模数据集上计算昂贵且远处重建不完整的问题。在渲染方面，PyGS 对高斯按簇分组（使用 mini-batch K-means 建立簇中心），每个簇配备可学习嵌入特征，通过紧凑加权网络 $w_{k,n} = \sigma(F_w(\text{cat}(E_k, \gamma(x_n)))) \in \mathbb{R}^L$ 根据相机位置动态决定该簇中各金字塔级别的影响力权重。渲染颜色公式将权重直接融入 $\alpha$-混合：$C = \sum_{i=1}^{N}(c_i + c_i') w_{k_i,n}[l_i] \alpha_i \prod_{j=1}^{i-1}(1 - w_{k_j,n}[l_j]\alpha_j)$，其中 $c_i'$ 是基于每相机外观嵌入的逐簇颜色校正。PyGS 在多个大规模数据集上实现渲染速度比当时最优方法快 400 倍以上。金字塔权重的关键优势在于：同一空间区域的多个金字塔级别通过加权混合而非硬切换实现尺度过渡，相机移动时各级别影响平滑变化，避免了离散 LOD 切换产生的视觉跳变。基于 NeRF 的初始化策略使金字塔各级别的点云密度自然匹配对应尺度的表达需求——低级别对应稀疏大尺度结构，高级别填充精细细节。

\paragraph{重要性驱动高斯剪枝与法向先验的块级精化}

HRGS（2025.06）~\cite{2506_14229} 提出内存高效的层次化框架：先从低分辨率数据生成全局粗略高斯表示，再分割场景为多个块用高分辨率数据精化。高斯分割将不规则场景归一化为有界立方空间，训练数据分割仅保留相关观测。粗略高斯先验引导块精化确保相邻块的无缝融合。重要性驱动高斯剪枝（IDGP）计算每个高斯的重要性分数移除低贡献基元。具体地，IDGP 为每个高斯 $G_k$ 定义重要性分数 $S_k = \sum_{v \in V_k} w_{k,v} / N_v$，其中 $V_k$ 为观测到该高斯的视角集合，$w_{k,v}$ 为其在视角 $v$ 下的加权渲染贡献，$N_v$ 为该视角下总高斯数。重要性分数低于阈值的低贡献基元被迭代移除，在加速收敛的同时显著降低内存使用。此外，HRGS 引入来自预训练法向估计模型的法向先验损失，将估计法向与高斯 surfel 法向对齐，增强表面重建质量。实验在 Mip-NeRF 360、Tanks and Temples 和 Deep Blending 三个基准上验证了该方法在超高分辨率（如 5K）新视角合成和表面重建任务上的有效性。

\subsubsection{LOD 渲染优化}

本路线介绍渲染时从层次结构中高效选择合适细节级别的优化方法。

建立层次结构只是 LOD 系统的一半工作；另一半——也是实际决定渲染帧率的关键——是渲染时如何高效地从层次结构中选择合适的 LOD 级别。一个直观的想法是逐点计算距离后选择细节层次，但在数千万高斯的场景中，这种逐点操作的代价不可接受。本节的方法从不同角度解决这一效率问题：FilterGS 彻底消除树遍历，LODGE 利用深度信息辅助 LOD 构建，Nebula 和 Voyager 分别面向 VR 和移动端的特殊约束。

\paragraph{深度感知 3D 平滑滤波与双块不透明度混合}

LODGE（2025.05）~\cite{2505_23158} 提出一种面向大规模场景且兼容移动端设备的层次 LOD 表示与分块缓存方案。该方法从已有 3DGS 重建之上构建 LOD 结构，而非采用从粗到细的渐进式训练，避免了粗粒度初始化过于稀疏导致细化失败的问题。

\textbf{基线做法与不足。}在 LODGE 之前，Octree-GS、FLoD 和 H3DGS 等方法虽然引入了 LOD 层次结构，但存在三个核心限制：（1）渲染时需要逐帧重新计算活跃高斯子集（active Gaussians），这一动态选择过程带来不可忽视的计算开销；（2）所有 LOD 层级的全部高斯必须常驻 GPU 内存，使得总高斯数量甚至超过原始 3DGS，在移动或低端设备上根本无法运行；（3）深度阈值等超参数需要逐场景手动调整。这些方法主要关注渲染速度的提升，却未从根本上解决内存占用问题。

\textbf{深度感知 3D 平滑滤波器。}LODGE 的核心创新是借鉴 Mip-Splatting 的 3D 平滑滤波思想，但将其用途从抗锯齿转向 LOD 构建。其出发点是 Nyquist 采样定理：当像素在深度 $d$ 处对应的 3D 世界空间采样间隔为 $T = d/f$（$f$ 为焦距）时，尺寸小于 $2T$ 的高斯只会产生混叠伪影，同时浪费内存和计算。因此，对于最小观测距离为 $d_l$ 的 LOD 层级 $G^{(l)}$，LODGE 将 $G^{(0)}$ 中的每个高斯与一个深度相关的平滑滤波器卷积，得到新的高斯函数：

\begin{equation}
\hat{G}(x) = \frac{|\Sigma + s \cdot d_l \cdot f \cdot I|^{-1/2}}{(2\pi)^{3/2}} \exp\left\{-\frac{1}{2}(x-\mu)^T (\Sigma + s \cdot d_l \cdot f \cdot I)^{-1}(x-\mu)\right\}
\end{equation}

其中 $\mu$ 和 $\Sigma$ 为原始高斯的均值和协方差，$s$ 是控制平滑强度的超参数，$I$ 为单位矩阵。该滤波器的效果是使远距离 LOD 层级中的高斯尺寸下界随距离增大，自然地使远处表示趋向低频、模糊，与人类视觉中近处清晰远处模糊的感知规律一致。与简单按距离阈值裁剪高斯的方法不同，这种距离驱动的平滑不会丢失远处场景的整体结构信息，而是在保留大尺度几何的同时抑制无法被像素采样的高频噪声。

\textbf{LOD 级别的迭代构建。}给定深度阈值序列 $d_1 < d_2 < \cdots < d_L$，LODGE 从最精细层级 $G^{(0)}$ 开始，由粗到精迭代构建各 LOD 层级。在第 $l$ 步，先对 $G^{(0)}$ 中所有高斯施加对应深度 $d_l$ 的 3D 平滑滤波器（仅修改高斯函数形态，不改变参数），然后执行三轮重要性剪枝（阈值分别为 0.2、0.6 和一个可调超参数），每轮剪枝后进行 1000 步微调优化。微调时采用 LOD 渲染程序，但将深度 $d_l$ 替换为均匀分布 $U(0.7d_l, 1.3d_l)$ 的随机采样值，增强对训练轨迹之外相机位置的鲁棒性。整个构建过程仅需 $N_{\text{levels}} \times 3 \times 1000$ 步优化迭代。

\textbf{活跃高斯选择与自动阈值。}渲染时，对相机中心 $c$，LODGE 从各层级选择满足距离条件的高斯子集：

\begin{equation}
\mathcal{A}(c) = \bigcup_{l=0}^{L-1}\left\{g_i \in G^{(l)} : d_l \leq \|\mu_i^{(l)} - c\|_2 < d_{l+1}\right\}
\end{equation}

其中 $d_0 = 0$、$d_L = \infty$。这使远处区域自动使用更低细节的高斯集，显著减少了像素级可见高斯数量的长尾分布——正是这些长尾像素拖慢了基于 $16 \times 16$ 瓦片的光栅化（同一瓦片内所有线程必须等待最慢线程完成）。深度阈值 $d_l$ 的选择采用贪心自动搜索：从 $d_1$ 开始迭代添加阈值，每次在训练视角子集上评估每个瓦片的平均高斯数量作为代价函数，选择使代价最小的阈值。实验观察到代价函数的等值线近似线性（$d_2 \approx a \cdot d_1 + b$），因此搜索复杂度降为线性。

\textbf{场景分块与预计算活跃高斯。}为进一步降低内存占用，LODGE 将场景分割为空间块（chunks）。具体地，对训练相机位置执行 K-means 聚类，聚类数量由公式 $N_{\text{clusters}} = 4 \cdot d_1 / \frac{1}{N_c}\sum_i\|c_i - \bar{c}\|$ 自动确定（$\bar{c}$ 为相机位置均值），在户外场景中对应约 5 米的块大小。关键设计是：每个块的活跃高斯集并非该块空间边界内的高斯子集，而是以块中心为虚拟相机位置、按上述距离公式选择的完整场景表示——近处以高细节、远处以低细节。这样渲染时只需加载最近块的预计算活跃高斯集，无需逐帧重新计算。每个块还额外执行可见性过滤：在块内训练视角的基础上添加随机扰动方向的额外视角，计算重要性分数并剪枝低贡献高斯。

\textbf{不透明度混合与块间过渡。}直接使用最近块的活跃高斯集会在相机跨越块边界时产生明显的清晰度跳变伪影。LODGE 提出基于双块的不透明度混合策略：渲染时找到距离当前相机位置 $c$ 最近的两个块中心 $m_f$（高斯 $i$ 所属块）和 $m_o$（另一块），对两个块活跃高斯集合并集中的高斯，按如下方式调制不透明度：

\begin{equation}
\hat{o}_i = o_i \cdot t, \quad t = \min\left(1, \max(0, \lambda)\right), \quad \lambda = \frac{(c - m_o)^T(m_f - m_o)}{\|m_o - m_f\|_2^2}
\end{equation}

其中 $\lambda$ 是 $(c - m_o)$ 在两块中心连线上的投影长度归一化值，使用投影长度而非欧氏距离是为了在相机不经过块中心时仍能实现平滑过渡。两个块中共同拥有的高斯保持原始不透明度不变，仅对非共有高斯施加调制。实际渲染中，当相机越过当前块中心后，系统在后台异步加载下一个块的活跃高斯——此时当前块的高斯不透明度已接近 1，因此加载延迟不会造成可见伪影。

\textbf{设计动机总结。}LODGE 的设计围绕三个核心洞察：（1）距离驱动平滑优于简单阈值裁剪——前者保留远处场景的大尺度结构，后者直接丢弃信息；（2）空间分块预计算优于逐帧动态选择——前者将计算开销从渲染时转移到预处理阶段，且将 GPU 内存占用限制为两个块的活跃高斯而非全部 LOD 高斯；（3）双块不透明度混合优于简单切换——前者在视觉上实现平滑过渡，后者在块边界产生清晰度跳变。这三者的结合使 LODGE 成为首个同时优化渲染速度、内存占用和视觉质量，且可在移动设备上实时运行的 LOD 方法。

\textbf{实验验证。}LODGE 在两组大规模数据集上评测：Hierarchical 3DGS 数据集的 SmallCity 和 Campus（户外街景），以及 Zip-NeRF 数据集的 Alameda、London 和 NYC（室内大场景），每个场景约 1000--2000 张图像。

与 Octree-GS、CityGS、H3DGS、FLoD 等方法相比，LODGE 的核心实验结论包括：（1）渲染质量方面，在所有场景上一致取得最优或次优的 PSNR/SSIM/LPIPS，同时近处细节和远处保真度均优于 Octree-GS（后者近处模糊、远处有色偏）和 H3DGS（后者远处出现锯齿状几何）；（2）渲染速度方面，比 Octree-GS 和 CityGS 快约 2 倍，原因是 LODGE 无需逐帧计算活跃高斯或运行 MLP 解码；（3）内存方面，GPU 中只需保持两个块的活跃高斯，而非所有 LOD 层的全部高斯，使移动端渲染成为可能；（4）移动设备实测方面，在 iPhone 13 Mini、iPhone 15 Pro、MacBook Air M3 和 HP Chromebook 上均可实时渲染，而 H3DGS 因内存超限无法在 iPhone 上运行；（5）消融实验验证了各组件的必要性：添加一个 LOD 层级即可获得约 3 倍渲染加速（代价仅是极小的 PSNR 下降），自动阈值选择优于手动设定（如 $d = 15, 70$m），可见性过滤进一步减少加载高斯数量，而不透明度混合消除了块边界的清晰度跳变伪影。

LODGE 的局限性在于：假设块间高斯加载可通过优化 Web 服务和压缩协议高效完成（目前未实现），且两块混合策略在理论上可能不足以覆盖三个以上块交汇的区域（实验中未观察到问题）。

\paragraph{Root\&Leaf 与 Ancestor 双级并行过滤}

FilterGS（2026.03）~\cite{2603_23891} 指出层次 LOD 方法中串行遍历消耗超过 60\% 的渲染时间——即使层次结构已经建立，逐节点判断"这个锚点是否该被渲染"的串行过程本身就成为瓶颈。FilterGS 提出并行过滤机制——使用两个互补滤波器高效选择高斯元素而无需树遍历。第一个是 Root \& Leaf (R\&L) 滤波器：对视锥体内所有高斯并行计算 2D 投影协方差的特征值，定义像素半径 $R_{2D} = 3\sqrt{\lambda_{\max}}$，当 $R_{2D}$ 小于阈值 $\theta_R$ 时保留该节点（包括叶节点和投影足够小的中间节点），这一步骤完全并行化在 GPU 上执行。第二个是 Ancestor 滤波器：沿祖先路径剪枝冗余节点，确保每个分支仅保留最合适的一个层级。两级滤波器级联的结果是一组可直接送入标准光栅化器的完整高斯基元，无需任何解码或预处理。同时提出 GTC（Gaussian-Tile Count）度量量化高斯-瓦片键值对冗余，基于此的场景自适应高斯收缩策略通过缩小高斯覆盖范围有效减少冗余键值对，在不损失视觉质量的前提下进一步降低光栅化负载。在多个大规模数据集上达到当时最优渲染速度。

\paragraph{时间感知 LOD 差量传输与立体光栅化}

Nebula（2025.12）~\cite{2512_20495} 解决城市级 3DGS 在 VR 中的可扩展性问题。核心创新是流式传输 LOD 搜索后的中间结果而非视频，将云端-客户端数据通信量降低至原来的约 5\%。云端采用时间感知 LOD 搜索利用帧间时间一致性减少冗余数据访问：通过对比前后帧的活跃 LOD 切割（cut），仅传输帧间差量高斯而非完整子集，将不规则的内存访问模式转化为可预测的流式读取。客户端提出立体光栅化使双眼在位精确质量下共享大部分计算：将左右眼视锥的高斯合并为单一排序序列，共享 $\alpha$-混合 的前缀累积计算，仅在对双眼渲染结果产生差异的深度处分叉计算，实现位精确质量下的近倍渲染效率。系统架构上，云端负责 LOD 层次遍历和高斯选择，客户端负责最终光栅化和显示，通信带宽需求从视频流的约 1 Gbps 降至约 50 Mbps，实现 2.7 倍运动-光子加速。

\paragraph{子树装箱与抢先 $\alpha$-过滤 的移动端渲染}

Voyager（2025.06）~\cite{2506_02774} 面向移动设备上的城市级 3DGS 实时渲染。其关键发现是：在正常用户运动下，视锥体内新可见高斯数量大致恒定。利用这一时间相关性，Voyager 提出时间感知 LOD 搜索：将 LOD 树预先划分为高度不超过 2 的小子树，通过受约束的装箱问题将子树打包为负载均衡的工作块，每帧仅需处理与前一帧切割的差量子树而非完整遍历。在光栅化加速方面，提出抢先 $\alpha$-过滤（preemptive $\alpha$-filtering）策略：在瓦片排序阶段提前计算累积不透明度，一旦透射率饱和即丢弃后续高斯，避免了对被遮挡基元的完整 $\alpha$-混合 计算。该方法在移动端 NVIDIA Ampere GPU 上将 SmallCity 场景从 12.2 FPS 提升至可用帧率，相比现有方案实现最多 6.6 倍加速和 85\% 能量节省。

\subsection{核外渲染与虚拟内存}
\label{ss:14_6}

本节介绍虚拟内存与核外渲染技术，通过 SSD-CPU-GPU 层次存储和智能数据调度，在消费级 GPU 上实现超大规模场景的实时渲染。

\subsubsection{虚拟内存与核外渲染}

本路线介绍虚拟内存与核外渲染技术，通过按需加载和智能缓存机制，在有限 GPU 显存上实时渲染超大规模场景。

想象一下：你有一座城市的三维模型，包含数十亿个高斯点，但你的 GPU 只有 24 GB 显存。就像操作系统用虚拟内存让程序使用比物理内存更大的地址空间一样，3DGS 的核外技术将全部高斯存储在 CPU 内存或 SSD 中，每次只将当前视点需要的一小部分传输到 GPU 进行渲染。关键挑战在于：如何快速判断哪些高斯是当前需要的、如何在数据传输和 GPU 计算之间实现流水线重叠以隐藏延迟。

\paragraph{顺序点树与 GPU 缓存的无分割核外训练}

"A LoD of Gaussians"（2025.07）~\cite{2507_01110} 提出在单张消费级 GPU 上训练和渲染超大规模高斯场景——无需分割。全场景存储在核外（如 CPU 内存），直接训练 LOD 表示，仅动态流式传输相关高斯。混合数据结构结合高斯层次与顺序点树（Sequential Point Trees, SPT）实现高效的视角依赖 LOD 选择：SPT 的切割条件利用共享的根-相机距离进行全并行评估，仅存储按最大距离排序的有序对，显著降低内存开销。致密化在 CPU 端执行，通过添加新叶节点和重生低不透明度叶节点来扩展层次结构，更新后重新转换为分层 SPT（HSPT）并传输到 GPU。轻量级 GPU 缓存追踪当前驻留的 SPT 子树及其细节层级，仅当请求的子树不在缓存中或细节层级不足时才触发 RAM-VRAM（Video RAM）传输，利用帧间时间一致性将缓存命中率维持在高位。该方法从根本上避免了分块边界伪影（chunk bleeding/ghosting）和冗余重叠存储，支持从航拍全景到地面细节的无缝多尺度重建和交互式可视化。

\paragraph{虚拟页表与按需流式传输的 3DGS 渲染}

Haberl et al.（2025.06）~\cite{2506_19415} 将成熟的虚拟内存和虚拟纹理技术应用于 3DGS 渲染。通过高效识别可见高斯并动态流式传输到 GPU 实现实时渲染，仅选择必要高斯用于存储和渲染，减少内存使用并有效加速复杂场景的渲染。具体地，该方法将高斯按空间位置组织为页（pages），维护一个页表映射三维空间到存储位置，渲染时通过视锥体剔除快速确定候选页集合，仅将可见页按需传输到 GPU。虚拟纹理机制进一步优化了存储——远处高斯使用低精度表示（降低球谐阶数和参数精度），近处高斯保持全精度，实现类似 Mipmap 的存储分层。进一步展示了如何将 LOD 集成到虚拟内存方法中以加速大规模场景渲染：结合 LOD 层次选择和按需流式传输，远处场景使用低细节页占用更少显存带宽，整体数据传输量随场景复杂度的增长被有效抑制。实验在桌面和移动设备上均验证了该方法对复杂场景渲染的显著加速效果。

\paragraph{共享 MLP 学习视角依赖可见性函数}

NVGS（2025.11）~\cite{2511_19202} 解决高斯基元半透明性导致传统遮挡剔除无法应用的问题。通过小型共享 MLP 学习训练模型中所有高斯的视角依赖可见性函数，在渲染前查询 MLP 丢弃被遮挡基元。其核心观察是标准体渲染隐式编码了软遮挡——一旦透射率饱和，后续高斯对像素颜色的贡献趋近于零，可以安全丢弃。可见性提取阶段从预训练 3DGS 资产中采样 2000 个方向的虚拟相机，用 Fibonacci 球面采样确保无偏覆盖，对每个采样视角计算逐高斯可见性标签后训练一个轻量 MLP 预测"给定视角方向和距离，该高斯是否可见"。在渲染管线中，利用 Tensor Core 高效执行 MLP 推理，将神经查询直接集成到实例化软件光栅化器中——实例化机制使同一资产的多份副本共享单个 MLP 参数集，场景中超过一亿高斯的组合场景只需评估一次可见性。该方法与现有 LOD 技术互补，在 VRAM 使用和图像质量上超越当前最优。

\paragraph{多尺度金字塔监督与物理尺寸下界正则化}

PrismGS（2025.10）~\cite{2510_07830} 针对大规模城市环境中的严重锯齿伪影和优化不稳定问题，提出物理接地正则化框架。金字塔多尺度监督（MSS）通过预滤波图像金字塔强制渲染一致性：在每次训练迭代中，对渲染图像和真值图像分别构建 $L$ 层图像金字塔，真值金字塔每层先施加高斯低通滤波 $F_\sigma$ 再下采样 $D_2$（即 $I_{gt}^{(l+1)} = D_2(F_\sigma(I_{gt}^{(l)}))$），渲染金字塔直接由 3DGS 光栅化器在各分辨率下生成，逐层计算光度损失迫使模型学习内在抗锯齿表示。显式尺寸正则化为 3D 高斯施加物理下界 $\min(s_i) \geq s_{\min} = \kappa \cdot d_{\min} / f$（$\kappa$ 为超参数，$d_{\min}$ 为最近观测距离，$f$ 为焦距），从 Nyquist 采样定理出发确保高斯尺寸不小于其在最近距离处对应的像素投影大小，防止退化视角依赖基元的形成。该正则化使几何表面更加稳定可信，显著减少锯齿边缘。PrismGS 即插即用且兼容现有分治管线，在 MatrixCity、Mill-19 和 UrbanScene3D 上验证，4K 渲染下比 CityGaussian 提升约 1.5 dB PSNR。

\begin{takeaway}
分层与 LOD 表示是大规模场景实时渲染的关键：通过构建高斯基元的层次结构并根据视距动态选择精度，配合神经可见性预测和物理尺寸正则化，在保持视觉质量的同时将内存和计算开销降至可控范围。
\end{takeaway}

% ======================================================================


\section{分布式与并行训练}
\label{sec:14_4}

当单张 GPU 无法容纳大规模场景的全部参数、梯度和优化器状态时，分布式训练成为必然选择。3DGS 的分布式训练面临独特挑战：每个高斯基元只影响渲染图像中的少量像素（稀疏性），且影响范围是动态变化的。这使得传统的数据并行和模型并行策略难以直接套用。

\subsection{并行训练范式}
\label{ss:14_7}

本节介绍模型并行与数据并行策略，将数千万乃至数十亿个高斯基元分配到多张 GPU 上进行分布式训练。

\subsubsection{模型并行与数据并行}

本路线介绍将大规模3DGS训练分布到多张GPU上执行的模型并行与数据并行方法。

分布式训练的核心问题是：如何将数千万乃至数十亿个高斯基元及其优化器状态分配到多张 GPU 上，同时保证训练结果与单 GPU 训练等价？本节的方法从三个不同角度回答了这一问题：Grendel 采用混合并行策略（高斯级分配 + 像素级分配），RetinaGS 通过数学证明将渲染方程分解为凸子空间上的层次合并，DOGS 用 ADMM 共识约束将优化问题分布式化，BlitzGS 在 Grendel 基础上引入三层工作负载管理实现数量级加速。

\paragraph{高斯-像素混合并行与稀疏全到全通信}

\textbf{背景与动机。}标准 3DGS 训练被限制在单张 GPU 上，造成严重的内存和计算瓶颈。以 Rubble 数据集为例（1657 张 4K 分辨率图像），单张 A100 40GB GPU 最多容纳约 1120 万个高斯——远低于质量饱和点。Grendel（2024.06）~\cite{2406_18533} 的实验表明，随着高斯数量增加，重建质量持续提升。因此，突破单 GPU 内存限制是提升大规模场景重建质量的关键。

然而，将分布式训练技术直接应用于 3DGS 面临根本性困难：3DGS 不是神经网络，其计算管线具有动态且不均衡的工作负载模式。不同阶段的并行化维度不同（像素级 vs.\ 高斯级），不同像素的渲染计算量差异巨大，且随训练进度持续变化。现有 DNN 分布式框架（Fully Sharded Data Parallel, FSDP、DeepSpeed）假设一致均衡的工作负载和规则密集张量操作，无法直接适配 3DGS 的不规则稀疏计算模式。

\textbf{核心思想：混合并行 + 空间局部性。}Grendel 的核心设计洞察是 3DGS 训练管线天然具备混合并行性（mixed parallelism）——高斯变换阶段适合高斯级并行（每个 GPU 操作不相交的高斯子集），图像渲染和损失计算阶段适合像素级并行（每个 GPU 操作不相交的像素子集）。3DGS 的空间局部性特征（每个像素仅受少量高斯影响）使得阶段间的数据交换可以通过稀疏通信高效完成。这两个特性共同构成了 Grendel 分布式训练的设计基础。

\textbf{技术细节。}Grendel 的系统设计包含三个关键机制。

（1）\textit{混合并行训练}。训练管线分为两个阶段分配工作。高斯级分配（Gaussian-wise distribution）：将高斯参数（包括位置、协方差、球谐系数）和优化器状态均匀分割到各 GPU，每张 GPU 独立计算所分配高斯的变换（投影、深度、覆盖半径）。由于单个高斯的计算量差异不大，均匀分配即可实现负载均衡并最大化高斯容纳数量。像素级分配（pixel-wise distribution）：将每张图像划分为 $16 \times 16$ 像素块，将连续像素块序列自适应分配给不同 GPU，用于渲染和损失计算。分配连续区域可以最大化空间局部性，减少跨 GPU 高斯传输量。批化训练时，每张 GPU 可从批次中不同图像获取像素块。

（2）\textit{稀疏全到全通信}。高斯变换完成后需要一次通信步骤——GPU 需要访问与其像素分区相交的高斯，但哪些高斯与哪些像素相交是视角依赖的，无法预先确定。Grendel 利用空间局部性优化通信：每个 GPU 先确定其像素分区所需的高斯集合，然后通过稀疏全到全通信（sparse all-to-all）仅传输必要的高斯。反向传播时执行反向的稀疏全到全通信。这一设计与 FSDP 的关键区别在于：Grendel 的高斯级分配不仅为存储也为计算服务（不像 FSDP 的权重分片仅为存储），且使用稀疏通信而非 FSDP 的密集全收集。

（3）\textit{迭代式负载重均衡}。由于物理场景全局稀疏（天空 vs.\ 树木区域密度差异巨大），渲染计算量在空间和时间上极不均匀。Grendel 采用两种自适应策略。像素级重均衡：在每个训练周期后记录每个像素块的渲染时间，利用训练过程中场景变化缓慢的特性，用上一周期的计时作为当前周期的负载估计，自适应重新分配像素块使各 GPU 工作量近似均衡。高斯级重均衡：每次致密化（densification）后，不同高斯的克隆和分裂速率差异导致分布不均，定期重新均匀分配高斯以恢复均衡。

\textbf{超参扩展策略。}批化训练需要调整 Adam 优化器的超参数。Grendel 提出"独立梯度假设"——各视角的梯度相互独立，推导出学习率和动量的缩放规则：$\lambda' = \sqrt{|\text{batch}|} \cdot \lambda$（$\sqrt{\text{batch size}}$ 缩放规则），以及一阶和二阶矩的相应调整 $\beta_1', \beta_2'$。直觉上，如果各视角梯度独立，则 $n$ 个视角的平均梯度方差是单视角的 $1/n$，因此学习率需要放大 $\sqrt{n}$ 倍以匹配更新步长。这一规则简单有效，无需繁琐的超参搜索即可在大批量训练中保持收敛特性。

\textbf{实验验证。}Grendel 在 Perlmutter GPU 集群上评测（4 A100 40GB/GPU 节点，NVLink 25GB/s，Slingshot 200Gbps 网络）。核心结果：（1）\textit{吞吐量扩展}——在 Rubble 场景（35M 高斯）上，4 GPU 批量 1 达 5.55 images/s，32 GPU 批量 64 达 38.03 images/s，近线性扩展；在小规模 Train 场景上，16 GPU 批量 32 仅用 2 分 43 秒完成 30K 图像训练；（2）\textit{内存扩展}——高斯容量随 GPU 数量线性增长，16 GPU 批量 16 在 Rubble 上支持超过 4000 万高斯；（3）\textit{质量随高斯数量提升}——在 Rubble、MatrixCity Block\_All 和 Bicycle 三个场景上，PSNR、SSIM 和 LPIPS 均随高斯数量增加而持续改善，其中 Rubble 从单 GPU 的 11.2M 高斯（PSNR 26.28）提升至 16 GPU 的 40.4M 高斯（PSNR 27.28）；（4）\textit{与 CityGaussian 对比}——在降采样场景上，Grendel 在相同 200K 图像训练量下达到或超过 CityGaussian 官方脚本的 PSNR，且训练速度快 3--6.7 倍，使用更简单（无需分治管线的多阶段超参调优）。

\textbf{局限与意义。}Grendel 的局限性在于：通信开销随 GPU 数量增加而上升（稀疏全到全通信的网络延迟），在超大规模集群上的扩展效率仍有提升空间；需要额外的 GPU 间通信带宽（NVLink 或高速网络）。尽管如此，Grendel 作为首个系统级分布式 3DGS 训练框架，其混合并行设计、稀疏通信策略和超参缩放规则为后续分布式 3DGS 工作（RetinaGS、DOGS、BlitzGS 等）奠定了系统设计范式，与分治方法（CityGaussian）形成互补——CityGaussian 的粗训练阶段即可利用 Grendel 进行多 GPU 加速。

\paragraph{渲染方程凸子空间分解与 KD-Tree 分布式训练}

\textbf{基线做法及其局限。}在分布式 3DGS 训练中，最直接的模型并行策略是让每张 GPU 负责完整图像的部分渲染——每张 GPU 持有高斯的一个分区，独立渲染后合并各 GPU 的输出。与 Grendel 通过混合并行（高斯级分配 + 像素级分配）和稀疏通信分发训练不同，这种纯模型并行策略设计更简洁，但存在根本性冗余计算问题：许多 GPU 在射线的不透明度已经饱和的深度之后仍在继续渲染像素，浪费大量算力。

更严重的是，此前的大规模 3DGS 方法（如 VastGaussian、CityGaussian）将高斯划分为互不相交的子集，各子集独立训练为独立的 3DGS 模型。这实际上偏离了原始渲染方程——射线穿越子空间边界时，$\alpha$-混合 的累积过程被人为截断。虽然这些方法通过引入额外的邻近相机和高斯来缓解边界伪影，但伪影在特定视角下仍会加剧并突破缓解策略的防护。这些方法因此主要假设鸟瞰城市场景的简单视角和高斯分布。Grendel 的混合并行策略虽然避免了边界问题，但其高斯级分配仍需处理大量跨 GPU 的高斯依赖关系，通信开销随 GPU 数量增加而上升。

\textbf{核心改造：渲染方程的凸子空间分解。}RetinaGS（2024.06）~\cite{2406_11836} 的核心理论贡献是严格证明了 3DGS 的 $\alpha$-混合 渲染方程可以被分解为凸子空间上的层次合并形式，且数学上与单 GPU 渲染完全等价。具体地，原始渲染方程为
\begin{equation}
\label{eq:retinags-original}
C(l) = \sum_{g_i \in \mathcal{N}_l} c_i \alpha_i \prod_{g_j \in \mathcal{N}_{l,<i}}(1-\alpha_j)
\end{equation}

其中 $\mathcal{N}_l$ 为对射线 $l$ 有贡献的高斯集合（按深度排序），$\alpha_i = \sigma_i \cdot g_i(l)$ 为有效不透明度。RetinaGS 将三维空间划分为 $K$ 个凸子空间 $S_1, S_2, \ldots, S_K$（子空间无缝覆盖整个空间，且每个子空间为凸集，即所有分界面为平面），为每个子空间构造一个允许重叠的高斯子集：$N_k = \{i : \mathbf{n}_{km}^\top u_i + d_{km} > -D_i, \forall m\}$，即子集 $N_k$ 包含所有其三维椭球与子空间 $S_k$ 相交的高斯基元（$u_i$ 为基元中心，$D_i$ 为截断阈值，取为高斯椭球主轴长度的三倍）。这一"重叠覆盖"设计至关重要——它确保边界附近的高斯同时属于相邻子集，从而保证渲染等价性。在每个子空间上独立计算部分颜色和部分不透明度：

\begin{align}
C_k(l) &= \sum_{g_i \in \mathcal{N}_{lk}} c_i \alpha_i \prod_{g_j \in \mathcal{N}_{lk,<i}}(1-\alpha_j) \label{eq:retinags-partial-color} \\
T_k(l) &= \prod_{g_i \in \mathcal{N}_{lk}} (1-\alpha_i) \label{eq:retinags-partial-trans}
\end{align}

其中 $\mathcal{N}_{lk} \subseteq N_k$ 为子空间 $S_k$ 内真正与射线 $l$ 相交的高斯子集（通过指示函数 $1(g_i \in \mathcal{N}_{lk}) = \prod_m 1(\mathbf{n}_{km}^\top x_i + d_{km} > 0)$ 在 GPU 核函数中高效判断，$x_i$ 为射线与 2D 高斯的交点）。然后按射线穿越子空间的顺序 $o_l$（排列置换）合并：

\begin{equation}
\label{eq:retinags-merge}
C(l) = \sum_{k \in o_l} C_k(l) \prod_{m \in o_{l,<k}} T_m(l)
\end{equation}

RetinaGS 在附录中给出了完整的等价性证明：将部分颜色和部分不透明度代入合并公式后，利用乘法分配律展开嵌套求和，再利用乘法交换律将各子空间的透射率项重新排列，最终可归约为单一求和 $\sum_{g_i \in \mathcal{N}_l} c_i \alpha_i \prod_{g_j \in \mathcal{N}_{l,<i}}(1-\alpha_j)$，与原始渲染方程完全一致。直观理解：就像一幅画可以按区域分别绘制后无缝拼接一样，射线穿越空间的渲染过程可以按射线经过的子空间顺序逐段累积颜色，每段独立计算"这段贡献了多少颜色"和"这段透过了多少光线"，最后按顺序合并即为完整渲染结果。

\textbf{为什么渲染方程可分解：凸性要求与设计动机。}分解的关键约束在于子空间必须是凸的。凸性保证了射线与子空间的交集是连续的线段（而非多个离散片段），从而使射线穿越子空间的顺序 $o_l$ 是良定义的——每条射线以确定的先后顺序穿越各子空间，不会在同两个子空间之间反复跳跃。这一性质是合并公式~\eqref{eq:retinags-merge} 成立的前提：如果子空间非凸，射线可能在同一子空间内有多个不连续的交集段，破坏按子空间独立渲染再顺序合并的逻辑。

分界面必须是平面（三维凸集的边界性质），这使得指示函数的计算极其高效——仅需检查交点到各分界面的带符号距离，在 GPU 核函数中只需几次乘加操作。这一理论等价性使 RetinaGS 可应用于任意场景和任意高斯分布，不假设鸟瞰视角或均匀分布，从根本上区别于 VastGaussian/CityGaussian 的互斥子集策略。

\textbf{实现架构：KD-Tree 分区与流水线。}在空间分区策略上，RetinaGS 采用三维 KD-Tree（深度 $L$，产生 $K = 2^L$ 个子空间）递归地以垂直于 X、Y、Z 轴的超平面平分空间并均分基元，使各子集基元数大致均衡，优化各 worker 的负载均衡和峰值内存。整个训练流水线为：（1）中心管理器维护 KD-Tree 和子空间定义，解析渲染请求并分发到相关子模型；（2）前向传播中，各 GPU 并行计算其子模型的部分颜色 $C_k$ 和部分不透明度 $T_k$（以与目标图像同分辨率的 2D map 表示，通信量极小）；（3）中心管理器按公式~\eqref{eq:retinags-merge} 合并完成渲染、计算损失，并将梯度分发回各子模型进行反向传播。每隔若干训练轮次重新执行 KD-Tree 分区以适应基元中心的显著漂移。属于多个子集交集区域的基元需要梯度同步，但实验发现省略此步对重建质量影响可忽略。对比实验表明，KD-Tree 分区相比固定大小网格分区在峰值内存和训练速度上均有显著优势。

\textbf{初始化策略：MVS 替代启发式致密化。}标准 3DGS 的初始化依赖 COLMAP 稀疏点云，训练过程中通过启发式致密化（分裂和克隆）控制高斯数量。但该策略在高分辨率训练中存在问题：过多的致密化迭代导致结果恶化，过度分裂使训练不稳定，且高斯数量难以精确控制。RetinaGS 放弃启发式致密化，改用多视角立体匹配（Multi-View Stereo, MVS）对所有训练数据进行密集深度估计，生成稠密三维点云后灵活采样为高斯基元，训练过程中固定基元数量不再调整。这种设计使基元数量完全可控可扩展，配合分布式训练方法可轻松扩展到十亿级别。由于 MVS 提供了较准确的位置参数，位置参数的学习率从原始 3DGS 的 $1.6 \times 10^{-4}$ 降至 $1.6 \times 10^{-6}$（百分之一），采用指数衰减函数。

\textbf{实验规模与关键结论。}RetinaGS 在 Mip-NeRF 360、ScanNet++（8408 像素宽的高分辨率图像）、Mega-NeRF 和 MatrixCity 上验证。最令人瞩目的实验是在完整 MatrixCity 数据集（MatrixCity-ALL：141{,}652 张图像、超过 1480 亿像素）上用 64 张 A100 GPU 训练 10 天，首次成功训练超过十亿个高斯基元的 3DGS 模型。系统实验揭示了两个关键规律：（1）在高斯数量相当时，RetinaGS 的 PSNR 与单 GPU 3DGS 训练结果高度一致，验证了等价性的实际有效性；（2）视觉质量随基元数量增加呈现清晰的上升趋势——固定分辨率下增加基元数可提升质量，固定基元数下提高分辨率亦可提升质量，两者协同效果更佳。这证实了基元数量决定了 3DGS 模型的容量上限，而高分辨率带来更多信息量需要更大模型容量才能充分拟合，因此视网膜级重建质量需要同时保证高分辨率和大规模模型。RetinaGS 的渲染方程等价性验证实验也值得关注：将相机光轴精确放置在两个子空间的分界平面上时，每个子模型的渲染结果在分界面处呈现清晰的明暗分界，合并后与单 GPU 渲染完全一致，而基于互斥子集的方法在此测试中失败。

\paragraph{ADMM 共识约束与自适应惩罚的分布式训练}

与 Grendel 保持原始 3DGS 算法不变、通过系统级并行化加速不同，DOGS（2024.05）~\cite{2405_13943} 从优化算法层面修改训练过程——将交替方向乘子法（Alternating Direction Method of Multipliers, ADMM）引入 3DGS 训练。场景被递归分解为 $K$ 个有重叠的空间块（每次沿最长轴二分并扩展边界框构造重叠区域），主节点维护全局 3DGS 模型 $z$，$K$ 个从节点维护各自的局部模型 $x_i^k$。每个从节点的训练目标被修改为原渲染损失加正则项：$\min_{x_i^k} L(x_i^k) + \frac{\rho}{2}\|x_i^k - z^t + u_i^k\|_2^2$，其中 $u_i^k$ 是对偶变量。正则项的作用是"拉向"全局共识——每个块的局部高斯不能偏离全局模型太远，偏离程度由惩罚参数 $\rho$ 控制。共识步骤通过平均所有从节点的局部高斯更新全局模型 $z^{t+1} = \frac{1}{K}\sum_{k=1}^{K}(x_i^k)^{t+1}$，再将对偶变量推近 $u_i^{t+1} = u_i^t + (x_i^{t+1} - z^{t+1})$。DOGS 引入自适应惩罚参数（根据原始残差和对偶残差的比例动态调整 $\rho$）和过松弛策略（$\alpha^t \in [1.5, 1.8]$）加速 ADMM 收敛，在 2000 次迭代后固定惩罚参数以保证收敛性。从节点训练完成后被丢弃，推理时仅查询全局模型。

值得注意的是，DOGS 的递归分割策略沿场景边界框的最长轴二分、重估计每个子块的紧凑边界框后再扩展以构造重叠区域，确保相邻子块在重叠带内共享高斯，从而在 ADMM 共识步骤中对这些共享高斯施加显式一致性约束。实验还揭示了分块数量的一个有趣规律：8 块配置（8 个从节点并行）在大规模场景上不仅训练速度最快，渲染质量也往往优于 4 块或 2 块配置，这是因为更细粒度的空间分解使各子块的局部优化问题更简单、收敛更稳定。DOGS 被接收至 NeurIPS 2024，训练速度提升 6 倍以上，同时达到最优渲染质量。

\paragraph{索引奇偶性分片与计划重要性评分的闪电重建}

BlitzGS（2026.05）~\cite{2605_13794} 在 Grendel 的双层并行基础上进一步引入三层工作负载管理，其核心原则是"城市级优化应避免触碰当前设备、当前视角或当前训练阶段所不需要的高斯"。系统层通过索引奇偶性（index parity）而非空间块分片高斯——每张 GPU 存储约 $|G|/M$ 个高斯，渲染时各 GPU 投影本地高斯后通过单次全到全通信将投影高斯路由到所属 tile 的 GPU 上完成光栅化，从根本上缓解了空间分区导致的跨块可见性冗余问题。模型层通过计划重要性评分（scheduled importance scoring）主动缩减全局高斯数量：在安排的简化扫描中，对每个高斯计算跨视角的归一化贡献密度 $s_i = \sum_{v: a_{i,v} > 0} w_{i,v} / (a_{i,v} + \epsilon)$，分别执行随机重要性加权采样和确定性累积质量截断（保留评分总和 99\% 的前缀子集），并将重要性信号反馈到致密化梯度重加权和视角级过滤中。视角层通过距离 LOD 门控 $L_v(d, t) = \text{clamp}(\lfloor\log_{f}(d_0/d)\rfloor, 0, L_{\max}(t))$ 和基于每视角冗余掩码的贡献性剔除，从几何和重要性两个独立维度裁剪每帧活跃集合。BlitzGS 在城市级场景上实现数量级加速，训练仅需数十分钟。具体而言，索引奇偶性分片使渲染步骤仅需一次跨 GPU 全到全通信（all-to-all exchange）即可完成投影高斯到 tile 所有者的路由，避免了 Grendel 中多轮稀疏全到全通信的延迟累积。在模型层，计划重要性评分与训练进度耦合——早期训练阶段执行较宽松的截断以保留足够的表达容量，后期逐步收紧以移除冗余高斯，确保在数量级加速的同时不损失渲染质量。

\paragraph{混合捕获梯度分析}

第~\ref{sec:14_1}~节已从跨视角融合角度讨论了 Cho（2026.05）~\cite{2605_00052} 对混合捕获场景的方差分解分析。从分布式训练的梯度角度，该工作的核心结论是：双视角累积（two-view accumulation）本身的方差减半效应是主导因素，结构性和随机性配对在期望上方差等价。

\subsection{联邦与多智能体协同}
\label{ss:14_8}

本节介绍联邦学习与多智能体重建技术，在保护数据隐私和不共享原始图像的前提下实现多智能体协同构建全局模型。

\subsubsection{联邦学习与多智能体重建}

本路线介绍在通信带宽有限且不允许传输原始图像的约束下，通过联邦学习实现多智能体协同 3DGS 重建的方法。

在分布式机器人场景中，多个智能体可能各自采集局部数据，需要在不共享原始数据的前提下协同构建全局模型。与上节的分布式训练不同，联邦学习的约束是通信带宽有限且不允许传输原始图像。每个智能体只能在本地训练后共享模型参数（高斯属性），由中心服务器聚合为全局模型。核心挑战是部分可观测性：每个智能体只能看到场景的一小部分，聚合时需要解决不同智能体对同一区域观测不一致的问题。

\paragraph{可见性感知聚合的联邦 3DGS 框架}

F3DGS（2026.04）~\cite{2604_01605} 是首个联邦 3DGS 框架，用于去中心化多智能体 3D 重建。首先通过注册多个客户端的局部 LiDAR 点云构建共享几何支架初始化全局模型——具体做法是利用 LiDAR SLAM（如 KISS-SLAM）构建全局锚点轨迹，再通过 Umeyama $\mathcal{S}im(3)$ 估计将各客户端的单目轨迹对齐到全局坐标系，并在拼接边界施加指数衰减的 SE(3) 平滑修正以消除位姿不连续性。联邦优化中固定高斯位置以保持几何对齐，每个客户端仅更新外观属性（尺度、旋转、不透明度、球谐系数），从结构上杜绝了跨客户端的几何漂移。每个客户端在本地训练中维护逐高斯可见性计数器 $v_{k,i}$，记录高斯 $i$ 在客户端 $k$ 的光栅化过程中具有正投影半径的次数。服务器使用可见性感知聚合——聚合权重为 $\alpha_{k,i} = v_{k,i} / (\sum_j v_{j,i} + \epsilon)$，即根据每个客户端观测每个高斯的频率加权——解决多智能体探索中的部分可观测性挑战。对于总可见性为零的高斯，保留上一轮全局值不做更新；四元数参数因符号歧义性需在加权平均前对齐符号并归一化。实验在自采集的多序列室内数据集（同步 LiDAR、RGB 和 IMU）上验证，联邦聚合模型的 PSNR 与集中式训练差距在 2 dB 以内。

\paragraph{云-边-端三层架构与自适应区域规划}

Radiant（2024.12）~\cite{2412_05546} 考虑系统异构性，基于环境信息为每个边缘设备分配合适的区域和工作负载。Radiant 首次将云-边-端（cloud-edge-device）三层层次化架构引入大规模 3DGS 重建：设备层负责本地图像采集和独立训练，边缘层对辖区内设备模型进行初步聚合，云端层执行最终融合以生成全局模型——该架构将计算负载从云端下沉到边缘，同时保护设备端的位置隐私和数据隐私。核心算法包括自适应区域规划（ARP）和资源感知任务分配（RTP）：ARP 根据各边缘的计算能力和带宽分配摄像机位姿数量和位置，RTP 进一步将各区域内的采集任务按设备的算力异构性分配到具体设备，以最小化端到端延迟并缓解木桶效应。在模型聚合方面，Radiant 发现传统的模型合并（直接切割边界后拼接）会在分界面处产生明显的视觉不连续，且问题随合并模型数量增加而加剧。为此，Radiant 提出基于少量地面真值图像的微调重训练策略，在聚合后对全局模型进行少量轮次的精调以消除边界伪影，同时保持隐私。实验在两个真实大规模场景数据集上验证，重建质量提升最多 25.7\%，端到端延迟降低最多 79.6\%。

\paragraph{局部模型压缩与自知识蒸馏的协作表面重建}

CoSurfGS（2024.12）~\cite{2412_17612} 提出多智能体协作的 3DGS 表面重建框架，采用与 Radiant 相似的设备-边缘-云三层分布式架构，但核心目标从渲染质量转向高质量表面几何恢复。CoSurfGS 的两个关键组件是局部模型压缩（LMC）和模型聚合方案（MAS）。LMC 在设备模型上传前执行剪枝：定义优先级评分筛选缺乏多视角一致性且不透明度低的高斯点，移除局部模型在重叠区域的冗余表示，从而大幅降低传输和聚合的内存开销。MAS 采用自知识蒸馏机制渐进式地将局部表面结构聚合为全局模型——在聚合过程中，先用单视角几何约束（法向量邻域一致性）和多视角几何约束（基于单应矩阵的重投影误差与光度一致性）训练每个局部模型，然后通过蒸馏保持各局部区域的关键几何结构并对齐相邻共视区域的表面。实验在 Urban3D、MegaNeRF 和 BlendedMVS 数据集上验证，在表面重建和新视角合成两个维度均达到 SOTA，同时训练时间和内存消耗显著低于所有对比方法。

\paragraph{分布式科学可视化}

Han~et~al.（2025.09）~\cite{2509_12138}~将分布式~3DGS~扩展到高性能计算（High-Performance Computing, HPC）科学可视化领域。该工作的完整流水线为：先用~ParaView~从体数据中提取等值面点云作为高斯基元初始位置，生成结构化轨道式合成相机视角集（所有节点使用相同设置保证训练一致性），然后将数据按节点分区并添加幽灵单元（ghost cells），各节点独立渲染训练图像和背景掩码后并行训练高斯模型，最终合并所有节点的~splats~进行全局渲染。幽灵单元通过在分区边界向邻域延伸一层额外的数据来消除缝隙和白色条纹伪影，背景掩码则过滤掉对当前分区无关的背景高斯。实验在~Argonne~的~Polaris~超算上进行，使用三种不同规模的科学数据集（Kingsnake~约~4M~点、Rayleigh-Taylor~约~18.2M~点、Richtmyer-Meshkov~约~106.7M~点），每节点配备~4~块~NVIDIA~A100~GPU。单节点内多~GPU~在~Kingsnake~数据集~$2048^2$~分辨率上实现~5.6~倍加速；多节点实验中，Richtmyer-Meshkov~数据集在~8~节点上训练~$2048^2$~图像实现~3.1~倍加速（相对~4~节点），且~PSNR~和~SSIM~保持稳定（PSNR~30.04，SSIM~0.97）。该工作发表于~SC 2025 Research Poster，为未来原位（in situ）科学可视化奠定了基础。

\begin{takeaway}
分布式与并行训练采用设备-边缘-云三层架构和自知识蒸馏聚合策略，将大规模 3DGS 重建从单机扩展到多智能体协作，已在 HPC 科学可视化等超大规模场景中验证了可扩展性。
\end{takeaway}

% ======================================================================

\subsection{近期方法脉络：分布式训练、多GPU与大规模训练}

这一组近期工作延续 3DGS training、Gaussian splatting，主要把改动集中在规模扩展等环节。下面只保留相对既有路线的变化，便于和前文方法直接对照。

\paragraph{Splaxel}
Splaxel（2026.06）~\cite{2606_18588} 延续 3DGS training 在“分布式训练”上的基本设定，主要改变是：Splaxel：基于像素级局部渲染和全局合成的分布式3DGS训练框架。3DGS 的显式高斯表示和可微渲染能力保持为基础。
其设计动机是：提出通信高效的分布式训练方案，避免区域分割导致的全局不一致和高斯级全局交换导致的通信爆炸。
具体而言，每个GPU执行局部渲染并在像素级合成全局图像/梯度，减少跨GPU传递高斯参数的需求。
适用边界是：全局合成策略依赖网络和GPU拓扑；极复杂遮挡可能增加同步需求。

\paragraph{A Scalable PyTorch Abstraction for Multi-GPU Gaussian Splatting}
A Scalable PyTorch Abstraction for Multi-GPU Gaussian Splatting（2026.06）~\cite{2606_11390} 延续 Gaussian splatting 在“多GPU与大规模训练”上的基本设定，主要改变是：多GPU PyTorch backend，通过CUDA unified memory和NVLink在operator层分布高斯参数和splatting算子。3DGS 的显式高斯表示和可微渲染能力保持为基础。
其设计动机是：提出把多GPU暴露为聚合PyTorch设备的后端，在operator层分布高斯参数和渲染操作，支持城市级重建和高分辨率训练。
具体而言，系统利用CUDA unified memory和NVLink管理跨GPU数据访问，模型代码仍像单设备运行，底层splatting operator自动分发计算和参数。
适用边界是：依赖NVLink/统一内存硬件条件；跨设备访问模式可能影响性能稳定性。
\section{其他大规模方法}
\label{sec:14_5}

除上述四个主要方向外，还有多种技术从不同角度推进了大规模 3DGS 的发展。

\subsection{核外训练与内存管理}
\label{ss:14_9}

本节介绍核外训练与内存优化技术，通过 SSD-CPU-GPU 层次存储和智能数据调度，在消费级 GPU 上训练十亿级高斯模型。

\subsubsection{核外训练与内存优化}

本路线介绍利用训练迭代的稀疏性，通过SSD-CPU-GPU层次存储实现消费级GPU上训练十亿级高斯的方法。

核外训练的核心观察是：3DGS 训练的每次迭代只涉及当前相机批次可见的一小部分高斯，其余高斯在本次迭代中完全不参与计算。这意味着 GPU 内存不需要同时容纳所有高斯，而只需作为"工作集缓存"保存当前迭代所需的高斯子集即可。本节的方法利用这一稀疏性，通过 SSD-CPU-GPU 的层次存储和智能数据调度，在消费级 GPU 上训练十亿级高斯模型。

\paragraph{SSD-CPU-GPU 三级存储与轨迹自适应差分流式传输}

TideGS（2026.05）~\cite{2605_20150} 是首个将 3DGS 训练扩展到十亿高斯量级的单 GPU 核外框架，通过将 SSD--CPU--GPU 构成三级存储层次，使消费级硬件上的超大规模训练成为可能。

\textbf{基线做法。}标准 3DGS 要求所有训练状态（参数、梯度、优化器动量）常驻 GPU 显存，单张 24\,GB GPU 仅能容纳约 1100 万高斯。已有的主机卸载方案（CLM、GS-Scale）将部分状态移至 CPU 内存，但仍需在 GPU 上保留关键几何信息或光栅化缓冲区，实际扩展上限约一亿高斯。超过此阈值后，全局基数排序等光栅化辅助缓冲区本身就会耗尽显存（$>$22\,GB）。多 GPU 分布式方案（Grendel-GS、CityGaussian）虽可聚合显存，但需要昂贵的硬件集群和复杂的工程实现，且引入分治边界处理问题。

\textbf{改了什么。}TideGS 将 GPU 显存从"持久参数存储"重新定义为"高带宽工作集缓存"——每次迭代仅在 VRAM 中驻留当前相机批次所需的高斯子集，完整参数表始终保存在 SSD 上。这一范式转换将显存需求从 $O(N)$ 解耦为 $O(|R_t|)$（驻留块数量），使单卡可训练规模不受总高斯数限制，而仅受 SSD 容量和带宽的约束。

\textbf{实现原理。}TideGS 由三项协同技术支撑 SSD--CPU--GPU 层次管理：

\begin{enumerate}
\item \textbf{块虚拟化几何。}将所有高斯参数打包为 SSD 对齐的空间块（$B=4096$ 个高斯/块，约 944\,KiB/块），按 Morton 曲线排序以保证空间邻近高斯映射到相邻块。采用两级可见性过滤：CPU 端对块包围球执行保守视锥体剔除，确定候选工作集 $K_t$；GPU 端在驻留块内执行精确光栅化，确定最终活跃高斯集 $I_t$。块分配在初始化后固定不变，高斯位置更新仅刷新块包围盒但不迁移跨块，因此保持标准 3DGS 渲染语义不变。

\item \textbf{层次化异步管线。}训练循环分为四个阶段：识别工作集 $\to$ 预取并驻留 $\to$ GPU 渲染与反向传播 $\to$ 淘汰并写回。SSD 读取、主机--设备传输、GPU 计算、SSD 写回在专用 I/O 线程和独立 CUDA 流上并行执行。SSD 存储采用追加式日志结构（log-structured），避免原地随机写，更新块顺序写入补丁段并通过索引指针定位最新版本。CPU DRAM 作为 LRU 热缓存，脏块仅在内存压力淘汰时异步刷入 SSD。双缓冲机制使下一迭代的传入块传输与当前迭代的 GPU 计算重叠。

\item \textbf{轨迹自适应差分流式传输（TIDE）。}利用聚类 TSP 排序的相机序列，使连续迭代访问高度重叠的块工作集。设当前驻留集为 $R_t$、下一迭代可见集为 $K_{t+1}$，通过集合差分仅传输增量：

\begin{equation}
S_t^{+} = R_{t+1} \setminus R_t, \quad S_t^{-} = R_t \setminus R_{t+1}
\end{equation}

即保留重叠部分 $R_t \cap R_{t+1}$，仅流入 $S_t^{+}$、淘汰 $S_t^{-}$。驻留评分结合下一步可见性和 LRU 近期性：$s(k) = \lambda \cdot \mathbf{1}[k \in K_{t+1}] + (1-\lambda) \cdot \mathrm{Recency}(k)$，并通过相机均衡 Top-$C$ 策略确保每个视角获得足够的块配额。
\end{enumerate}

\textbf{设计动机。}3DGS 训练本质上是稀疏的——在大规模场景中，单次迭代仅访问约 0.39\% 的高斯（MatrixCity 数据），且沿平滑相机轨迹的连续视角具有高度重叠的可见集。这种可见性稀疏性和时间局部性使得 SSD 层级的核外训练可行：跨层数据流量可随可见块工作集的增量变化缩放，而非随全模型规模 $N$ 缩放。消融实验证实，关闭差分流式传输使 PCIe 流量增加 8.5 倍（0.10 $\to$ 0.85\,GB/iter），关闭异步管线使迭代耗时增加 2.3 倍（90.7 $\to$ 210.5\,ms），去除 Morton 排序使 CPU 缓存命中率从 95.2\% 骤降至 42.1\%。

\textbf{关键参数与实验结果。}在单张 NVIDIA RTX A5000（24\,GB）上，TideGS 成功训练 11 亿高斯的 MatrixCity 场景，达到 26.1\,dB PSNR，远超 CLM 在一亿高斯上可达到的 25.0\,dB。在显存可容纳的规模上（Mip-NeRF 360），TideGS 与原生 3DGS 的 PSNR 差距仅 0.11\,dB，吞吐量开销 $<$15\%。优化器状态（Adam 动量）仅针对驻留块实例化，淘汰时丢弃、重新驻留时冷启动，以此换取更低的跨层流量和更小的显存占用。TideGS 的局限性在于对相机轨迹连续性的依赖——非结构化图像集的块复用率较低会导致更高的块淘汰率和跨层流量。

\paragraph{选择性卸载与延迟优化器更新的主机卸载训练}

GS-Scale（2025.09）~\cite{2509_15645} 将所有高斯存储在主机内存，每次前向/反向传播仅传输子集到 GPU。其核心观察是：大规模场景中每次训练迭代仅访问约 8\% 的高斯，而高斯参数、梯度和优化器状态占 GPU 总内存的约 90\%。因此，将占主体的参数和优化器状态卸载到 CPU 可大幅降低显存需求。三项系统级优化各有针对性：

（1）\emph{选择性卸载}——仅将几何属性（位置、缩放、旋转共 10 个参数，占总参数 59 个的 17\%）保留在 GPU 上用于快速视锥体剔除，其余外观属性全部存储在主机内存。

（2）\emph{参数前向（Parameter Forwarding）}——提前更新下一迭代所需的高斯参数并传输到 GPU，打破 CPU 优化器更新与 GPU 前向/反向传播之间的数据依赖，实现两者管线化并行执行。

（3）\emph{延迟优化器更新}——利用 Adam 优化器在梯度为零时的确定性衰减行为（动量和方差仅乘以 $\beta_1^{d+1}$、$\beta_2^{d+1}$），通过 4 位计数器延迟零梯度高斯的更新，仅在梯度非零或计数器饱和时恢复状态，将不必要的内存访问减少至约 6.7\%。

GPU 内存需求降低 3.3--5.6 倍，训练速度与纯 GPU 训练可比。在 RTX 4070 Mobile GPU 上，GS-Scale 可将高斯数量从 400 万扩展到 1800 万，LPIPS 改善 23--35\%。

\paragraph{属性级卸载与 TSP 调度的精确高斯缓存}

CLM（2025.11）~\cite{2511_04951} 通过将高斯卸载到 CPU 内存并在需要时加载到 GPU，在单张 RTX 4090 上渲染一亿高斯场景（MatrixCity BigCity 数据集）。CLM 的设计基于四项关于 3DGS 内存访问模式的关键观察：（1）每个视角的可见高斯集可提前计算，允许将下一迭代的加载与当前迭代的计算重叠；（2）相邻视角间存在大量高斯重叠，可缓存复用高斯以减少通信量；（3）训练过程具有空间局部性，同一区域的视角倾向于访问相同高斯。CLM 提出三项核心优化：\emph{属性级卸载（Attribute-Wise Offload）}仅将关键属性（位置、旋转、缩放共 10 个浮点数，不足总参数的 20\%）保留在 GPU 用于视锥体剔除，非关键属性存储在 CPU 按需加载；\emph{精确高斯缓存（Precise Gaussian Caching）}利用视锥体剔除已计算的可见集交集，避免连续微批次间重复传输重叠高斯；\emph{管线顺序优化}将微批次调度建模为旅行商问题（Traveling Salesman Problem, TSP），以高斯集对称差为距离度量，最大化连续微批次的访问重叠和早期 Adam 更新完成率。该方法在数学上保持与原始 3DGS 完全一致的渲染公式，已被 ASPLOS 2026 接收。

\paragraph{视角感知隐式编码与解码高斯混合的整体优化}

MixGS（2025.05）~\cite{2505_23280} 提出整体优化框架——不分割场景，而是将相机位姿和高斯属性集成为视角感知表示，解码为精细高斯。具体流程包括四个步骤：首先根据相机位姿提取视锥体内的高斯基元，利用不透明度信息进一步过滤低贡献高斯；然后通过 3D 哈希编码将筛选后高斯的位置映射为空间特征，结合高斯其他属性和相机位姿作为辅助输入丰富特征空间；接着通过解码器生成精化高斯，配合设计的偏移池（offset pool）增强细节表达能力；最后将解码高斯与原始高斯通过混合操作组合——原始高斯捕获粗粒度场景结构，解码高斯补偿缺失的细粒度细节。这种隐式高维特征编码替代了显式致密化，将 GPU 内存需求从高斯数量驱动转化为特征网络参数驱动，在单张 24 GB GPU 上实现大规模场景训练，同时保持与分治方法可比的渲染速度，在大规模场景基准上达到最优渲染质量。

\paragraph{动量教师解码器与重建引导块加权的自蒸馏}

Momentum-GS（2024.12）~\cite{2412_04887} 利用动量自蒸馏促进分块训练间的一致性和准确性。该方法解决混合表示（hybrid representation）在并行分块训练中的两个核心问题：（1）独立训练各块导致数据多样性不足，重建精度下降；（2）并行训练将分块数量限制为可用 GPU 数量。

Momentum-GS 通过动量更新维护一个教师高斯解码器（momentum Gaussian decoder），该解码器以指数移动平均方式聚合所有块的知识，为每个块提供稳定的全局参考。在线高斯解码器（student）在每个训练步骤中同时最小化重建损失和与教师解码器的一致性损失，从而在保持各块独立训练灵活性的同时实现跨块空间一致性。此外，\emph{重建引导块加权}机制动态调整各块的训练权重——重建质量较差的块获得更高权重，使共享解码器优先关注薄弱区域，防止收敛到局部最优。

该方法将分块数量从 GPU 数量解耦：从 $n$ 个块中周期性采样 $k$ 个块分配到 $k$ 个 GPU，实现超 GPU 数量的灵活扩展。相比 CityGaussian，LPIPS 提升 18.7\%，且使用更少的分割块，已被 ICCV 2025 接收。

\subsection{多模态与特殊场景}
\label{ss:14_10}

本节介绍多模态与特殊场景重建技术，将 3DGS 扩展到事件相机、热红外、声纳、射频等非可见光模态和特殊应用场景。

\subsubsection{多模态与特殊场景}

本路线介绍面向事件相机、热红外、声纳等非可见光模态以及水下、文化遗产等特殊场景的3DGS扩展方法。

大规模 3DGS 的应用远不限于可见光 RGB 图像。在许多专业领域中，传感器模态（事件相机、热红外、声纳、射频）和工作场景（水下、城市射频覆盖、文化遗产）各有独特的物理特性和重建需求，需要针对性的 3DGS 扩展。

\paragraph{事件相机大规模场景}

E-3DGS（2025.02）~\cite{2502_10827} 首次将 3D 高斯用于事件相机的新视角合成，重建大型无界场景。

\textbf{动机。}事件相机输出异步脉冲信号而非标准图像帧，无法直接用于基于图像的 3DGS 训练；同时，事件数据的动态范围和噪声特性与传统 RGB 图像截然不同，需要专门的数据转换和损失函数设计。事件相机的高动态范围（$>$120\,dB vs.\ 标准 RGB 相机约 60\,dB）和对运动模糊的鲁棒性使其在弱光、高速运动等挑战场景中具有独特优势。E-3DGS 将曝光期间产生的事件流积分转换为强度帧，作为 3DGS 优化的监督信号，从而将事件相机的异步脉冲数据适配到基于图像的 3DGS 训练框架中。该方法还贡献了首个面向大规模无界场景的事件数据集，包含真实和合成事件数据，涵盖室内、室外和 UAV（Unmanned Aerial Vehicle）场景。E-3DGS 相比 EventNeRF 在 PSNR 上提升 11--25\%，且重建和渲染速度高出数个数量级，已被 3DV 2025 接收。

\paragraph{热红外场景}

Thermal3D-GS（2024.09）~\cite{2409_08042} 引入物理诱导的 3D 高斯方法，通过神经网络建模大气传输效应和热传导，配合温度一致性约束增强重建精度。热红外成像受大气传输效应（路径上的热辐射衰减和散射）和热传导（相邻物体间的热量扩散）双重物理因素影响，导致合成图像中出现浮游体和边缘模糊等问题。Thermal3D-GS 用神经网络显式建模这两种物理过程，将其融入 3DGS 的渲染管线，使高斯基元不仅学习场景外观，还学习热物理特性。温度一致性约束利用热力学中相邻空间位置温度连续的先验，惩罚高斯基元间不合理的温度跳变。该工作还创建了首个大规模热红外新视角合成数据集 TI-NSD，包含 20 个真实热红外视频场景（室内、室外和 UAV），共计 6664 帧数据，覆盖全天候条件，PSNR 相比基线提升 3.03\,dB，显著缓解了浮游体和边缘模糊问题，已被 ECCV 2024 接收。

\paragraph{声纳水下场景}

NAS-GS（2026.01）~\cite{2601_06285} 提出噪声感知声纳高斯框架，引入双向泼溅（Two-Ways Splatting）技术精确建模声纳成像中的强度累积和透射率计算方向，配合基于高斯混合模型的噪声模型捕获旁瓣、散斑和多径噪声。

\textbf{动机。}声纳成像与光学成像在物理原理上存在根本差异——声纳信号经历入射和回波两个传播方向，且水下环境中的旁瓣干扰、散斑噪声和多径反射远比光学噪声严重，直接套用光学 3DGS 的渲染模型会产生系统性偏差。Two-Ways Splatting 的核心创新在于区分声纳信号的入射和回波两个传播方向，分别建模强度累积和透射率衰减，而非光学渲染中的单向累积。GMM 噪声模型将渲染图像与噪声观测的似然作为训练损失，而非简单的 $L_1/L_2$ 损失，从而防止高斯基元过拟合到噪声模式上。该方法显著提升了渲染速度且不牺牲质量，在模拟和真实大规模海上声纳场景中达到最优的新视角合成和三维重建性能，对水下自主导航、海洋考古和环境监测等应用具有重要意义。

\paragraph{射频场景重建}

PropSplat（2026.05）~\cite{2605_08035} 将 3D 高斯用于无地图射频（Radio Frequency, RF）传播建模，每个高斯编码相对于可学习基线路径损耗模型的标量路径损耗偏移。

\textbf{动机。}传统的 RF 传播预测依赖高精度地理数据（楼层平面图、地形数据库），这些数据在许多场景中不可获取或严重过时；3DGS 的空间分布特性恰好可以隐式学习传播环境的空间结构，无需任何先验地图。基线路径损耗模型采用可学习的路径损耗指数，使高斯偏移量专注于捕捉局部传播异常（如反射、绕射、阴影效应）而非全局衰减趋势。

在大规模户外驱动测试中（6 个 Sub-6 GHz 频率、测量点间距 300\,m、跨多个地形区域），RMSE 达到 5.38\,dB，超越 WRF-GS+（5.87\,dB）、GSRF（7.46\,dB）和 NeRF$^2$（14.76\,dB）。在室内蓝牙低功耗定位实验中，PropSplat 实现平均 0.19\,m 定位误差，比 NeRF$^2$（1.84\,m）优一个数量级。RF-PGS（2025.08）~\cite{2508_16849} 引入平面高斯作为几何基元，配合全结构化射频辐射场和定制多视角损失，实现从稀疏路径损耗谱重建高保真无线传播路径。

\paragraph{无人机资源分配}

LAGS（2026.04）~\cite{2604_16910} 研究多无人机低空高斯泼溅中的资源分配问题，提出分组异构图神经网络（GW-HGNN）自动平衡数据保真度和传输成本。计算延迟相比 MOSEK 求解器降低约 100 倍，达到毫秒级推理。

\paragraph{自监督场景表征学习}

Gaussian2Scene（2025.06）~\cite{2506_08777} 利用 3DGS 的效率和显式特性进行场景级自监督预训练。

\textbf{动机。}3D 目标检测等下游任务需要丰富的 3D 场景理解能力，但标注大规模 3D 数据的成本极高；现有基于 NeRF 的自监督方法受限于体渲染的高计算开销和仅在 2D 空间施加重建目标，无法有效学习底层 3D 几何结构。Gaussian2Scene 采用渐进式两阶段训练策略：第一阶段通过双分支掩码自编码器学习 2D（图像 patch）和 3D（点云）场景表示，实现跨模态对齐；第二阶段用第一阶段重建的点云初始化 3DGS 模型，利用高斯基元的几何位置提供直接 3D 空间监督，同时通过渲染 RGB 图像提供 2D 视觉监督，两者联合强化骨干网络的几何理解和跨模态学习能力。3DGS 的显式特性使 3D 场景重建成为可能（不同于 NeRF 需要密集射线采样），同时避免了体渲染的高内存开销。在多个下游 3D 目标检测基准上，Gaussian2Scene 相比现有预训练方法取得持续改善。

\paragraph{数据扩展与增强}

FreeScale（2026.04）~\cite{2604_10512} 提出确定性感知自由视角采样策略，从不完美重建场景中识别语义有意义且受重建误差影响最小的新视角，将有限的现实图像序列转化为可扩展的高质量训练数据源。

\textbf{动机。}3D 视觉领域面临训练数据稀缺的根本瓶颈——高质量的多视角 3D 数据采集成本高昂，而不完美重建场景中蕴含的几何信息若能安全利用，则可作为廉价的数据生成引擎。其核心发现是：不完美的重建场景可作为丰富的几何代理，但朴素采样会放大重建伪影。FreeScale 通过确定性感知机制评估每个候选视角的重建置信度，筛选出同时具有语义意义和最小重建误差的视角，从而安全地将重建场景作为数据生成引擎使用。该方法生成的高质量合成数据具有双重用途：（1）扩大前馈 NVS 模型的训练数据规模，在 OOD（Out-of-Distribution）基准上 PSNR 提升 2.7\,dB；（2）主动增强逐场景 3DGS 优化，通过合成新视角补充稀疏的真实观测，在多个数据集上一致改善重建质量。FreeScale 为解决 3D 视觉领域训练数据稀缺这一根本瓶颈提供了实用且强大的数据生成方案，已被 CVPR 2026 接收。

\begin{takeaway}
自监督场景表征学习和数据扩展策略从不同角度突破大规模重建的瓶颈：3DGS 的显式特性使跨模态自监督预训练成为可能，而确定性感知采样则将不完美重建转化为高质量训练数据生成引擎。
\end{takeaway}

% ======================================================================

\section{近期研究（2026 年 8--11 月）}
\paragraph{HiCo-GS：让八叉树锚点共享层级上下文并约束几何}
2026.08~\cite{2608_14136}。八叉树多层原语能应对城市视距变化，但不同层的锚点各自优化，会割裂粗结构与细节。HiCo-GS 在上层、当前层和下层间查询并融合上下文，再约束由高斯协方差得到的法向与由深度估计的法向接近。边缘感知平滑用来维持大片表面，同时避免把明显图像边界抹平。

论文在 Mill19、UrbanScene3D、MatrixCity 和 China-Pagoda 比较渲染和几何结果。几何约束利用城市表面近似平面的性质，在强曲率、非平面和纹理与几何边界不一致的区域，应单独观察其作用。

\paragraph{ABCD：用预合成的前景和背景训练当前空间分块}
2026.08~\cite{2608_27735}。ABCD 每轮只优化一个空间分块。对相机和活动凸分块，暂时冻结的场景按前后关系预合成为图像，当前分块的渲染再与它们合成。优化因此能保留来自冻结部分的遮挡与外观，而更新显存主要放活动参数和合成缓冲。

论文在 NeRF 的 kitchen 与 garden 比较常规 3DGS、消融合成的分块方法和完整方法。作者指出样本规模小，光栅、优化器和图像缓冲在显存中占比不小。文中的渐近显存推断与这两个场景实际测得的节省应分开看。

\paragraph{Atlas：为头显双目和磁盘流式调度共同设计高斯渲染}
2026.09~\cite{2609_02352}。Atlas 将多数城市高斯放在磁盘，按头显当前所需内容异步调入，利用帧间连贯性减少细节层级搜索，并复用双眼渲染的部分计算。这里优化的是大场景的运行时存储和显示路径，而非从图像学习高斯。

论文对软件双目渲染做图像质量比较，并用 RTL 实现与移动 GPU 和专用渲染硬件比较。双目协议使用固定眼间距和分辨率。专用架构的速度与能耗结果依赖模拟或硬件配置，不能直接视为任意现有头显都可达到的实测值。

\paragraph{InceptionGS：将场景几何作为生成补全的条件}
2026.09~\cite{2609_02747}。大场景的随手采样会造成视角密度高度不均，常规高斯拟合在弱观测区域容易留下细节空缺。InceptionGS 先优化基础 3DGS，同时把场景特定的层次几何用作生成先验的条件，再以光度混合和随视角变化的采样将生成信号逐步反馈到高斯。

论文在 GigaNVS 与 Mip-NeRF 360 等场景比较重建和生成基线。生成细节应与真实图像可验证的内容区分。作者在结论中还列出两个阶段均需逐场景计算，限制其即时处理大批量场景。

\paragraph{OceanXL：按点预算分块并轻量密度控制}
2026.09~\cite{2609_29985}。大范围水下视频重建既要控制单块优化规模，又要保留远处结构。OceanXL 用二元空间划分使各块点数符合 GPU 预算，并在块内采用水下适配的密度控制与选择性高斯差分激活。划分边界随点分布调整，而不是始终等距切分。

原文在 ISRO 分块例中将预算设为每块 4 万点，并和固定中点的 BlockGaussian 切分作对比；这个预算是实验配置。多水下数据对照给出画质与计算成本折中。作者说明当前深度分层固定，也未显式处理随时间变化的场景。

\paragraph{Budgeted-GS：以容量预算约束训练和细节层级渲染}
2026.10~\cite{2610_03162}。Budgeted-GS 一方面将已有的高斯模型整理成可按预算调度的细节层级树，另一方面在训练时按估计的容量要求设定原语上限，在增密过程中控制模型大小。两条路径分别服务已有模型的部署和已知预算下的直接训练。

作者报告训练侧匹配预算的比较以及大场景渲染实验。论文限于单 GPU 设置，透明场景、极端遮挡和跨存储流式渲染尚未验证。容量分析只应按文中定义的渲染器和一致性泛函理解。

\paragraph{RoMo-GS：在线建图中对曝光运动做解析积分}
2026.10~\cite{2610_24605}。RoMo-GS 把曝光期间投影高斯沿近似线性轨迹的贡献写成解析积分，用一次光栅化逼近以往多时刻采样平均的结果。在线后端联合修正位姿、速度和高斯，另以局部优化与基于 Hessian 的评分抑制原语增长。

作者在合成模糊、公开视觉数据和机器人现场录制中比较图像质量和速度。当前积分依赖曝光时长和曝光期线性运动假设，轨迹弯曲、相机参数不明或前端跟踪失败均会影响映射质量。

\paragraph{ReScene：把稀疏复访中的高斯变化写成可回放记录}
2026.11~\cite{2611_20433}。复访通常只提供少量新图像，又未必能取回早期帧。ReScene 根据图像差异在高斯合成时的贡献，把变化归因给部分原语，用增补和可回放的更新记录保留过去版本。已经保存的版本可以按记录恢复，这里的精确指对历史模型的复现。

论文分开比较当前重建、历史回放、几何与存储，并以 Multi-State、WAT 和 ChronoScene 为测试。方法要求提供统一坐标中的标定位姿，复访间缺少连续观测。它能保留各次重建，但不能反推中间的真实运动，也会保存早期模型已有的错误。

\paragraph{DOOM-GS：以时空分块和流式泼溅训练城市级动态高斯}
2026.11~\cite{2611_22805}。动态场景离开显存训练时，一个视角的可见原语也可能塞不进显存。DOOM-GS 按空间与时间组织数据块，利用相邻相机可见集合重叠来排程，并以流式泼溅和梯度估计在有界显存下更新模型。

论文在加入程序生成车辆和行人的 DynamicMatrixCity 比较单一动态模型与逐帧静态基线，并使用留出相机质量、静态区闪烁和训练时间衡量。结果针对该模拟城市场景和相机划分，真实交通动态的效果需另行验证。


\section*{本章小结}
\addcontentsline{toc}{section}{本章小结}
\label{sec:14_summary}

\begin{futurebox}{本章总结与展望}{ch14_summary}
本章系统梳理了大规模与在线 3DGS 重建领域的 106 篇论文，覆盖五个核心方向。

城市场景重建（43 篇）是论文最密集的方向，分治策略已成为处理城市场景的标准范式。从 CityGaussian 的基础分治+LOD 方案，到 BlockGaussian 的内容感知分割、OccluGaussian 的遮挡感知分割、TraGraph-GS 的轨迹图引导分割，分割策略不断精细化。几何约束方面，LI-GS 展示了 LiDAR 融合对几何精度的巨大提升，UrbanGS 和 MetroGS 通过深度正则化和混合优化进一步增强几何质量。跨视角融合方面，Horizon-GS 和 DRAGON 分别从统一训练和迭代桥接角度解决航拍-地面视角差异。语义辅助方面，GS4City 和 GS4Buildings 展示了结构化城市模型先验的巨大价值。

在线与增量重建（23 篇）从流式重建和场景更新两个角度解决时效性问题。流式重建方法从 Gaussian On-the-Fly Splatting 的渐进式优化到 StreamGS 的无位姿在线重建、LongSplat 的长序列压缩、PLANING 的松耦合混合表示，逐步逼近实时在线重建的目标。场景更新方面，CL-Splats、GaussianUpdate、GaME 等方法通过变化检测和局部优化避免全局重计算。

分层与 LOD 表示（13 篇）是大规模渲染效率的关键。从 Hierarchical 3D Gaussians 的开创性层次结构，到 Octree-GS 的八叉树组织、PyGS 的金字塔表示，层次结构设计日趋成熟。渲染优化方面，FilterGS 的无遍历并行过滤、Nebula 的 VR 协作渲染、Voyager 的移动端加速代表了 LOD 渲染的前沿。虚拟内存技术（A LoD of Gaussians、Virtual Memory for 3DGS）使消费级硬件能够处理城市级场景。

分布式与并行训练（9 篇）是突破单 GPU 内存限制的核心手段。Grendel 和 RetinaGS 分别从分布式系统和模型并行角度探索了高斯数量扩展极限。DOGS 引入 ADMM 共识约束保证分布式收敛，BlitzGS 通过三层工作负载管理实现城市级场景的快速训练。F3DGS 将联邦学习引入 3DGS，在保护数据隐私的前提下实现多智能体协同重建。

其他大规模方法（18 篇）覆盖了核外训练（TideGS、GS-Scale、CLM）、多模态重建（事件相机、热红外、声纳、射频）、自监督学习（Gaussian2Scene）和数据增强（FreeScale）等方向，拓展了大规模 3DGS 的应用边界。

展望未来，大规模与在线 3DGS 重建仍面临多个开放问题：（1）\textbf{十亿级高斯的实时渲染}——TideGS 已实现在单 GPU 上训练十亿高斯，但实时渲染仍需 LOD 和核外技术的进一步突破；（2）\textbf{持续终身更新}——如何在场景持续数月甚至数年变化中维护高质量的 3DGS 模型；（3）\textbf{异构设备协同}——从云端 GPU 到边缘设备的无缝协作渲染；（4）\textbf{语义-几何联合优化}——将高层语义理解融入大规模重建以提升质量和效率。随着这些问题的逐步解决，3DGS 有望成为城市数字孪生、自动驾驶高精地图、虚拟现实/增强现实等应用的核心基础设施。
\end{futurebox}
